% 本文件由 LaTeX 排版 Agent 根据有效 translation.json 自动生成。
% author=Theodoret; work_id=2703; title=对话录

\chapter{对话录}

% structure: p0001 source=h1 target=omitted reason=page-title-already-rendered-by-parent

\Emph{\Quote{Eranistes}或\Quote{Polymorphus}}\par

序言\newline{}对话一\newline{}对话二\newline{}对话三\par

\SourceRecord{Translated by Blomfield Jackson.；Nicene and Post-Nicene Fathers, Second Series；Vol. 3.；Edited by Philip Schaff and Henry Wace.；Buffalo, NY: Christian Literature Publishing Co.,；1892.；<http://www.newadvent.org/fathers/2703.htm>.}

% structure: p0001 source=h1 target=section reason=page-title

\section{对话录（序言）}

真福\Place{居鲁士}主教\Person{狄奥多勒}的\WorkTitle{厄拉尼斯忒斯}或\WorkTitle{多形者}。\par

有些人，既无家世，又无教育，也没有因正直生活带来的可敬名声，却以邪恶的方式追求出名。\newline{}著名的铜匠\Person{亚历山大}便是其中之一：他毫无出众之处——没有高贵的出身，没有雄辩的口才，从未领导过政治派别或军队，也从未在战斗中展现英勇，只是日复一日地操持他那卑微的行业，而他的出名仅仅是因为对圣\Person{保禄}的疯狂暴力。\par

再者，一个身份卑微、奴仆等级的\Person{史米}，因他对圣\Person{达味}的鲁莽攻击而变得非常出名。\par

据说，摩尼教异端的创始人不过是一个常受鞭打的奴隶，出于对名声的渴望，编写了他那可憎而迷信的著作。\par

现在许多人也追随同样的行径，他们背弃了因需要经历劳作才能获得的荣誉之德，却为自己买来了由羞耻和耻辱带来的名声。因为他们急于成为新学说的捍卫者，收集并拼凑了许多异端的不敬之辞，编撰出这个死亡的异端。\par

现在，我将试图与他们简短地辩论，目的有二：如果可能，治愈他们的不健全，并向全体发出一个警告。\par

我称我的著作为《\WorkTitle{厄拉尼斯忒斯}，或\WorkTitle{多形者}》，因为他们从许多不幸的源头聚集了他们有害的教义，然后产生出拼凑而不协调的虚妄。因为称我们的主基督仅为天主，是\Person{西蒙}、\Person{塞尔多}、\Person{玛尔西安}以及分享这可憎意见的其他人的方式。\par

承认祂由童贞女所生，但同时断言这出生只是一个过渡的过程，且天主圣言并未取走童贞女本性的任何东西，这是从\Person{瓦伦提努斯}、\Person{巴尔得撒纳}及其神话追随者那里偷来的。\par

称主基督的天主性和人性为一个本性，这是从\Person{阿波利纳里}的愚昧中窃取的错误。\par

此外，将受苦的能力归于基督的天主性，是对\Person{亚略}和\Person{欧诺弥}亵渎的剽窃。因此，他们教导的主要原则就像乞丐的破袍——一件由不协调的碎布拼成的杂衣。\par

因此，我称这部著作为\WorkTitle{厄拉尼斯忒斯}或\WorkTitle{多形者}。我将以对话的形式写作，包含问答、命题、解答、对辩以及对话应有的一切。我不会像古代有智慧的希腊人那样，将提问者和回答者的名字插在对话正文中，而是将它们写在段落开头的页边。\newline{}他们确实将他们的作品交到受过高度、多方面教育且视文学为生命的读者手中。相反，我希望我所写的内容的阅读以及其中任何好处的发现，即使对未受教育者也容易。我认为，如果对话者的名字清楚地写在页边，那么这将更容易实现：因此，为教令辩护的论辩者被称为\Quote{正统者}，而他的对手被称为\Quote{厄拉尼斯忒斯}。\newline{}一个靠许多人的施舍生活的人，我们通常称为\Quote{乞丐}；一个知道如何聚敛钱财的人，我们称为\Quote{敛财者}。因此，我们根据他的性格和行为给我们的论辩者起了这个名字。\par

我恳求所有拿到我这本书的人，放下一切先入之见，检验真理。为清晰起见，我将我的书分为三篇对话。\newline{}第一篇将包含对独生子天主性不变的论证。第二篇将藉天主的助佑，说明主基督的天主性和人性的结合是无混淆的。第三篇将论证我们的救主的天主性不受苦难。\newline{}在这三场辩论之后，我们将附加几个对话作为补充，在每个标题下提供正式证明，并完全清楚地表明宗徒的教义由我们保存。\par

\SourceRecord{Translated by Blomfield Jackson.；Nicene and Post-Nicene Fathers, Second Series；Vol. 3.；Edited by Philip Schaff and Henry Wace.；Buffalo, NY: Christian Literature Publishing Co.,；1892.；<http://www.newadvent.org/fathers/27030.htm>.}

% structure: p0001 source=h1 target=section reason=page-title

\section{对话一}

\Emph{正统者与厄兰尼斯特。}\par

\Emph{正统者：}— 我们最好还是同意并持守纯洁的宗徒教导。但是，不知为何，你破坏了和谐，现在又向我们提出新的教义，如果你愿意，让我们不要争吵，来探究真理。\par

\Emph{厄兰尼斯特：}— 我们不需要探究，因为我们确切地持守真理。\par

\Emph{正统者：}— 每个异端者都这样认为。是的，连犹太人和外教人也认为自己是在维护真理的教义；不仅如此，柏拉图和毕达哥拉斯的追随者，甚至伊壁鸠鲁派，以及那些完全不信神或没有信仰的人，也是如此。然而，我们不应该成为先入为主之见的奴隶，而应该寻求认识真理。\par

\Emph{厄兰尼斯特：}— 我承认你所说的话有道理，并准备按你的建议行事。\par

\Emph{正统者：}— 既然你毫不困难地听从了我这个初步的劝告，我接下来要求你不要让人来判定对真理的探究，而是追踪宗徒、先知以及跟随他们的圣人的足迹。因为行路者在迷路时，往往会仔细察看小路，如果那里有任何人、马、驴或骡子往这个或那个方向留下的印记，一旦发现，他们就像狗一样追踪这些足迹，直到重新走上正路。\par

\Emph{厄兰尼斯特：}— 我们就这样做吧。你既然开始了讨论，就由你来引导。\par

\Emph{正统者}— 那么，让我们首先仔细深入地探究这些神性名称——我是指本体、本质、位格和特性，让我们学习并确定它们彼此之间的区别。然后我们再处理后续的问题。\par

\Emph{厄兰尼斯特：}— 你为我们的论证给出了一个非常出色且恰当的开端。这些点清楚了，我们的讨论就会毫无阻碍地顺利展开。\par

\Emph{正统者：}— 既然我们已经确定这必须是我们的行事程序，那么请告诉我，朋友，我们是否承认天主的一个本体，既属于圣父，也属于独生子及圣神，正如我们从新旧约圣经以及从在\Place{尼西亚}召开的大公会议的神长们所学的？还是追随亚略的亵渎之言？\par

\Emph{厄兰尼斯特：}— 我们宣认天主圣三的一个本体。\par

\Emph{正统者：}— 我们是否认为位格表示本体以外的意思，还是把它当作本体的另一个名称？\par

\Emph{厄兰尼斯特：}— 本体和位格之间有什么区别吗？\par

\Emph{正统者}— 在基督宗教之外的哲学中，没有区别，因为\OriginalTerm{本体}{\GreekText{οὐσία}}表示\OriginalTerm{存在}{\GreekText{τὸ} \GreekText{ὄν}}，即那存在的，而\OriginalTerm{位格}{\GreekText{ὑπόστασις}}表示那自立存在的。但是根据教父们的教导，\OriginalTerm{本体}{\GreekText{οὐσία}}和\OriginalTerm{位格}{\GreekText{ὑπόστασις}}之间的区别就如同共通与特殊、种类与个体之间的区别。\par

\Emph{厄兰尼斯特：}— 请你更清楚地告诉我，什么是种类、类别和个体。\par

\Emph{正统者：}— 我们就动物来说种类，因为它同时表示许多事物。它既指理性的，也指非理性的；非理性的动物又有许多类：飞行的、两栖的、行走的和游泳的。这些类别中的每一种又有很多分支：行走的动物中有狮子、豹子、公牛和无数的其他种类。同样，飞行动物和其他种类也有许多类；然而，尽管类别如前所述，它们都属于同一个动物种类。同样，\Quote{人}这个名字是人类共通的名称；因为它表示罗马人、雅典人、波斯人、萨尔马提亚人、埃及人，总而言之所有人类；而\Quote{保禄}或\Quote{伯多禄}这个名字则不表示种类的共通性，而是特定的某个人；因为没有人听到保禄会想到亚当、亚巴郎或雅各伯，而只会想到他所听到名字的那个人。但如果他仅仅听到\Quote{人}这个词，他不会将意念集中于某个人，而是想到印度人、斯基泰人、马萨格特人，以及所有人类种族。我们不仅从自然界学到这一点，也从圣经中学到，因为我们读到天主说：\Quote{我要从地面上消灭人}，这话是对无数群众说的；在亚当之后过了两千两百多年，他通过洪水给人类带来了普遍的毁灭。因此，蒙福的达味说：\Quote{人在尊贵中而不醒悟}，他不是指控这一个或那一个，而是指责所有人类共同。还可以找到上千个类似的例子，但我们不能冗长。\par

\Emph{厄兰尼斯特：}— 共通与特殊之间的区别已经清楚地显示出来了。现在让我们回到关于\OriginalTerm{本体}{\GreekText{οὐσία}}和\OriginalTerm{位格}{\GreekText{ὑπόστασις}}的讨论。\par

\Emph{正统者：}— 那么，正如\Quote{人}这个名称是人性共通的，我们理解天主本体指示天主圣三；但\Quote{位格}表示每一位，如圣父、圣子和圣神；因为遵循圣教父们的定义，我们说位格和个体是同一回事。\par

\Emph{厄兰尼斯特：}— 我们同意是这样。\par

\Emph{正统者：}— 那么，凡是对天主本性所言的，对圣父、圣子、圣神都是共通的，例如\Quote{天主}、\Quote{主}、\Quote{创造者}、\Quote{全能者}等等。\par

\Emph{厄兰尼斯特：}— 毫无疑问，这些词对圣三都是共通的。\par

\Emph{正统者：}— 但所有自然指示位格的东西就不再是圣三共通的，而是指示它所归属的位格，例如\Quote{父}、\Quote{无生}这些名称是圣父特有的；而\Quote{子}、\Quote{独生子}、\Quote{天主圣言}这些名称既不指示圣父，也不指示圣神，而是指示圣子；而\Quote{圣神}、\Quote{护慰者}这些词自然指示圣神的位格。\par

\Emph{厄兰尼斯特：}— 但是圣经不也称圣父和圣子为\Quote{圣神}吗？\par

\Emph{Orth.}——是的，圣经称圣父和圣子都为\Quote{\OriginalTerm{神}{Spirit}}，以此术语表示神性本体的非形体性和无限性。圣经只称圣神的位格为\Quote{\OriginalTerm{圣神}{Holy Ghost}}。\par

\Emph{Eran.}——这是无可争议的。\par

\Emph{Orth.}——既然我们断言有些术语是至圣圣三共有的，有些是每个位格特有的，那么，我们说\Quote{\OriginalTerm{不变}{immutable}}这个术语是属于本体共有的，还是属于某个位格特有的呢？\par

\Emph{Eran.}——\Quote{不变}这个术语是圣三共有的，因为不可能本体的一部分可变，另一部分不变。\par

\Emph{Orth.}——你说得好，因为正如\Quote{可朽}这个术语是人类共有的，同样，\Quote{不变}和\Quote{\OriginalTerm{不变}{invariable}}是至圣圣三共有的。因此，独生子是不变的，生祂的圣父和圣神也是如此。\par

\Emph{Eran.}——不变。\par

\Emph{Orth.}——那么，你为何提出福音中的话：\Quote{圣言成了血肉}，并将变化归于不变的本性呢？\par

\Emph{Eran.}——我们断言祂成为血肉不是通过变化，而是照祂自己所知道的方式。\par

\Emph{Orth.}——如果祂不是因取了血肉而被称为成为血肉，那么就必须断言两件事之一：要么祂经历了变成血肉的变化，要么祂只是表象上看起来如此，而实际上祂是无血肉的天主。\par

\Emph{Eran.}——这是瓦伦提努、马西昂和摩尼教门徒的教义，但我们确实被教导说，圣言成了血肉，这是没有争议的。\par

\Emph{Orth.}——但你说\Quote{成了血肉}是什么意思？是\Quote{取了血肉}，还是\Quote{变成了血肉}？\par

\Emph{Eran.}——正如我们听见福音传道者所说：\Quote{圣言成了血肉}。\par

\Emph{Orth.}——你如何理解\Quote{成了}？\par

\Emph{Eran.}——那经历了变成血肉的变化的，就成了血肉，并且，正如我刚才所说，照祂自己所知的方式。但我们知道，在祂一切都是可能的，因为祂将尼罗河的水变为血，将白昼变为黑夜，使海成为干地，使干涸的旷野充满水，我们也听见先知说：\Quote{\BibleQuote{凡主所喜悦的，祂在天上、地上、海洋和一切深渊中都行了}}。\par

\Emph{Orth.}——受造物被造物主随意改变，因为它是可变的，服从于造物主的命令。但祂的本性是不变和不变的，因此先知论及受造物说：\Quote{\BibleQuote{祂创造并改变万有}}。但论及圣言，伟大的达味说：\Quote{\BibleQuote{祢原是一样，祢的岁月没有穷尽}}。而且，同一上帝论及自己说：\Quote{\BibleQuote{因为我上主是不变的}}。\par

\Emph{Eran.}——隐藏的事不应该\Quote{被探究}。\par

\Emph{Orth.}——也不应该完全忽略明了的事。\par

\Emph{Eran.}——我不知道道成肉身的方式。我听见说圣言成了血肉。\par

\Emph{Orth.}——如果祂是通过成了血肉而发生变化，那么祂就不再是祂以前所是的，这可以通过几个类比容易理解。例如，沙子受热时，先变成流体，然后变化并凝结成玻璃，在变化时它的名字也改变了，不再被称为沙，而是玻璃。\par

\Emph{Eran.}——是这样。\par

\Emph{Orth.}——我们称葡萄树的果实为葡萄，一旦我们压榨了它，就不再称它为葡萄，而称为酒。\par

\Emph{Eran.}——当然。\par

\Emph{Orth.}——酒本身在发生变化后，我们习惯上不再称它为酒，而称它为醋。\par

\Emph{Eran.}——正确。\par

\Emph{Orth.}——同样，石头经过燃烧和溶解后，不再被称为石头，而称为石灰。还可以找到无数其他类似的例子，其中变化伴随着名称的改变。\par

\Emph{Eran.}——同意。\par

\Emph{Orth.}——因此，如果你断言圣言经历了变成血肉的变化，你为什么称祂为天主，而不称祂为血肉？因为名称的改变与本质的变化相一致。因为即使那些经历变化的东西与它们以前的状态有某种关系（例如醋与酒，酒与葡萄的果实，玻璃与沙有一定近似），它们在变化后也会得到另一个名字。那么，当它们之间的差别是无限的，像蚊虫与整个可见和不可见受造物之间的差别那样大（实际上血肉与神性本质之间的差别更大，甚至更大），怎么可能在变化后仍然使用同一个名称呢？\par

\Emph{Eran.}——我已经不止一次说过，祂成为血肉不是通过变化，而是继续是祂以前所是的，祂成为了祂以前所不是的。\par

\Emph{Orth.}——但除非\Quote{成了}这个词变得非常清楚，否则它会暗示变化和改变，因为除非祂是通过取了血肉而成了血肉，否则祂就是通过经历变化而成了血肉。\par

\Emph{Eran.}——但\Quote{取}这个词是你自己的发明。福音传道者说圣言成了血肉。\par

\Emph{Orth.}——你似乎要么对圣经无知，要么故意错解。如果你是無知，我会教导你；如果你是故意错解，我会责备你。那么请回答：你承认神圣保禄的教导是来自圣神吗？\par

\Emph{Eran.}——当然。\par

\Emph{Orth.}——你同意同一位圣神通过福音传道者和宗徒工作吗？\par

\Emph{Eran.}——是的，因为我从宗徒的圣经中学会了：\Quote{\BibleQuote{恩赐虽有区别，却是同一圣神}}，又说：\Quote{\BibleQuote{这一切都是这同一圣神所行的，祂随意分给各人}}，又说：\Quote{\BibleQuote{我们既然有同一信德的圣神}}。\par

\Emph{Orth.}——你引用宗徒的见证正是时候。如果我们断言福音传道者和宗徒的教导都出于同一位圣神，那么请听宗徒如何解释福音的话，因为在《希伯来人书》中他说：\Quote{\BibleQuote{祂确实没有援助天使，而是援助了亚巴郎的后裔}}。现在告诉我，你所说的\Quote{亚巴郎的后裔}是什么意思？那自然属于亚巴郎的，不也同样属于亚巴郎的后裔吗？\par

\Emph{Eran.}——不，并非无一例外，因为基督没有犯罪。\par

\Person{Orth.}——罪并非出于本性，而是出于败坏的意志。正因如此，我没有泛泛地说亚巴郎有什么，而是说他按照本性有什么，即身体和理性的灵魂。现在请明白告诉我：你肯承认亚巴郎的后裔具有身体和理性的灵魂吗？若不肯，你在这点上与阿波里纳里的狂言一致。但我将用其他方式迫使你承认这一点。现在告诉我：犹太人是否有身体和理性的灵魂？\par

\Person{Eran.}——当然有。\par

\Person{Orth.}——那么，当我们听见先知说：\Quote{但你、以色列，我的仆人，雅各伯，我所拣选的，我朋友亚巴郎的后裔}，我们是把犹太人理解为仅有身体吗？难道我们不把他们理解为由身体和灵魂构成的人吗？\par

\Person{Eran.}——对。\par

\Person{Orth.}——而亚巴郎的后裔并非没有灵魂或智力，乃是一切表征亚巴郎后裔的特征都具备？\par

\Person{Eran.}——这样说的人提出了两个儿子。\par

\Person{Orth.}——但那位说圣言变成了血肉的人，连一个儿子也不承认，因为单是血肉本身并非儿子；但我们宣认一位圣子，祂按照神圣宗徒的话，取了亚巴郎的后裔，并为人类成就了救恩。但若你不接受宗徒的宣讲，就请公开说明。\par

\Person{Eran.}——但我们坚持宗徒的言论不一致，因为\Quote{圣言成了血肉}与\Quote{取了亚巴郎的后裔}之间似乎存在某种不一致。\par

\Person{Orth.}——这是因为你缺乏智识，抑或是你为争论而争论，才使一致显得不一致。对于使用正确理性的人而言，并不如此；因为神圣宗徒教导说，圣言成了血肉，并非通过改变，而是通过取了亚巴郎的后裔。同时，他也想起了给亚巴郎的恩许。难道你不记得宇宙的天主给圣祖的恩许吗？\par

\Person{Eran.}——什么恩许？\par

\Person{Orth.}——当日主将他从他父家领出来，命他往巴勒斯坦地去时，岂没有对他说：\Quote{我要祝福祝福你的人，诅咒诅咒你的人；地上万民都要因你的后裔蒙受祝福}？\par

\Person{Eran.}——我记得这些恩许。\par

\Person{Orth.}——也要记得天主与依撒格和雅各伯所立的盟约，因为祂也给了他们同样的恩许，以第二次和第三次确认了前一次。\par

\Person{Eran.}——这些我也记得。\par

\Person{Orth.}——正是针对这些盟约，神圣宗徒在写给迦拉达人的书信中说：\Quote{恩许原是向亚巴郎和他的后裔说的。} 他没有说\Quote{后裔们}，如同指许多人，而是如同指一位……就是基督，非常清楚地表明基督的人性出自亚巴郎的后裔，并实现了给亚巴郎的恩许。\par

\Person{Eran.}——宗徒是这样说的。\par

\Person{Orth.}——所说的话已足以消除对此点所引发的一切争议。但我仍愿提醒你另一个预言。赐予圣祖雅各伯及其父祖的祝福，由他单独给了他的儿子犹大。他说：\Quote{权杖不离犹大，统帅不离他的腰间，直到那应得者来到，他是外邦人所期待的。} 难道你不接受这预言是指救主基督说的吗？\par

\Person{Eran.}——犹太人对这类预言作了错误的解释，但我是基督徒；我信靠圣言；我毫无怀疑地接受预言。\par

\Person{Orth.}——既然你承认你相信预言，并承认这些预言是神圣地关于我们救主而说的，那么请思考宗徒话语的意图；当他指出给圣祖们的恩许已得成就时，他说了那些著名的话：\Quote{他没有承袭天使的本性}，几乎是在说：恩许是真实的；主实现了祂的承诺；祝福的源泉向外邦人敞开；天主取了亚巴郎的后裔；藉着这后裔，祂带来所应许的救恩；藉着这后裔，祂证实了对外邦人的恩许。\par

\Person{Eran.}——先知的话与宗徒的话契合得极妙。\par

\Person{Orth.}——同样，神圣宗徒提醒我们犹大的祝福，并指出它如何得到应验，高声说：\Quote{显然我们的主由犹大出生。} 米该亚先知和圣史玛窦也是这样。前者说出了他的预言，后者将预言与其叙述相连。特别之处在于，他说真理的公开敌人们明确告诉黑落德，基督在白冷出生，因为经上记着：\Quote{犹大地白冷啊，你在犹大的郡邑中绝非最小的，因为将由你出来一位领袖，他将牧养我的百姓以色列。} 现在让我们补充犹太人出于恶意省略的部分，从而使见证不完整。因为先知在说了\Quote{将由你出来一位统治以色列的领袖}之后，又加上\Quote{他的由来是从亘古，从永远就有}。\par

\Person{Eran.}——你引述先知全部的证据做得很好，因为他指出在白冷出生的那位是天主。\par

\Person{Orth.}——不仅是天主，也是人；按肉身出自犹大并在白冷出生的人；而作为天主，则万世之前就已存在。因为\Quote{将由你出来一位统治我的领袖}这话显示了他末后日子按肉身的诞生；而\Quote{他的由来是从亘古，从永远就有}则清楚地宣告了他万世之前的存在。同样，神圣宗徒在写给罗马人的书信中，哀叹犹太人昔日福乐的恶化，并追想他们的神圣恩许和立法，接着说：\Quote{圣祖们是他们的，按肉身说，基督也是从他们来的，基督是在万有之上，永远受赞美的天主，阿们。} 在同一段话中，他既显明祂是万物的创造者、主宰和统治者，作为天主；又显明祂按肉身从犹太人中出生，作为人。\par

\Emph{Eran.}——好了；你已经解释了这些段落，对于耶肋米亚的预言你要说什么？因为这预言宣告他是唯一的天主。\par

\Emph{Orth.}——你说的是哪个预言？\par

\Emph{Eran.}——\BibleQuote{这就是我们的天主，没有别的可与他相比——他找出了全部知识的道路，并将它赐给了他的仆人雅各伯和他所爱的以色列。此后他在地上显现，与世人交往。}\par

在这些话里，先知既没有讲肉体，也没有讲人性，也没有讲人，而只讲天主。\par

\Emph{Orth.}——那么辩论有什么用？我们说天主本性是不可见的吗？还是我们不同意宗徒所说的\BibleQuote{不可朽坏、不可见、唯一的天主}？\par

\Emph{Eran.}——毫无疑问，天主本性是不可见的。\par

\Emph{Orth.}——那么，没有身体，不可见的本性如何能被看见？还是你不记得宗徒清楚教导天主本性不可见的那句话？他说\BibleQuote{没有人见过，也不能看见}。因此，如果天主本性对人不可见，我还要加上对天使也不可见，请告诉我，那位不能被看见和注视的，如何在地上被看见了？\par

\Emph{Eran.}——先知说他在地上被看见了。\par

\Emph{Orth.}——而宗徒说\BibleQuote{不可朽坏、不可见、唯一的天主}以及\BibleQuote{没有人见过，也不能看见}。\par

\Emph{Eran.}——那么怎样？先知在说谎吗？\par

\Emph{Orth.}——绝对不可。两种说法都是圣神的话。\par

\Emph{Eran.}——那么让我们查考如何不可见的被看见了。\par

\Emph{Orth.}——我求你，不要引入人的理性。我只服从圣经。\par

\Emph{Eran.}——你不会得到任何未经圣经证实的论证，如果你给我任何从圣经推导出的问题解答，我会接受，绝不反对。\par

\Emph{Orth.}——你知道前一刻我们如何藉宗徒的见证阐明了福音的话；神圣的宗徒清楚地告诉我们圣言如何成了血肉，明确地说\BibleQuote{因为他没有取得天使的本性，而是取得了亚巴郎的后裔}。这位老师也将教导我们神圣的圣言如何在地上被看见，并在人间居住。\par

\Emph{Eran.}——我服从宗徒和先知的话。那么请照你的应许，给我展示预言的解释。\par

\Emph{Orth.}——神圣的宗徒写信给弟茂德也说\BibleQuote{毫无疑问，伟大的虔敬奥秘：天主在肉身显现，在圣神内成义，被天使看见，被传于外邦人，被世人信服，被接在荣耀中。}\par

由此可见，天主本性是不可见的，而肉身是可见的，藉着可见的，不可见的被看见了，藉着肉身行奇事，并显明自己的能力，因为祂用手造了视觉，并治好了生来瞎眼的人。祂又赐给聋子听觉，解开了哑巴的舌头，用指头作工具，用唾沫作某种医治的药。同样，当祂在海面上行走时，显示了天主性的全能。因此，宗徒贴切地说\BibleQuote{天主在肉身显现}。因为藉着它，不可见的本性藉着肉身被天使军旅看见，因为他说\BibleQuote{他被天使看见}。\par

无形体者的本性也与我们同享了这一恩惠。\par

\Emph{Eran.}——那么天使在救主显现之前没有看见天主吗？\par

\Emph{Orth.}——宗徒说他\BibleQuote{在肉身显现，并被天使看见}。\par

\Emph{Eran.}——但主说\BibleQuote{你们要小心，不可轻视这些小子中的一个，因为我告诉你们，他们的天使在天上常见我父的面}。\par

\Emph{Orth.}——但主又说\BibleQuote{不是有人见过父，惟独从天主来的，他见过父}。因此福音作者明确宣告\BibleQuote{从来没有人见过天主}，并证实了主的话，因为他说\BibleQuote{独生子，在父怀里，他表明了父}。伟大的梅瑟渴望看见不可见的本性时，听见主天主说\BibleQuote{没有人见了我还能活着}。\par

\Emph{Eran.}——那么\BibleQuote{他们的天使在天上常见我父的面}这话如何理解？\par

\Emph{Orth.}——正如我们通常理解那些据说见过天主的人所说的话一样。\par

\Emph{Eran.}——请说得更清楚些，因为我不明白。天主也能被人看见吗？\par

\Emph{Orth.}——当然不能。\par

\Emph{Eran.}——然而我们听见圣经说天主在玛默勒橡树那里向亚巴郎显现；依撒意亚说\BibleQuote{我看见主坐在高高的宝座上}；米该亚、达尼尔和厄则克耳也说了同样的事。关于立法者梅瑟，记载说\BibleQuote{上主同梅瑟面对面说话，如同人同朋友说话}；天主亲自说\BibleQuote{我要与他口对口说话，明明地，不用谜语}。那么我们要说什么？他们看见了天主本性吗？\par

\Emph{Orth.}——绝不，因为天主亲自说\BibleQuote{没有人见了我还能活着}。\par

\Emph{Eran.}——那么那些说见过天主的人是说谎吗？\par

\Emph{Orth.}——绝对不可——他们看见了他们所能看见的。\par

\Emph{Eran.}——那么仁慈的主按看见祂者的能力适应祂的启示？\par

\Emph{Orth.}——是的；这祂已藉先知表明，因为他说\BibleQuote{我，我增加了异象，藉先知的手作了比喻}。\par

祂没有说\BibleQuote{被看见}，而是说\BibleQuote{作了比喻}。而比喻并不显示所见之物的本性。即使是皇帝的肖像也不显示皇帝的本性，尽管它清晰地保存了他的容貌。\par

\Emph{Eran.}——这模糊不清，不够明白。那么，看见那些显现的人没有看见天主的本质吗？\par

\Emph{Orth.}——没有；谁是疯狂的胆敢这样说？\par

\Emph{埃兰}—— 但毕竟经上说他们看见了。\par

\Emph{奥尔托}—— 的确经上有此说；但我们在虔敬理性的运用下，并依靠那明确宣告\Quote{从来没有人见过天主}的神谕，断言他们所见的并非天主性本身，而是与他们的能力相称的某种显现。\par

\Emph{埃兰}—— 我们也这样说。\par

\Emph{奥尔托}—— 那么，当我们听到天使们每日看见你们天父的面时，也当如此理解。他们所见的并非那不可界定、不可领会、不可把握、包罗万有的神圣本体，而是某种与其本性相称的光荣。\par

\Emph{埃兰}—— 这是公认的。\par

\Emph{奥尔托}—— 然而在降生成人之后，祂也被天使看见了，正如神圣宗徒所说，但不是通过光荣的相似，而是以肉身的真实活体作为某种帷幔。\Quote{天主}，他说，\Quote{在肉身显现，在圣神内成义，被天使看见。}\par

\Emph{埃兰}—— 我接受这是圣经，但我不能接受措辞的新奇。\par

\Emph{奥尔托}—— 我们引入了什么措辞的新奇？\par

\Emph{埃兰}—— 就是那个\Quote{帷幔}。圣经何处称主的肉体为帷幔？\par

\Emph{奥尔托}—— 你似乎不太勤读圣经；否则你便不会非难我们以比喻方式所说的话。首先，神圣宗徒说那不可见的本性藉着肉身得以彰显，这事实就允许我们将肉身理解为神性的帷幔。其次，神圣宗徒在致希伯来人书中明确使用了这个措辞，因为他说：\Quote{弟兄们！我们既因耶稣的血，得以坦然进入至圣所，藉着祂给我们开辟的一条又新又活的路，由帷幔，即祂的肉身，通过；并且我们既然有一位治理天主房屋的大司祭，我们就怀着真诚的心，以完备的信德来亲近天主。}\par

\Emph{埃兰}—— 你的论证无可反驳，因为它建立在宗徒权威之上。\par

\Emph{奥尔托}—— 那就不要指控我们标新立异。我还要引证另一个先知权威，它明确称主的肉身为长袍和外氅。\par

\Emph{埃兰}—— 只要不显得隐晦含混，我们便不反对，并为此感激。\par

\Emph{奥尔托}—— 我要让你亲自为这承诺的真实性作证。你知道族长雅各伯在训谕犹大时，如何将犹大的主权限定在主的诞生之前。\Quote{权杖不离犹大，柄杖不离他脚间，直到那应得的人来到，万民都要归顺他。}你已经承认这预言是关于救主的。\par

\Emph{埃兰}—— 我承认了。\par

\Emph{奥尔托}—— 那么请记住下文所说的话：\Quote{他要在酒中洗自己的衣服，在葡萄汁中洗自己的外氅。}\par

\Emph{埃兰}—— 族长说的是衣服，不是身体。\par

\Emph{奥尔托}—— 那么请告诉我，他何时何地在葡萄汁中洗了自己的外氅？\par

\Emph{埃兰}—— 请反过来说，他何时将身体染红于其中？\par

\Emph{奥尔托}—— 我恳求你恭敬作答。或许有些未入教的人在场。\par

\Emph{埃兰}—— 我愿以奥妙的语言来听和回答。\par

\Emph{奥尔托}—— 你可知主称自己为葡萄树？\par

\Emph{埃兰}—— 是的，我知道祂说过\Quote{我是真葡萄树}。\par

\Emph{奥尔托}—— 那么葡萄树的果实经压榨后如何称呼？\par

\Emph{埃兰}—— 称为酒。\par

\Emph{奥尔托}—— 当士兵用长矛刺透救主肋旁时，福音记载从其中流出了什么？\par

\Emph{埃兰}—— 血和水。\par

\Emph{奥尔托}—— 那么，族长称救主的血为葡萄汁，便是恰当而适宜的，因为主被称为葡萄树，葡萄树的果实是酒，而从主肋旁涌出的血与水顺着祂全身流下。于是族长预言说：\Quote{他要在酒中洗自己的衣服，在葡萄汁中洗自己的外氅。}因为正如我们在祝圣后称葡萄树的奥妙果实为主的血，同样他也称真葡萄树的血为葡萄汁。\par

\Emph{埃兰}—— 眼前的问题已用既奥妙又清晰的言语陈述清楚了。\par

\Emph{奥尔托}—— 虽然这些话已足够坚固你的信德，但我还要为信德的坚固再提供另一个证据。\par

\Emph{埃兰}—— 你若如此做，我将感激不尽，因为你将加倍对我的恩惠。\par

\Emph{奥尔托}—— 你可知天主如何称自己的肉体为饼？\par

\Emph{埃兰}—— 是的。\par

\Emph{奥尔托}—— 又在另一处称自己的肉为麦粒？\par

\Emph{埃兰}—— 是的，我知道。因为我曾听祂说：\Quote{人子受光荣的时辰到了}，以及\Quote{一粒麦子如果不落在地里死了，仍只是一粒；如果死了，才结出许多籽粒来。}\par

\Emph{奥尔托}—— 是的；并且在分施奥迹时，祂称饼为身体，称调和的酒为血。\par

\Emph{埃兰}—— 祂确实如此做了。\par

\Emph{奥尔托}—— 然而按本性，身体自当称为身体，血自当称为血。\par

\Emph{埃兰}—— 同意。\par

\Emph{奥尔托}—— 但救主更换了名称，将象征的名称给予自己的身体，又将身体的名称给予象征。同样，在自称葡萄树之后，祂将象征称为血。\par

\Emph{埃兰}—— 不错。但我愿知道更换名称的原因。\par

\Emph{奥尔托}—— 对于领受神圣奥迹的人来说，意图是清楚的。因为祂愿参与神圣奥迹的人不留意可见之物的本性，而是藉着名称的变化，相信恩宠所成就的改变。因为我们知道，那位称自己的自然身体为麦粒和饼，又自称葡萄树的主，用身体和血的称呼尊崇了可见的象征，不是因为祂改变了它们的本性，而是因为祂向它们的本性添加了恩宠。\par

\Emph{埃兰}—— 奥迹以奥妙的言语讲述，对非所有人知晓的真理作了清晰宣告。\par

\Person{正统者}——既然大家都同意主的身体被圣祖称为\Quote{外衣}和\Quote{斗篷}，而我们已进入对神圣奥迹的讨论，请如实告诉我，你认为圣体是何种象征与预像？是主基督的天主性，还是祂的身体和祂的血？\par

\Person{异端者}——显然，是那些它们得名的事物。\par

\Person{正统者}——你是指身体和血？\par

\Person{异端者}——正是。\par

\Person{正统者}——你说话真像一个热爱真理的人。因为主拿起象征物时，祂并没有说\Quote{这是我的天主性}，而是\BibleQuote{这是我的身体}；又说\BibleQuote{这是我的血}；在另一处说\BibleQuote{我所要赐给的食粮，就是我的肉，是为世界的生命而赐给的}。\par

\Person{异端者}——这些话是真实的，因为它们是神圣的训示。\par

\Person{正统者}——如果它们是真实的，那么我认为主有一个身体。\par

\Person{异端者}——不；因为我主张祂是无身体的。\par

\Person{正统者}——但你不是承认祂有一个身体吗？\par

\Person{异端者}——我说圣言成了血肉，因为我所受的教导是这样。\par

\Person{正统者}——看来，正如俗语所说，我们好像在用有洞的桶打水。因为在所有的论证和困难解答之后，你又将同样的论点绕回来了。\par

\Person{异端者}——我不是在给你我的论点，而是福音书的论点。\par

\Person{正统者}——难道我没有从先知和宗徒的话中给你解释福音的言语吗？\par

\Person{异端者}——它们无助于澄清争论的要点。\par

\Person{正统者}——然而我们已表明，祂虽是不可见的，却藉着血肉彰显出来，而这血肉的关系，我们已从神圣作者那里学到：\BibleQuote{祂取了亚巴郎的后裔。}上主天主对圣祖说：\BibleQuote{地上万民都要因你的后裔蒙受祝福。}宗徒说：\BibleQuote{显然我们的主出自犹大支派。}我们还引用了其他几个类似的见证。但既然你渴望再听别的，那就听宗徒的话：\BibleQuote{因为每一位大司祭是从人间选立，被委派献上供物和祭品，因此这人也必须有所奉献。}\par

\Person{异端者}——那么，请指出祂取了身体之后是如何奉献的。\par

\Person{正统者}——神圣宗徒本人在同一段落中清楚地教导了；因为说了几句话后，他说：\BibleQuote{为此，基督进入世界时说：\InnerQuote{牺牲和素祭你都不想要，却为我预备了一个身体。}}他没有说\Quote{你变成了一个身体}，而是说\BibleQuote{你为我预备了一个身体}，并清楚表明身体的构造是由圣神完成的，符合福音的宣告：\BibleQuote{不要怕娶你的妻子玛利亚，因为那在她内受孕的是出于圣神。}\par

\Person{异端者}——那么，童贞女只生了一个身体？\par

\Person{正统者}——看来你连词语的构成都不懂，更不用说它们的含义了；因为他教导若瑟的是方式，不是产生的样式，而是受孕的样式。因为他并没有说那从她产生的——即造成或形成的——是出于圣神。若瑟不明白这个奥迹，怀疑是奸情；所以他清楚地被教导这形成是由圣神完成的。这正是天主藉先知所预表的，当时祂说\BibleQuote{你为我预备了一个身体}，因为充满圣神的神圣宗徒解释了这预言。如果奉献供物是司铎的特殊职能，而基督按人性被称为司铎，且祂所奉献的祭品除祂自己的身体外别无其他，那么主基督当然有一个身体。\par

\Person{异端者}——这一点连我也一再肯定，我并没有说圣言没有身体就显现了。我所主张的不是祂取了一个身体，而是祂成了血肉。\par

\Person{正统者}——依我看，我们的争辩针对的是瓦伦提努、马尔西翁和摩尼的支持者；但连他们也不敢说那不变的本性变成了血肉。\par

\Person{异端者}——辱骂是不合乎基督徒的。\par

\Person{正统者}——我们不是在辱骂，而是在为真理争辩；我们厌烦你竟然把无可争辩的事当作可争辩的。不过，我将努力结束你这种无理的争论。现在回答我：你记得天主向达味所作的许诺吗？\par

\Person{异端者}——哪些？\par

\Person{正统者}——先知插入第八十八篇圣咏的那些许诺。\par

\Person{异端者}——我知道向达味作了许多许诺。你现在问的是哪些？\par

\Person{正统者}——那些关乎主基督的。\par

\Person{异端者}——你自己回想那些话吧，因为你答应提出你的证据。\par

\Person{正统者}——现在请听，先知在圣咏的开头如何赞美天主。他用先知的眼睛看到了他百姓将来的不义，以及随之注定的被掳；但他仍然为自己的主而赞美那从不落空的许诺。他说：\BibleQuote{我要歌唱上主的仁慈直到永远，我的口要世世代代传扬你的信实；因为你曾说：\InnerQuote{仁慈将要永远建立，你的信实要在天上立定。}}\par

通过这些，先知教导我们，这许诺是由天主出于慈爱而作的，并且这许诺是信实的。然后他继续说出天主所许诺的是什么，对谁许诺，并让天主自己发言。（\BibleQuote{我与我所拣选的人立了盟约。}）祂称圣祖们为拣选的人；接着祂继续说：\BibleQuote{我向我的仆人达味起了誓}，并说明祂所发誓的内容：\BibleQuote{我要永远坚立你的后裔，世世代代建立你的宝座。}\par

现在你认为谁被称为达味的后裔？\par

\Person{异端者}——那许诺是关于撒罗满的。\par

\Person{正统者}——那么，祂是与圣祖们订立盟约关于撒罗满吗？因为在说到达味之前，祂提到了向圣祖们所作的许诺：\BibleQuote{我与我所拣选的人立了盟约。}祂应许圣祖们要在他们的后裔中祝福万民。请指出万民是如何通过撒罗满蒙受祝福的。\par

\Emph{厄冉。——}那么天主藉着我们的救主，而非藉着撒罗满，实现了这应许。\par

\Emph{正统。——}那么我们的主基督实现了给达味的应许。\par

\Emph{厄冉。——}我认为这些应许是天主关于撒罗满或则鲁巴贝耳所作的。\par

\Emph{正统。——}方才你引用了\Person{马尔西雍}、\Person{华伦提努}和\Person{玛尼斯}的论点。现在你又倒向了完全对立的阵营，为犹太人的狂妄辩护。这正像所有偏离正道的人；他们先是错失，然后迷途，东漂西荡，在荒野中徘徊。\par

\Emph{厄冉。——}辱骂者被宗徒排除于天国之外。\par

\Emph{正统。——}是的，如果他们的辱骂是枉然的。有时，那位神圣的宗徒自己也适时地使用这种言辞。他称迦拉达人\BibleQuote{无知}，论及其他人时又说\BibleQuote{心术败坏，背弃信德的人}，再论及另一等人时又说\BibleQuote{他们的天主是肚腹，以羞辱为光荣}，等等。\par

\Emph{厄冉。——}我给了你什么辱骂的理由？\par

\Emph{正统。——}你真的不认为，主动为公开的真理敌人辩护，给敬虔者提供了非常合理的愤怒理由吗？\par

\Emph{厄冉。——}我又资助过哪些真理的敌人？\par

\Emph{正统。——}此刻，犹太人。\par

\Emph{厄冉。——}何以见得？\par

\Emph{正统。——}犹太人将这类预言与撒罗满和则鲁巴贝耳联系起来，以表明基督徒立场的无根据；但单是这些言语就足以控告他们的不义，因为经上记载：\BibleQuote{我要永远坚立我的宝座}。如今，不仅撒罗满和则鲁巴贝耳——犹太人将这些预言应用于他们身上——度过了他们命定的时限，到达了生命的终点，而且达味的整个家族都已断绝；因为当今有谁曾听说过任何出自达味根株的人呢？\par

\Emph{厄冉。——}那么，那些被称为犹太人宗主教的人，难道不都是达味家族的吗？\par

\Emph{正统。——}当然不是。\par

\Emph{厄冉。——}那么，他们是从哪里出来的？\par

\Emph{正统。——}出自外邦人黑落德，他父系是阿市刻隆人，母系是依杜默雅人；但他们也全都消失了，他们的主权结束已有多年。然而我们的主天主不仅应许使达味的后裔永远存续，还要建立他永不毁坏的王国；因为祂说：\BibleQuote{我要建立我的宝座，直到万代}。\par

但我们看到他的家族已逝，王国也已终结。然而，尽管我们看到了这一点，我们知道宇宙的天主是真实的。\par

\Emph{厄冉。——}天主是真实的，这是明显的。\par

\Emph{正统。——}那么，如果天主是真实的，祂事实上就是，并且祂应许达味要永远坚立他的家族，使他的王国永存，而如今既不见家族也不见王国，因为两者都已终结，我们如何能说服我们的对手相信天主是真实的呢？\par

\Emph{厄冉。——}那么，我想这预言实际上是指向主基督。\par

\Emph{正统。——}那么，如果你承认这一点，让我们一同查考圣咏中间的一段；我们便会更清楚地看到这预言的含义。\par

\Emph{厄冉。——}请引导；我必虔诚地追随你的脚步。\par

\Emph{正统。——}在关于这后裔作出许多应许之后——说他应成为海陆之主，高于地上的众王，被称为天主的长子，并应大胆地称天主为父——天主又加上了这句话：\BibleQuote{我要永远为他保存我的仁慈，我的盟约与他坚定不渝。我也要使他的后裔永远存续，使他的宝座如天之永存。}\par

\Emph{厄冉。——}这应许超越了人性的界限，因为生命和尊荣都是不可朽坏、永恒的。但人只存续一时；他们的本性是短命的，他们的王国即使在有生之年也要经历许多多样的变迁，所以这预言的高超确实只适合于救主基督。\par

\Emph{正统。——}那么，继续看下文，你在这点上的看法将得到全面证实，因为宇宙的天主又说：\BibleQuote{一次我曾指我的圣洁起誓，我决不对达味说谎：他的后裔必永远存续，他的宝座在我面前如太阳；它必如月亮永远确立。}\par

然后，为了指出这应许的真实性，祂加上了：\BibleQuote{并且在天上的见证是忠信的。}\par

\Emph{厄冉。——}我们必须毫不怀疑地相信这位忠信见证人所给的应许，因为如果我们惯于相信那些应许说实话的人，即使他们没有用誓词证实自己的话，那么当宇宙的创造主在祂的话上加添了誓词时，谁还能疯狂到不信祂呢？因为那禁止他人起誓的，却以誓词确认了祂计划的不变性，\BibleQuote{好让我们这些逃去避难的人，抓住那摆在面前的希望，得着强有力的鼓励，藉着两件不可更改的事，在这些事上天主决不可能说谎。}\par

\Emph{正统。——}那么，如果这应许是不可反驳的，而在犹太人中间如今既不见达味先知的家族也不见他的王国，让我们相信我们的主耶稣基督按祂的人性明确地称为达味的后裔，因为祂的生命和国度都是同样永恒的。\par

\Emph{厄冉。——}我们毫不怀疑；我承认这就是真理。\par

\Emph{正统。——}这些证据足以清楚地显示我们的主救主从达味的后裔所取的人性。但为了藉多数见证人消除一切怀疑的余地，让我们听听天主如何经由先知依撒意亚的口提及给达味的应许。他说：\BibleQuote{我要与你们订立永久的盟约}，并且指明那立法者，他加上：\BibleQuote{就是达味确实的慈爱。}\par

既然祂向\Person{达味}作了这项承诺，又通过\Person{依撒意亚}发言，祂必定会实现此承诺。预言之后的内容与我所说的相符，因为祂说：\BibleQuote{看，我已立祂为万民的证人，为万民的领袖和司令。看，不认识你的民族要呼求你，不明白你的国民要奔向你。}这并不适用于任何\Person{达味}的后裔，因为按\Person{依撒意亚}所说，\Person{达味}的后裔中谁曾成为万民的统治者？又有什么民族在祈祷中曾呼求\Person{达味}的后裔为天主呢？\par

\Emph{厄兰.}——对于完全清楚的事，争辩是不合适的，这明确指的是主基督。\par

\Emph{正统.}——那么让我们转向另一个先知见证，听听同一位先知说：\BibleQuote{从叶瑟的树干必生出一根枝条，从他的根上必发出一根嫩芽。}\par

\Emph{厄兰.}——我认为这个预言是关于\Person{则鲁巴贝耳}的。\par

\Emph{正统.}——若你听到下文，就不会坚持你的意见。犹太人从未这样理解这一预言，因为先知继续说：\BibleQuote{上主的神将住在他身上，智慧和聪敏的神，超见和刚毅的神，明达和敬畏上主的神。}任何人都不会把这归于一个纯粹的人，因为即使是圣人，圣神的恩赐也是分开赐予的，正如神圣的宗徒所见证的：\BibleQuote{这人藉圣神得了智慧的言语，那人因同一圣神得了知识的言语，}等等。先知描述那从\Person{叶瑟}之根而出的那一位，拥有圣神的一切能力。\par

\Emph{厄兰.}——反驳这一点是纯粹的愚昧。\par

\Emph{正统.}——现在请听下文。你将看到一些超越人性的事，他接着说：\BibleQuote{祂不凭眼见审判，也不按耳闻断案；但祂要按正义审判穷人，以公平为地上的权贵断案；以口中的棍杖击打大地，以唇中的气息诛杀恶人。}这些预言中有一些是属人的，有一些是属神的。正义、真理、公平和审判中的正直，彰显了人性中的德性。\par

\Emph{厄兰.}——至此我们清楚地知道先知预言了我们救主基督的来临。\par

\Emph{正统.}——下文会更清楚地向你显示这解释的真实性。因为他继续说：\BibleQuote{豺狼将与羔羊同住}等等，由此他同时教导了生活方式的区别和信仰的和谐；经验证实了这个预言，因为富足者、贫穷者、仆人与主人、统治者和被统治者、士兵与公民以及掌握世界权杖的人，都在同一泉源中领受洗礼，都受教于同一教义，都被接纳于同一奥迹之筵席，每位信徒都享有同等的份额。\par

\Emph{厄兰.}——这样就证明了天主被言说。\par

\Emph{正统.}——不仅是天主，也是人。所以在这一预言的开头，他说一根枝条将从\Person{叶瑟}的根长出。然后在预言的结尾，他再次拾起开头的主题，因为他说：\BibleQuote{到那日，叶瑟的根必成为万民的旗帜，外邦人必寻求祂，祂的安息之所大有荣耀。}\Person{叶瑟}是\Person{达味}的父亲，誓言中的承诺是给\Person{达味}的。若先知只知道主基督是天主，就不会说祂是从\Person{叶瑟}长出的枝条。预言也预告了世界的改变，因为他说：\BibleQuote{大地必充满认识上主的知识，像海水覆盖海洋一样。}\par

\Emph{厄兰.}——我听到了先知的言论。但我急切想知道，宗徒的神圣团体是否也说出主基督按肉身是\Person{达味}的后裔。\par

\Emph{正统.}——你询问的信息非但不困难，反而非常容易。只要听宗徒之长宣告：\BibleQuote{达味既是先知，也知道天主曾向他发誓，要从他的腰腹中按肉身兴起基督，坐他的宝座；他预先看见了，就论及基督的复活，说祂的灵魂没有被遗弃在阴府，祂的肉身也没有见到腐朽。}\par

由此你可以看出，主基督按肉身出自\Person{达味}的后裔，祂不仅有肉身，也有灵魂。\par

\Emph{厄兰.}——还有哪位宗徒宣讲了这事？\par

\Emph{正统.}——单是伟大的\Person{伯多禄}就足以见证真理，因为主在接受了\Person{伯多禄}单独给出的真理的宣认后，以令人难忘的认可予以确认。但由于你渴望听到其他人宣讲同一件事，请听\Person{保禄}和\Person{巴尔纳伯}在\Place{丕息狄雅的安提约基雅}宣讲：当他们提到\Person{达味}时，接着说：\BibleQuote{天主从这人的后裔中，按照承诺，给以色列兴起了一位救主，耶稣}等等。在写给\Person{弟茂德}的书信中，神圣的\Person{保禄}说：\BibleQuote{你要记住耶稣基督，按我的福音，从达味的后裔中复活了。}在写给罗马人的书信中，他一开始就提到\Person{达味}的亲族，因为他说：\BibleQuote{耶稣基督的仆人保禄，蒙召为宗徒，被选为传天主的福音；这福音是天主先前藉着先知在圣经中所预许的，论及祂的儿子，按肉身是从达味的后裔生的，}等等。\par

\Emph{厄兰.}——你的证据很多且令人信服；但请告诉我，你为什么省略了后续内容？\par

\Person{Orth.}——因为使你感到困难的并非关于天主性，而是关于人性。假若你对天主性有所怀疑，我会向你提供证据。只要说\Quote{按肉体}就足以表明那不用术语表达的天主性。当论及一般人的亲属关系时，我不会说某人之子\Quote{按肉体}，而只说\Quote{子}。因此，那位神圣的福音作者在他的族谱中说\BibleQuote{亚巴郎生了依撒格}，并没有加上\Quote{按肉体}，因为他仅仅是凡人，他同样地提到其余的人，因为他们都是人，没有超越其本性的品质。但当真理的宣告者谈论我们的主基督，并向无知者指出祂较低卑的亲属关系时，他们会加上\Quote{按肉体}这句话，藉此指明祂的天主性，并教导说主基督不但是人，而且是永生的天主。\par

\Person{Eran.}——你已经从宗徒和先知引用了许多证据，但我遵循福音作者的话：\BibleQuote{圣言成了血肉}。\par

\Person{Orth.}——我也遵循这神圣的教导，但我以虔诚的意义理解它，意指祂藉取人性（包括肉体和理性的灵魂）而成为血肉。但如果神圣的圣言没有取我们本性中的任何东西，那么天主用誓言与圣祖们所立的盟约就不真实，犹大的祝福也是虚空的，对达味的应许是假的，童贞玛利亚也是多余的，因为她没有为我们本性贡献任何东西给降生成人的天主。那么先知的预言就没有应验。那么我们的宣讲是虚空的，我们的信德是虚空的，复活的望德也是虚空的，因为宗徒似乎说谎，当他说\BibleQuote{祂使我们一同复活，并使我们一同坐在天上基督耶稣内}。因为如果主基督没有我们本性的任何东西，那么祂被描述为我们的初果就是虚假的，祂的身体本性没有从死者中复活，也没有在天上坐在右边；如果祂没有得到这些，那么天主怎能祂使我们一同复活，并使我们一同与基督同坐，而我们与祂在本性上毫无关联呢？但说这话是不敬的，因为神圣的宗徒，尽管普遍的复活尚未发生，尽管天国尚未赐予信徒，却宣告说：\BibleQuote{祂使我们一同复活，并使我们一同坐在天上基督耶稣内}，为教导我们，既然我们的初果已经复活，祂已坐在右边，我们众人也都要获得复活，而且所有分享祂的本性并接受祂信仰的人，也分享祂荣耀的初果。\par

\Person{Eran.}——我们讨论了许多有力的论据，但我渴望知道这句福音话语的力量。\par

\Person{Orth.}——你不需要外来的解释。福音作者自己解释自己。因为说了\BibleQuote{圣言成了血肉}，接着就说\BibleQuote{并居住在我们中间}。也就是说，藉着居住在我们内，并使用了从我们取来的肉体如同圣殿，祂被称为成了血肉；并且为教导祂没有改变，福音作者又加上\BibleQuote{我们见过祂的荣耀——正如父独生者的荣耀，充满恩宠和真理}。因为虽然披着血肉，祂却彰显了祂父的尊贵，放射出天主性的光芒，发出了主大能的荣光，藉着祂奇妙的作为揭示了祂隐藏的本性。神圣的宗徒致斐理伯人的话也提供了类似的说明：\BibleQuote{你们该怀有基督耶稣所怀有的心情：祂虽具有天主的形体，却没有以自己与天主同等为应当把持的，反而空虚了自己，取了奴仆的形体，成为人的模样；既取了人的样式，祂就谦卑自下，服从至死，且死在十字架上。}\par

请看这些陈述的关系。福音作者说\BibleQuote{圣言成了血肉并居住在我们中间}，宗徒说\BibleQuote{取了奴仆的形体}；福音作者\BibleQuote{我们见过祂的荣耀，正如父独生者的荣耀}——宗徒说\BibleQuote{祂虽具有天主的形体，却没有以自己与天主同等为应当把持的}。简而言之，两者都教导：祂是天主、天主子，披戴着祂父的荣耀，与生祂者具有相同的本性和权能；祂起初就在，与天主同在，祂是天主，是世界的创造者，却取了奴仆的形体；人们所看到的似乎只有这些；但却是天主披上了人性，为人类施行救恩。这就是\BibleQuote{圣言成了血肉}和\BibleQuote{成为人的模样，既取了人的样式}所指的意思。所有这些就是犹太人所看见的，所以他们向祂说：\BibleQuote{我们不是为了善事而用石头砸你，而是为了亵渎，并且因为你是一个人却把自己当作天主。}，又说：\BibleQuote{这个人不是从天主来的，因为他并不守安息日。}\par

\Person{Eran.}——犹太人因不信而盲目，因此说了这些话。\par

\Person{Orth.}——如果你发现连宗徒们在复活之前也这样说，你会接受这个解释吗？我听见他们在船上，在平息风浪的大奇迹之后，说\BibleQuote{这是什么人？连风和海也听从祂！}\par

\Person{Eran.}——这已经清楚了。但现在告诉我这一点——神圣的宗徒说祂\BibleQuote{成为人的模样。}\par

\Emph{Orth.}— 祂所取的并非人的样式，而是人的本性。因为所谓\Quote{奴仆的形体}之理解，正如\Quote{天主的形体}被理解为天主之本性。祂取了此形体，遂成了人的模样，且以人的形态显现。因为祂原是神，却因所取的本性而看似为人。然而，福音作者称祂成了人的模样，亦即成了肉身。但为使你知晓那些否认救主肉身的人属于相反之神，请听伟大的\Person{若望}在他公函中所言：\Quote{凡明认耶稣基督是成了肉身而来的，便是出于天主；凡不明认耶稣基督是成了肉身而来的，便不是出于天主，这就是那假基督的精神。}\par

\Emph{Eran.}— 你的解释合乎情理，但我渴望知道教会的古老导师们如何理解\Quote{圣言成了肉身}这段话。\par

\Emph{Orth.}— 你本应被宗徒与先知之证所说服；但既然你还需要圣教父们的解释，我也将在天主助佑下为你提供这一药物。\par

\Emph{Eran.}— 不要为我引述名声不显或教义可疑之人。我不接受这类人的解释。\par

\Emph{Orth.}— 那么，声名远扬的\Person{亚大纳削}，\Place{亚历山大}教会最耀眼的光芒，在你看来是否值得信赖？\par

\Emph{Eran.}— 当然，因为他藉为真理所受的苦难印证了他的教导。\par

\Emph{Orth.}— 那么请听他致\Person{埃皮克提图}的书信：\Quote{\Person{若望}所说的\InnerQuote{圣言成了肉身}，其解释如下，据我们从圣\Person{保禄}类似的段落\InnerQuote{基督为我们成了诅咒}中所能发现的，并非因为祂成了诅咒，而是因为祂为我们接受了诅咒，故被称为成了诅咒；同样，并非因为祂变成了肉身，而是因为祂为我们取了肉身，故被称为成了肉身。}神圣的\Person{亚大纳削}如此说。此外，在所有人中享有盛誉的\Person{额我略}，他曾治理\Place{博斯普鲁斯}海峡口的帝国都城，后居于\Place{纳齐盎}，也曾致信\Person{克莱多尼乌斯}，驳斥\Person{阿波利纳里}似是而非的谬论。\par

\Emph{Eran.}— 他是一位杰出人物，是虔诚事业中最卓越的斗士。\par

\Emph{Orth.}— 那么请听他所说。他说：\Quote{\InnerQuote{祂成了肉身}这一表述，似乎与祂被称为\InnerQuote{成了罪}和\InnerQuote{成了诅咒}相似，并非因为主被转变成这些——祂怎么可能呢？——而是因为当祂承担我们的过犯并背负我们的病弱时，祂接受了这些。}\par

\Emph{Eran.}— 两种解释是一致的。\par

\Emph{Orth.}— 我们已向你展示了南方和北方的牧者在和谐一致；现在，让我们也介绍西方杰出的导师们，他们虽以另一种语言，却以同一心意写下了解释。\par

\Emph{Eran.}— 我听说\Person{盎博罗削}，他装饰了\Place{米兰}的主教宝座，在所有异端中首当其冲地战斗，写出了极优美且与宗徒教导相符的著作。\par

\Emph{Orth.}— 我将向你提供他的解释。\Person{盎博罗削}在他的著作\WorkTitle{论信德}中说：\Quote{经上记载：\InnerQuote{圣言成了肉身}。我不否认这记载，但请注意所用措辞；因为接着有\InnerQuote{并居住在我们中间}，即居于人的肉身中。因此，当经上论及圣言取肉身而成肉身时，你感到惊讶；同样，当论及祂所没有的罪，却说祂\InnerQuote{成了罪}时，并非说祂成了罪的本性与运作，而是要使祂在我们的肉身中钉死我们的罪。那么就让他们停止断言圣言的本性经历了变化与改变吧，因为取者为一，所取者为另一。}\par

现在，你应当聆听东方导师们的教导；这是东方唯一的区域，也是我们迄今为止唯一未提及的世界一角，尽管他们本可首先为真理作证，因为宗徒的教导首先传给了他们。但因你已在谎言磨刀石上磨快了舌头，攻击虔诚的长子们，我们便将他们留到最后，好让你先聆听其余见证，再将见证与见证并列，从而既钦佩他们的和谐，又终止你无休止的争论。那么，请听\Person{弗拉维亚诺}，他长期睿智地掌舵\Place{安提约基雅}教会，并使其引领的众教会安全度过\OriginalTerm{亚流}{Arian}风暴，为他们讲解了福音的话：\Quote{圣言成了肉身，居住在我们中间；祂并非变成了肉身，也未停止成为天主，因为祂从永远就是天主，并在降生成人的\OriginalTerm{救恩计划}{dispensation}中成了肉身，祂自己建造了自己的殿宇，并住在能受苦的受造物中。}若你愿听\Place{巴勒斯坦}先贤之言，请侧耳聆听令人钦佩的\Person{格拉西乌斯}，他在\Place{凯撒勒雅}教会辛勤耕耘。以下是他\WorkTitle{论主显节讲道}中的话：\Quote{从渔夫\Person{若望}的话中学习真理：\InnerQuote{圣言成了肉身}，祂自身并未经历变化，而是居住在我们中间。住所是一回事，圣言是另一回事；圣殿是一回事，住在其中的天主又是另一回事。}\par

\Emph{Eran.}— 我深深惊叹于这种一致。\par

\Emph{Orth.}— 难道你不认为\Person{若望}持守了宗徒信德的准则？他首先丰盛地浇灌了\Place{安提约基雅}人的教会田地，后来又成为\Place{帝国都城}教会的智慧耕耘者。\par

\Emph{Eran.}— 我认为这位导师在各方面都是令人钦佩的。\par

\Emph{正统者}：——那么，这位最杰出的人已经解释了福音的这段经文。他写道：\Quote{当你们听说圣言成了血肉时，不要惊惶或沮丧，因为本质并没有堕落为血肉——那是极其不敬的想法——而是继续保有它本来的样子，如此取了奴仆的形体。正如当宗徒说\BibleQuote{基督由法律的咒骂中赎出了我们，为我们成了咒骂}时，他并不是说基督的本质离开了自己的荣耀，取了咒骂的本质——这连魔鬼也不会想象，连最没有理智、天生愚钝的人也不会——不敬与疯狂的联系如此显著。他所肯定的乃是：在我们应得的咒骂被祂接受之后，祂不让我们将来再受咒骂。正是在这个意义上，才说祂成了血肉，不是因为祂把本质变成了血肉，而是因为祂取了血肉，本质始终不受损害。}\par

你们或许也愿意听听\Place{加巴拉}的主教\Person{塞维里亚诺}。若是如此，我将引用他的证言，请你们侧耳倾听。\par

\Quote{\BibleQuote{圣言成了血肉}这经文并不表示本性的退化，而是对我们本性的摄取。假设你们把\BibleQuote{成了}理解为变化，那么当你们听到保禄说\BibleQuote{基督由法律的咒骂中赎出了我们，为我们成了咒骂}时，你们是否理解为祂变成了咒骂的本性？正如\BibleQuote{成了咒骂}除了指祂亲自承担了我们的咒骂之外别无他意，同样\BibleQuote{圣言成了血肉，并居住在我们中间}也不过是指摄取了血肉。}\par

\Emph{收集者}：——我钦佩这些人的精确一致。他们对福音经文的解释如此一致，仿佛是在同一地点聚会并共同写下意见似的。\par

\Emph{正统者}：——山与海将他们远远分开，但距离并不损害他们的和谐，因为他们都由同一圣神的恩赐所感动。我本也应把胜利的虔诚斗士\Person{狄奥多若}和\Person{戴奥多若}的解释提供给你们，若不是看到你们对他们怀有敌意、并继承了\Person{阿波利纳里}的敌意的话；你们会看到他们表达了类似的体验，从神圣的源泉汲水，自身也成了圣神的活泉。但我将略过他们，因为你们对他们宣战不休。然而，我要向你们展示教会那位著名的教师对天主降生的看法，好让你们知道他对所取的本性持何种意见。无疑你们听说过杰出的\Person{依纳爵}，他由伟大的\Person{伯多禄}手中接受了主教恩宠，在治理\Place{安提约基雅}教会后戴上了殉道的冠冕。你们也听说过\Person{依勒内}，他得到了\Person{波利卡普}的教导，成了西部\Place{高卢}的光明；还有\Person{希玻里}和\Person{美多德}，主教和殉道者，以及其他许多人，我将把他们的名字附在各自的意见之后。\par

\Emph{收集者}：——我非常渴望听一听他们的证言。\par

\Emph{正统者}：——现在请听他们提出宗徒的教导。\Emph{圣\Person{依纳爵}的证言，\Place{安提约基雅}主教，殉道者。}\par

\Emph{取自\WorkTitle{致斯米纳人书}（第一章）：}——\par

\Quote{我们确信我们的主按肉体真正是由达味所生，按天主性和德能是天主子，真正由童贞女所生，由若翰受洗，好使一切义德由祂完成，在\Person{般雀·比拉多}和分封侯\Person{黑落德}之时，真正为我们在肉身内被钉十字架。}\par

\Emph{同一书信中同一位：}——\par

\Quote{若有人称赞我，却亵渎我的主，不承认祂是血肉的承担者，这对我有何益处？但谁不如此宣认，就是真正否认祂，并且自己成了尸体的承担者。}\par

\Emph{同一书信中同一位：}——\par

\Quote{因为如果这些事只是表面上由我们的主所做，那么我带上锁链也只是表面上的；我为何把自己交付给死亡、烈火、刀剑和野兽？但谁接近刀剑，就是接近天主。只有在耶稣基督的名内，我才忍受一切，为分享祂的苦难，而祂——完全的真人——给因无知而否认祂的人赐予我力量。}\par

\Emph{同一位自\WorkTitle{致厄弗所人书}：}——\par

\Quote{因为我们的天主耶稣基督，按天主的措施，由达味的后裔和圣神，在玛利亚的胎中被孕育，祂被生、受洗，好使我们的可朽性得以净化。}\par

\Emph{自同一书信：}——\par

如果你们每个人按恩宠，名符其实地聚集于同一信德和同一耶稣基督内，祂按肉体是达味族系的天主子、人子。\par

\Emph{同一书信中同一位：}——\par

\Quote{只有一位医师，血肉和神性的，受生和非受生的，在人之内的天主，在死亡中的真生命，玛利亚之子和天主之子，先可受苦后不可受苦，耶稣基督我们的主。}\par

\Emph{最后同一位在其\WorkTitle{致特拉里亚人书}中：}——\par

\Quote{因此，若有人对你们讲论耶稣基督之外的事，你们要成为聋子——祂是达味族系、玛利亚所生，真正出生，真正吃喝，在\Person{般雀·比拉多}时受迫害，被钉十字架而死，当时地上、天上和地下的生灵都在观看。}\par

\Emph{\Person{依勒内}主教的证言，\Place{里昂}主教，自其\WorkTitle{反对异端}第三卷：}——\par

\Quote{那么他们为何加上\BibleQuote{在达味城中}这些话？岂非为宣扬天主对达味所许诺的、从他的腰中所出的永远之王，已经应验；这许诺确实是世界的创造者所立的。}\par

\Emph{自同一书卷中同一位：}——\par

\Quote{当祂说\BibleQuote{达味家啊，请听}时，意思是天主应许达味要从他的身体兴起的那位永远之王，正是由达味的童贞女所生。}\par

\Emph{自同一书卷中同一位：}——\par

\Quote{如果第一个亚当有人的父亲，并由种子而生，那么说第二个亚当由若瑟所生就是合理的。但若前者从泥土中取来，其创造者是天主，那么那在自己内更新天主所创造的人的那一位，也必须与前者有相同的生育样式。那么天主为何不再取泥土？祂又为何预定要从玛利亚形成肉身？为的是不再有另一个受造物；为的是那被拯救的不再是别的；而是使前者自己得以更新，而不失其样式。因为那些声称祂从童贞玛利亚未取任何东西的人，也正因此跌倒，他们否认血肉的继承，抛弃了这样式。}\par

\Emph{同一位作者，同书：}—\par

\Quote{因为祂降到玛利亚那里是徒劳的；若祂打算不从她那里取得任何东西，为何要降到她里面？再者，若祂从玛利亚未取任何东西，祂就不会接受取自泥土的食物，因为取自泥土的身体是靠这食物滋养的；祂也不会像梅瑟和厄里亚那样，禁食四十天后感到饥饿，因为祂的身体需要自己的食物；祂的门徒若望也不会在写关于祂的事时说—\Quote{耶稣因旅途疲倦，就坐下}；达味也不会预言关于祂的\Quote{他们加重了我伤口的疼痛}；祂也不会为拉匝禄哭泣；祂也不会滴下血汗；祂也不会说\Quote{我的心神忧闷得要死}；当祂被刺穿时，血和水也不会从祂的肋旁流出。因为这一切都是取自泥土之肉身的明证，这肉身在祂自己内为拯救祂自己的受造物而更新了。}\par

\Emph{同一位作者，同书：}—\par

\Quote{因为正如因那一个从粗劣泥土中首先被形成的人的悖逆，许多人成了罪人并失去了生命；同样，也适宜因一个人的服从，即一位童贞女的首生子，许多人应被成义并领受他们的救恩。}\par

\Emph{同一位作者，同作品：}—\par

\Quote{\Quote{我曾说你们是神，都是至高者的儿子；但你们要死，如同世人一样。}这话是对那些不接受义子恩赐、反而羞辱圣言那纯洁诞生的降生成人、剥夺人上升至天主的途径、并忘恩于那为他们而成为血肉的圣言的人说的。为此，圣言成为人，为使人领受圣言并接受义子身份，而成为天主的儿子。}\par

\Emph{同一位作者，同书：}—\par

\Quote{既然因预定的计划，圣神降下，天主的独生子，也是圣父的圣言，在时期圆满时，在人间成了血肉，我们的主耶稣基督——同一位——完成了全部人的计划，正如主自己作证，宗徒们宣认，所有那些发明了八度、四度及相似之说的人的教导，都显明是虚假的。}\par

\Emph{圣依玻理，主教兼殉道者的见证，出自他关于} \Emph{\BibleQuote{上主是我的牧者}} \Emph{的讲道：}—\par

\Quote{救主自己就是一块不朽坏木头的约柜，因为祂的居所的不朽坏和不毁坏，意味着它不产生罪的腐朽。因为承认自己罪的罪人说：\Quote{因我的愚妄，我的伤口发臭流脓。}但主是没有罪的，在祂的人性中由不朽坏的木头所造，即由童贞玛利亚和圣神所生，仿佛里外都镀上了圣言的最纯之金。}\par

\Emph{同一位作者，出自他关于厄耳卡纳和亚纳的讲道：}—\par

\Quote{那么，撒慕尔，请给我那被牵到白冷的母牛，好使你显示那由达味所生、并由圣父傅油为王和司祭的君王。}\par

\Emph{出自同一讲道：}—\par

\Quote{请告诉我，荣福玛利亚，你在胎中所怀的是什么？你在童贞女之胎中所负的是什么？那是天主的圣言，从天而来的首生者，降临在你身上，而人在胎中作为首生者被形成，为使首生的圣言与首生的人得以结合而彰显。}\par

\Emph{出自同一讲道：}—\par

\Quote{第二次，即通过先知如撒慕尔的那次，祂撤销了，并使祂的子民从外邦人的奴役中回转。第三次，即祂取了童贞女的人性，亲自在肉身中临在；祂看见那城时，就为它哭泣。}\par

\Emph{同一位作者，出自他关于依撒意亚开端的讲道：}—\par

\Quote{祂把世界比作埃及；它的偶像崇拜比作偶像；它的除灭和毁灭比作地震。祂称圣言为\Quote{主}，而\Quote{快云}意指那正直的纯洁居所，我们的主耶稣基督驾临其上进入世界，以消除堕落。}\par

\Emph{圣默多狄，主教兼殉道者的见证，出自他关于殉道者的讲道：}—\par

\Quote{殉道是如此奇妙而珍贵，以致我们的主耶稣基督自己，天主子，为此作证，祂不以自己与天主同等为抢夺的，为使祂所降入的人性能以此恩宠得冠冕。}\par

\Emph{圣欧斯塔基，安提约基雅主教、精修者的见证，出自他对第十六篇圣咏的诠释：}—\par

\Quote{耶稣的灵魂经历了二者。因为它处于人的灵魂的地位，并且离开肉身后仍存活并持续。所以它是理性的，与人的灵魂同质，正如肉身与人的肉身同质，出自玛利亚。}\par

\Emph{同一位作者，出自他关于灵魂的作品：}—\par

\Quote{当看到那孩子的教养，或他身量的增长，或时间的延长，或身体的成长时，他们会说什么？但暂且不论在地上所行的奇迹，让他们看看死人复活、苦难的记号、鞭打的痕迹、青肿与击打、被刺的肋旁、钉痕、血流出、死亡的证据，总之，那同一个身体的实在复活。}\par

\Emph{出自同一作品：}—\par

\Quote{确实，若有人注视祂身体的诞生，便会清楚发现，祂在伯利恒出生后，被裹在襁褓里，因凶恶的希律的恶谋，被带到埃及一段时间，并在纳匝肋长大成人。}\par

\Emph{同书：}—\par

\Quote{因为圣言和天主的帐幕并非同一个，有福的斯德望借此看见了天主的荣耀。}\par

\Emph{同一位作者在其讲道\Quote{主在祂道路的开端创造了我}中：}—\par

\Quote{如果圣言从祂经过母亲的子宫穿上人形时才开始祂的受生，那么显然祂是由女人所生；但如果祂起初就是圣言，是偕同圣父的天主，并且我们断言宇宙是由祂造成的，那么祂——即那自有、是万有之因的——并非由女人所生，而是本性为天主，自有的、无限的、不可理解的；由女人所生的是人，在童贞女的胎中由圣神形成。}\par

\Emph{同书：}—\par

\Quote{因为一座绝对圣洁无玷的帐幕，就是圣言按肉体居住的帐幕，天主在其中有形可见地居住并住下，我们如此断言并非出于猜测；因为那位本性是天主圣子的，在预言圣殿的毁灭与复活时，用祂的教导明确训示我们，当祂对凶残的犹太人说：\InnerQuote{你们拆毁这座圣殿，三天之内我要把它重建起来。}}\par

\Emph{同书：}—\par

\Quote{当圣言建造了一座殿并取人性时，祂以身体与人共处，无形地展示了各种奇迹，并派遣宗徒作为祂永恒王国的报信者。}\par

\Emph{同一位作者，在其对第九十二篇圣咏的诠释中：}—\par

\Quote{显然，如果\Quote{那傅油者}指天主，祂的宝座被称为\Quote{永远的}，那么傅油者显然是天主，由天主所生。但受傅者领受了一种获得的德性，因一个被选为天主性居所之殿的装饰而美丽。}\par

\Emph{圣亚大纳削，亚历山大主教及精修者的证词。取自亚历山大主教狄约尼削的辩护：}—\par

\Quote{\InnerQuote{我是葡萄树，你们是枝条。我父是园丁。}因为我们在身体上与主有血缘关系，因此祂自己说：\InnerQuote{我要向我的弟兄宣扬你的名。}正如枝条与葡萄树同质且出于它，同样我们，由于拥有与主的身体相似的身体，从祂的圆满中领受它们，并以祂作为我们复活与救恩的根。父被称为园丁，因为祂自己藉着圣言栽培了作为葡萄树的主的身体。}\par

\Emph{同一位作者，出自同一论著：}—\par

\Quote{主因祂与枝条（即我们）的身体关系而被称为葡萄树。}\par

\Emph{同一位作者，出自其论信仰的大公演说：}—\par

\Quote{圣经\Quote{在起初已有圣言}明确指明天主性。经文\Quote{圣言成了血肉}显示了主的人性。}\par

\Emph{出自同一讲道：}—\par

\Quote{\InnerQuote{祂在酒中洗净自己的衣服}，即祂的身体，那是天主性的衣服，在祂自己的血中。}\par

\Emph{同一位作者，出自同一讲道：}—\par

\Quote{动词\Quote{存在}指祂的天主性，\Quote{成了血肉}指祂的身体；圣言成了血肉，并非因缩减成血肉，而是因承载了血肉，正如有人会说某人变成或成了老人，虽然起初并非如此出生，或士兵成了老兵，先前并非如此。若望说：\InnerQuote{我成为或曾在帕特摩岛上，在主日。}并不是他在那里产生或出生，而是他说\InnerQuote{我成为或在帕特摩}代替说\InnerQuote{我到达}；同样，圣言\InnerQuote{到达}了血肉，正如经上说\InnerQuote{圣言成了血肉}。请听\InnerQuote{我成了一个破漏的器皿}和\InnerQuote{我成了一个无力的人，在死者中自由。}}\par

\Emph{同一位作者，出自其致埃皮克特托书：}—\par

\Quote{谁曾听过这样的事？谁教导过它们？谁学习过它们？\InnerQuote{法律将出自熙雍，上主的话将出自耶路撒冷。}但这些事从何而来？什么地狱呕吐出它们？说由玛利亚取来的身体与圣言的天主性同质，或圣言变成了血肉、骨头、头发和整个身体；有谁曾在教会中或在基督徒中听说过，天主是靠着收养而非靠着出生而拥有身躯？}\par

\Emph{同一位作者，出自同一书信：}—\par

但是，谁听到说圣言为自己造了一个可受苦的身体，并非出自玛利亚，而是出自祂自己的本体，还会称呼说这话的人为基督徒？谁发明了如此无根据的不敬之罪，以至于甚至想和说：那些确认主的身体出自玛利亚的人，不再设想一个天主圣三，而是在天主性中设想一个四位一体？仿佛持此意见的人将救主由玛利亚所穿上的血肉描述成出自天主圣三的本体。\par

\Quote{此外，从哪里有些人呕吐出与前一种同样恶劣的不敬之罪，并断言身体并非晚于圣言的天主性，而是始终与它同永存，因为它是从智慧的实体形成的。}\par

\Emph{同一位作者，出自同一书信：}—\par

\Quote{因此由玛利亚取来的身体按照圣经是人性的，且是真实的，因为它与我们的相同。因为玛利亚是我们的姊妹，既然我们都出自亚当，凡记得路加的话的人都不能怀疑这事实。}\par

\Emph{圣巴西略，凯撒勒雅主教的证词：}—\par

\Emph{出自对第六十篇圣咏的诠释。}\par

\Quote{所有外邦人都屈身并服从了基督的轭，因此\InnerQuote{祂向厄东抛出了祂的鞋。}}\par

\Emph{同一位作者，出自其论圣神致安菲洛基的书信：}—\par

\Quote{祂用了\InnerQuote{由她}一词，而非\InnerQuote{透过她}；正如保禄所说的\BibleQuote{由女人所生}。祂在另一处清楚为我们作了这种区分，说：属于男人是由女人所生，而属于女人是由男人所成，正如经上所说\BibleQuote{因为女人原是由男人而出，男人也是由女人而生}。但祂既指明这些表达的不同用法，又附带纠正某些人以为主的身体是属神的错误，为要显示承载天主的血肉是由人的质料所成，便突出了更强调的表达；因为\InnerQuote{透过女人}的表达有危险暗示受生的意义仅仅是经过的过程，而\InnerQuote{由女人}则清楚表明孩子与母亲的共同本性。}\par

\Emph{圣额我略·纳齐安主教之见证。取自先前致克肋多尼的论述：}—\par

\Quote{若有人说那血肉是从天降下，而非来自这大地并出于我们，愿他受绝罚。因为\BibleQuote{第二个人来自天上}、\BibleQuote{那属天的怎样，凡属天的也怎样}、\BibleQuote{没有人升过天，除了那自天降下的人子}以及其他类似段落，必须理解为因着与人的结合而说的；同样，\BibleQuote{万有都是藉基督而受造}和\BibleQuote{基督住在我们心里}也必须理解，不是按照可感觉的，而是按照对神性理智的领悟，各词相互混合，如同本性混合一般。}\par

\Emph{同一位作者，出自同一著作：}—\par

\Quote{让我们从他们自己的话来看，他们如何解释天主成为人（即降生成人）的缘由。若果真天主（否则不受空间限制）是为了受空间限制，并好像蒙着面纱在肉身中与人交谈，那么他们的面具和戏剧倒是精巧；更不必说，祂本可以像古时在燃烧的荆棘中和以人的形象那样，以其他方式与我们交谈。但若祂是为了藉相似者解除罪的诅咒，那么正如祂因着被判罪的血肉而需要血肉，因着灵魂而需要灵魂，同样祂也因着理智而需要理智，这理智在亚当身上不仅堕落了，而且——借用医生论疾病所用的术语——患了原病。因为那没有遵守诫命的，正是那接受了诫命的；那胆敢犯禁的，正是那没有遵守诫命的；那特别需要救恩的，正是那犯了罪的；那被取纳的，正是那需要救恩的；因此理智被取纳了。现在这一点已经用近乎数学和必然的论证证明了，无论他们愿意与否。但你们所做的，正好像一个人眼睛有疾、脚又瘸，你却只治脚而不治眼；或者，一个画家画了一幅糟糕的画，你却修改画作而不管画家本人，好像他画得很好。但若他们被这些论证所迫，以致躲进\InnerQuote{没有理智，天主也能救人}的说法，那么显然，祂甚至没有血肉也能做到，只凭祂的旨意，正如祂在宇宙中无身体地工作并已经工作那样。那么，连血肉也除掉吧，连同理智一起！不要让你的荒谬自相矛盾。}\par

\Emph{圣额我略·尼撒主教之见证。取自关于亚巴郎的讲道：}—\par

\Quote{因此圣言降下并非赤裸，而是在成为血肉之后，不是以天主的形象，而是以仆人的形象。这就是那位曾说\BibleQuote{我凭自己不能做什么}的一位。因为不能做是软弱的一部分。因为正如黑暗与光相对，死亡与生命相对，软弱也与能力相对。然而基督却是天主的德能。能力与不能做完全不相容。因为若能力无能为力，那么什么是有能力的呢？因此当圣言宣称祂不能做什么时，显然祂不是将祂的不能归因于独生子的神性，而是将祂的不能与我们的本性的软弱联系起来。血肉是软弱的，正如经上所记\BibleQuote{心神固然愿意，肉体却软弱}。}\par

\Emph{同一位作者，出自其著作《论生命的成全》：}—\par

\Quote{再者，那真正的立法者（梅瑟是祂的预象）用我们的泥土为自己凿出了本性的石版。没有婚姻为祂塑造那将要承受神性的血肉，而是祂自己成为自己血肉的雕刻者，这血肉由天主的手指所刻。因为圣神临于童贞女，至高者的能力荫庇了她。当这成就之后，本性再次恢复其不朽的特质，因着天主手指的印记而成为不朽的。}\par

\Emph{同一位作者，出自其反对欧诺弥的著作：}—\par

\Quote{因此我们断言，当祂在上面说\BibleQuote{智慧为自己建造了房屋}时，这话暗示了主血肉的成形；因为智慧本身没有居住在陌生的居所，而是用童贞女的身体为自己建造了居所。}\par

\Emph{同一位作者，出自同一论著：}—\par

\Quote{圣言是在万世之前，而血肉是在末期造成；没有人会反过来主张血肉在万世之前，或圣言在末期造成。}\par

\Emph{同一位作者，出自同一论著：}—\par

\Quote{\BibleQuote{创造了我}这个表达不应理解为指神性和无玷者，而是如所说过的，指按照降生成人的安排，我们的受造本性。}\par

\Emph{同一位作者，出自《论真福》第一讲：}—\par

\Quote{\BibleQuote{祂虽具有天主的形体，却没有以自己与天主同等为应当把持不舍的，反而空虚了自己，取了仆人的形体。}就天主而言，有什么比仆人的形体更贫穷？就万王之王而言，有什么比接近我们卑微本性的共融更卑微？万王之王、万主之主自愿穿上了奴仆的形体。}\par

\Emph{圣弗拉维安·安提阿主教之见证。取自关于洗者若翰的讲道：}—\par

\Quote{不要以任何肉体的方式来思考结合，也不要抱有夫妻交合的想法。因为你的创造者正在创造祂自己由你而生的身体圣殿。}\par

\Emph{同一位从他的著作\WorkTitle{主的圣神在我身上}中：}—\par

\Quote{你们听祂说：\InnerQuote{圣神在我身上，因为祂用油膏了我。}祂说：你们不知道你们所读的是什么，因为我，被圣神所膏的那一位，来到了你们中间。然而，与我们相似的，而不是那不可见的本性，是被圣神所膏的。}\par

\Emph{\Person{安菲洛基乌斯}，\Place{依科尼雍}主教，从他的讲道\WorkTitle{圣父比我大}中的见证：}—\par

\Quote{现在请为我区分两个本性：神圣的本性和人的本性。因为人并非因堕落而出于天主，天主也非因提升而成为人。我是在谈论天主和人。然而，当你把苦难归于肉体，把奇迹归于天主时，你就必然且不自觉地卑微的称号归于那由玛利亚所生的人，而崇高的、神圣的称号归于那在起初就与天主同在的圣言。因此，在某些情况下我说出崇高的话，在另一些情况下我说出卑微的话，为借崇高的话显明内住的圣言的本性，借卑微的话承认卑微肉体的软弱。因此，有时我称自己与圣父平等，有时又说自己比圣父大，并非自相矛盾，而是表明我是天主也是人；因为天主属于崇高，人属于卑微；但若你想知道我的圣父如何比我大，我是就肉体而言，而非就神性的位格而言。}\par

\Emph{同一位从他的讲道\WorkTitle{圣子什么也不能凭自己作}中：}—\par

\Quote{亚当如何在天上不服从？他又如何是那天上身体的最初形成？但起初形成的是地上的亚当；地上的亚当不服从；地上的亚当被取了。因此，地上的亚当得救了，从而降生成人的理由得以显为必要和真实。}\par

\Emph{\Person{圣若望}，\Place{君士坦丁堡}主教，当哥特使节在他面前发言时所作的演讲中的见证：}—\par

\Quote{看祂从起初做了什么。祂穿戴了我们软弱且被征服的本性，为借它来作战和奋斗，并从起初就根除了反叛的本性。}\par

\Emph{同一位从他的讲道\WorkTitle{圣诞庆节}中：}—\par

\Quote{因为难道不是极其愚蠢的事：他们把自己的神明降到石头和廉价木像中，仿佛把它们关在某种监狱里，并幻想自己所说或所做的一切毫无可耻之处，然后指责我们说天主为自己用圣神造了一个活的圣殿，借此祂给世界带来了救助？因为如果天主住在人的身体中是可耻的，那么，比起人来，石头和木头更无价值，祂住在石头和木头中岂不更可耻？但也许在他们看来，人类比这些无感觉的对象更无价值。他们把天主的本体降到石头中，甚至降到狗中；但许多异端者却降到比这些更污秽的东西中。然而我们甚至无法容忍听到这些事。但我们所说的是：基督从童贞女的子宫中取了纯净、神圣、无玷的肉体，使之不受任何罪的侵染，并恢复了祂自己的身体。}\par

稍后又说：\Quote{我们断言，神圣的圣言为自己塑造了一座圣殿，借此将天上的状态带入我们的生命。}\par

\Emph{同一位从他的演讲\WorkTitle{论基督卑微的言行并非出于能力的缺乏，而是出于经世之区别}。}\par

\Quote{那么，关于祂，无论是祂自己还是祂的宗徒们，为何说了许多卑微事情的原因是什么？第一个也是最大的原因是祂穿戴了肉体的事实，并希望祂所有同时代的人以及此后所有活着的人都相信：祂不是阴影，所见的也不仅仅是形式，而是本性的现实。因为如果当祂自己和祂的宗徒们如此频繁地用卑微的、人间的意义谈论祂时，魔鬼还能说服一些可怜可悲的人否认降生成人的理由，胆敢说祂没有取肉体，从而摧毁祂爱人情的全部基础，那么，如果祂从未说过类似的话，有多少人不会陷入这个深渊呢？}\par

我现在已从真理宣讲者的众多权威中为你引证了几条，以免用太多使你震惊。这些已足以显示这些杰出作者的思想倾向。现在该由你来说出他们的著作似乎具有什么力量。\par

\Emph{Eran.}— 他们彼此和谐一致，西方葡萄园的工人也同意他们，这些工人的耕作在于日出之地。然而我察觉到他们的言论中有相当大的差异。\par

\Emph{Orth.}— 他们是神圣宗徒的继承者；甚至一些宗徒有幸听到了那神圣的声音，看到了那美好的景象。他们中的大多数还戴上了殉道者的冠冕。你认为向他们挥舞亵渎的舌头是合理的吗？\par

\Emph{Eran.}— 我对此退避三舍；但同时我不赞成他们的巨大分歧。\par

\Emph{Orth.}— 但现在我将为你带来一个意想不到的补救办法。我将引证你们自己美丽的异端中的一位——你们的老师\Person{阿波利纳里}，我将向你们展示他理解经文\BibleQuote{圣言成了血肉}，正如圣教父们一样。现在请听他关于此事在\WorkTitle{摘要}中所写的。\par

\Emph{\Person{阿波利纳里}从他的\WorkTitle{摘要}中的见证：}—\par

\Quote{如果没有人变成他所取的，而基督取了肉体，那么祂就没有变成肉体。}\par

随后他接着说：—\par

\Quote{因为祂也借着身体将自己赐给我们，为与我们建立关系以拯救我们。那拯救的远比那被拯救的更卓越。那么，在取肉体内，祂远比我们更卓越！但如果祂变成了肉体，就不会更卓越了。}\par

稍后他说：—\par

\Quote{单纯者是单一的，但复合者不能是单一的；因此，声称祂成了血肉的人，便是在肯定独一圣言的变化。但如果复合者也是单一的，如同人，那么因着与血肉的结合而说圣言成了血肉的人，指的是复合中的一体。}\par

\Quote{成了血肉就是成了空虚，但成了空虚所宣告的不是人，而是人子，祂\BibleQuote{空虚了自己}\BibleRef{斐 2:7}，不是通过变化，而是通过穿戴。}\par

\Quote{我们相信祂成了血肉，同时祂的天主性保持不变，为的是更新人性。因为在天主的神圣权能中，既没有变化，也没有位置的改变，更没有包容}——随后不久——\Quote{我们朝拜那从荣福童贞女取了血肉的天主，因此，在血肉中祂是人，但在灵里是神。}在另一篇阐述中他说——\Quote{我们宣认天主子成了人子，不是名义上，而是真实地，因从童贞玛利亚取了血肉。}\par

\Emph{Eran.}——我没有想到\Person{亚坡里纳尔}持有这些见解。我对他另有看法。\par

\Emph{Orth.}——好了，现在你已得知，不仅先知、宗徒以及他们之后被立为世界导师的人，就连\Person{亚坡里纳尔}——那个异端胡言乱语的作者——也宣认神圣的圣言是不变的，声明祂没有变成血肉，而是取了血肉，并且你已听到他反复这样说。所以，不要试图用你自己的话来掩盖你师傅的亵渎。因为主说：\BibleQuote{门徒不得高过师傅。}\par

\Emph{Eran.}——是的，我承认天主的圣言是不变的，并且取了血肉。抗拒如此众多且伟大的权威是极其愚蠢的。\par

\Emph{Orth.}——你愿意探讨其余难题的解答吗？\par

\Emph{Eran.}——让我们把考查推迟到明天。\par

\Emph{Orth.}——很好；我们的会议就此解散。我们离开吧，并记住我们所达成的一致。\par

\SourceRecord{Translated by Blomfield Jackson.；Nicene and Post-Nicene Fathers, Second Series；Vol. 3.；Edited by Philip Schaff and Henry Wace.；Buffalo, NY: Christian Literature Publishing Co.,；1892.；<http://www.newadvent.org/fathers/27031.htm>.}

% structure: p0001 source=h1 target=section reason=page-title

\section{对话二}

\Emph{厄拉尼斯特和正统者.}\par

\Emph{厄拉.}——我已如约而来。你面临两种选择：要么解答我的疑难，要么同意我和我朋友们所提出的。\par

\Emph{正.}——我接受你的挑战，因为我认为这是正确而公平的。但我们必须先回想昨天我们讨论到何处，以及我们论证的结论是什么。\par

\Emph{厄拉.}——我来提醒你结论。我记得我们一致同意：神圣的圣言保持不变，取了肉身，祂自己并未变成肉身。\par

\Emph{正.}——你似乎对已同意的要点感到满意，因为你忠实地回想起了它们。\par

\Emph{厄拉.}——是的，而且我已经说过，抗拒如此众多且伟大的教师的人无疑神志不清。此外，我颇感羞愧地发现，\Person{阿波利纳里}使用了与正统者相同的术语，尽管他在关于降生成人的著作中，其意图显然朝向另一个方向。\par

\Emph{正.}——那么，我们确认神圣的圣言取了肉身？\par

\Emph{厄拉.}——是的。\par

\Emph{正.}——那么，\Quote{肉身}指什么？仅指身体，如\Person{亚略}和\Person{欧诺米}所主张的，还是指身体和灵魂？\par

\Emph{厄拉.}——身体和灵魂。\par

\Emph{正.}——什么样的灵魂？理性的灵魂，还是某些人所称的\OriginalTerm{植物的}{phytic}、生长的，即生命的？因为阿波利纳里派编造谎言的骗术迫使我们问些不相宜的问题。\par

\Emph{厄拉.}——那么\Person{阿波利纳里}区分灵魂吗？\par

\Emph{正.}——他说人由三部分组成：身体、生命灵魂，以及理性灵魂（他称之为理智）。相反，圣经只知道一个灵魂，而非两个；这在第一个人的形成中清楚教导我们。因为经上记载：天主从泥土取了尘土，\BibleQuote{造了人}，又\BibleQuote{在他鼻孔内吹了生气，人就成了一个有灵的生物。}在福音中，主对圣门徒说：\BibleQuote{你们不要害怕那些杀害身体而不能杀害灵魂的；但更要害怕那能把灵魂和身体都毁灭在地狱里的。}\par

而神圣的\Person{梅瑟}在叙述下到埃及的人并说明每个族长带同哪些人时，补充说：\BibleQuote{从埃及出来的所有灵魂共七十五个，}为每位移民计一个灵魂。而神圣的宗徒在\Place{特洛阿}，当众人都以为\Person{厄乌提苛}死了时，说：\BibleQuote{你们不要惊慌，因为他的灵魂还在他身上。}\par

\Emph{厄拉.}——这清楚表明每个人只有一个灵魂。\par

\Emph{正.}——但\Person{阿波利纳里}说有两个；并说神圣的圣言取了无理性的灵魂，而代替理性的，祂在肉身中成了。正是为此我问你，你断言与身体一同被取的是什么样的灵魂。\par

\Emph{厄拉.}——我说是理性的。因为我随从神圣的圣经。\par

\Emph{正.}——那么，我们同意神圣的圣言所取的\Quote{奴仆的形体}是完整的。\par

\Emph{厄拉.}——是的，完整的。\par

\Emph{正.}——而且合理；因为整个第一个人成了罪的奴仆，失去了天主肖像的印记，而全人类随从，结果，造物主意欲更新模糊的肖像，取了整个本性，并印上了远优于第一个的印记。\par

\Emph{厄拉.}——确实。但现在我首先请求你，使所用术语的意义十分清楚，这样我们的讨论可以无阻碍地进行，没有对可疑点的探究插入打断我们的谈话。\par

\Emph{正.}——你说得好。现在请问你喜欢的任何问题。\par

\Emph{厄拉.}——我们该如何称呼耶稣基督？人？\par

\Emph{正.}——不能单用任何一个名称，而要用两者。因为那位神人成为人后，被称为耶稣基督。\Quote{因为}经上记载：\BibleQuote{你要给祂起名叫耶稣，因为祂要把自己的民族从他们的罪恶中拯救出来，}以及\BibleQuote{今天在达味城中为你们诞生了主基督。}这些是天使的声音。但在降生成人之前，祂被称为天主、天主子、独生子、主、圣言和创造者。因为经上记载：\BibleQuote{在起初已有圣言，圣言与天主同在，圣言就是天主，}以及\BibleQuote{万物是藉着祂造成的，}以及\BibleQuote{祂是生命，}以及\BibleQuote{祂是那普照每人的真光。}还有其他类似经文，宣告天主性。但在降生成人之后，祂被称为耶稣和基督。\par

\Emph{厄拉.}——因此主耶稣只是天主。\par

\Emph{正.}——你听见圣言成了人，难道你只称祂为天主吗？\par

\Emph{厄拉.}——既然祂成了人而没有改变，仍保持祂先前之所是，我们就必须称祂为祂先前之所是。\par

\Emph{正.}——圣言过去、现在和将来都是不变的。但当祂取了人的本性时，祂成了人。因此我们应该宣认两个本性：既宣认那取者，也宣认那被取的。\par

\Emph{厄拉.}——我们必须以更尊贵的名称呼祂。\par

\Emph{正.}——人——我指作为动物的人——是单纯体还是合成体？\par

\Emph{厄拉.}——合成体。\par

\Emph{正.}——由哪些组成部分构成？\par

\Emph{厄拉.}——由身体和灵魂。\par

\Emph{正.}——在这两个本性中，哪一个更尊贵？\par

\Emph{厄拉.}——显然是灵魂，因为它是理性的、不朽的，并被委托掌管那动物的主权。但身体是有死、可朽的，没有灵魂时是无理性的，是一具尸体。\par

\Emph{正.}——那么圣经应该按它更优越的部分来称呼那动物。\par

\Emph{厄拉.}——它确实这样称呼，因为圣经称那些从埃及出来的人为灵魂。因为\BibleQuote{以七十五个灵魂，以色列下到埃及。}\par

\Emph{正.}——但圣经从未按身体称呼任何人吗？\par

\Emph{厄拉.}——它称那些是肉身的奴仆的人为肉身。因为经上记载：\BibleQuote{天主说：我的神不会永远留在这些人里面，因为他们是肉身。}\par

\Emph{正.}——但没有人不受责备被称为肉身吗？\par

\Emph{厄拉.}——我不记得。\par

\Person{正统派}——那么，我要提醒你，并向你指出，就连圣人们也被称为\Quote{肉身}。现在请回答：你会怎么称呼宗徒们？是属神的，还是属肉身的？\par

\Person{厄兰尼斯特}——属神的，而且他们是属神之人的领袖和导师。\par

\Person{正统派}——现在请听圣保禄说的话：\BibleQuote{然而，那位从我母腹中分别出来、又藉他的恩宠召叫我的天主，既然乐意将他的圣子启示在我里面，叫我将他传扬在外邦人中，我就立刻没有与血肉之人商议，也没有上到那些在我以前作宗徒的人那里去。}他这样说，是因为他责备那些宗徒吗？\par

\Person{厄兰尼斯特}——当然不是。\par

\Person{正统派}——难道他不是按照他们可见的本性来称呼他们，并将那出于人的召叫与那出于天上的召叫相对比吗？\par

\Person{厄兰尼斯特}——对。\par

\Person{正统派}——那么也请听圣咏作者达味的话：\BibleQuote{凡有血气的，都要来就你。}也请听依撒意亚先知预言说：\BibleQuote{凡有血气的，都要看见我们天主的救恩。}\par

\Person{厄兰尼斯特}——已充分证明，圣经用\Quote{血肉}来指称人性，并无丝毫责备之意。\par

\Person{正统派}——我还要给你进一步的证明。\par

\Person{厄兰尼斯特}——还有什么？\par

\Person{正统派}——有时，当圣经责备人时，它只用了\Quote{灵魂}这个名称。\par

\Person{厄兰尼斯特}——你在圣经的哪里能找到这一点？\par

\Person{正统派}——请听上主天主藉厄则克耳先知说：\BibleQuote{犯罪的那灵魂，它必要死亡。}又藉伟大的梅瑟说：\BibleQuote{若有一个灵魂犯罪……}又说：\BibleQuote{将来，凡不肯听从那先知的灵魂，必被除灭。}还可以找到许多类似的经文。\par

\Person{厄兰尼斯特}——这已清楚证明。\par

\Person{正统派}——那么，在存在着某种自然结合，以及受造物之间、因着服务与时间而相连的存在物之间的组合时，圣经并不惯于只用那较高贵的本性来称呼这个存有；它既用那较低贱的，也用那较高贵的来无区别地称呼它。既然如此，当我们承认主基督是天主之后，称祂为人，而许多因素又迫使我们如此行时，你们怎能因此指责我们呢？\par

\Person{厄兰尼斯特}——有什么因素迫使我们称救主基督为\Quote{人}呢？\par

\Person{正统派}——异端者们那些彼此不同、互为矛盾的见解。\par

\Person{厄兰尼斯特}——什么见解？与什么相矛盾？\par

\Person{正统派}——亚略的见解与撒伯里乌的见解相反。前者分割实体；后者混淆位格。亚略引入三个实体，撒伯里乌则用一个位格代替三个。请告诉我，我们应该如何医治这两种病症？是用同一种药物来治两种病，还是各用其对症的药？\par

\Person{厄兰尼斯特}——各用其对症的药。\par

\Person{正统派}——那么我们要努力说服亚略，要他承认天主圣三的实体，并且我们会从圣经中为这一立场提出证据。\par

\Person{厄兰尼斯特}——是的，应该这样做。\par

\Person{正统派}——但在与撒伯里乌辩论时，我们将采取相反的做法。关于实体，我们不会提出任何论点，因为连他也只承认一个实体。\par

\Person{厄兰尼斯特}——显然。\par

\Person{正统派}——但我们会尽力医治他学说中有缺陷的部分。\par

\Person{厄兰尼斯特}——我们说过，他跛脚的地方在于位格方面。\par

\Person{正统派}——既然他断言圣三是一个位格，我们将向他指出，圣经宣告了三个位格。\par

\Person{厄兰尼斯特}——这是应该采取的步骤。但我们离题了。\par

\Person{正统派}——一点也没有。我们正在收集它的证据，你马上就会明白。但请告诉我，你是否认为所有以基督命名的异端，都承认基督的天主性和人性？\par

\Person{厄兰尼斯特}——绝不。\par

\Person{正统派}——难道不是有些只承认天主性，有些只承认人性吗？\par

\Person{厄兰尼斯特}——是的。\par

\Person{正统派}——而有些只承认人性的部分？\par

\Person{厄兰尼斯特}——我想是的。但最好让我们列出这些不同意见的持有者的名字，以便使讨论更加清楚。\par

\Person{正统派}——我会告诉你名字。西门、梅南德、马西昂、瓦伦提努、巴西里德、巴德桑、塞尔多和摩尼，公开否认基督的人性。另一方面，阿尔特蒙、提奥多图、撒伯里乌、萨莫萨特的保禄、玛尔塞鲁和佛提诺，则陷入了完全相反的亵渎；因为他们宣扬基督只是人，并否认那在万世以前就存在的天主性。亚略和欧诺米把独生子的天主性当作受造的天主性，并主张祂只取了身体。阿波利纳里承认所取的身体是活的身体，但在他的著作中，却剥夺了理性灵魂的尊荣与救恩。这就是这些败坏见解的相反之处。但请你本着对真理应有的爱，告诉我们：我们必须与这些人进行讨论，还是任由他们惨叫着、头朝下坠向灭亡？\par

\Person{厄兰尼斯特}——忽视病人是不人道的。\par

\Person{正统派}——很好；那么我们必须怜悯他们，并尽力医治他们。\par

\Person{厄兰尼斯特}——当然。\par

\Person{正统派}——如果你已科学地学会了如何医治身体，而你周围站着许多人向你求医，并展现出各种病痛，例如眼疾、耳伤、牙痛、关节挛缩、瘫痪、胆汁或痰涎，你会怎么做？请告诉我：你会对所有人都用同一种疗法，还是对每种病用其对症的疗法？\par

\Person{厄兰尼斯特}——我当然会给每个人对症的药物。\par

\Person{正统派}——那么，通过对热的施以冷疗，对冷的施以热疗，对紧张的使之松弛，对松弛的使之紧张，对潮湿的使之干燥，对干燥的使之湿润，你就能驱除疾病，恢复它们所驱走的健康。\par

\Person{厄兰尼斯特}——这是医学所规定的疗法，因为人们说，相反者是相反者的解药。\par

\Person{正统派}——如果你是园丁，你会对所有植物用同一种护理吗？还是对桑树、无花果树各自给予它们的护理？对梨树、苹果树和葡萄树，各给予适合它们的？总之，对每种植物给予其特有的栽培？\par

\Emph{埃兰}——显然，每株植物都需要各自的处理方法。\par

\Emph{正统者}——又假如你从事造船，见桅杆需要修理，你会用修理舵柄的方式来修理它吗？还是会给它桅杆应有的处理？\par

\Emph{埃兰}——这些事毫无疑问：每样事物都需要自身的处理，无论是植物、肢体、器具或索具。\par

\Emph{正统者}——那么，对身体和无生命之物各施以适宜的处理，而对灵魂却不遵守这一处理规则，岂不荒谬？\par

\Emph{埃兰}——极不正当；不，与其说是不义，不如说是愚蠢。采用其他方法的人，在治疗艺术上完全无技。\par

\Emph{正统者}——那么，在反对每一种异端时，我们是否要用适当的疗法？\par

\Emph{埃兰}——当然。\par

\Emph{正统者}——而适当的疗法，就是补其不足、去其多余吗？\par

\Emph{埃兰}——是的。\par

\Emph{正统者}——那么，在努力治疗\Person{福提诺}、\Person{玛尔塞禄}及其追随者时，为执行治疗规则，我们应添加什么？\par

\Emph{埃兰}——承认基督的\OriginalTerm{天主性}{Godhead}，因为他们所缺的正是这个。\par

\Emph{正统者}——但关于\OriginalTerm{人性}{manhood}，我们不会对他们说什么，因为他们承认主基督是人。\par

\Emph{埃兰}——你说得对。\par

\Emph{正统者}——而在与\Person{亚略}和\Person{欧诺米}辩论关于\OriginalTerm{独生子}{only begotten}的\OriginalTerm{降生成人}{incarnation}时，我们应说服他们在自己的宣认中添加什么？\par

\Emph{埃兰}——\OriginalTerm{取了灵魂}{assumption of the soul}；因为他们说\OriginalTerm{天主圣言}{divine Word}只取了肉体。\par

\Emph{正统者}——那么，\Person{阿波利纳里}为使他的教导在降生成人上准确，所缺的是什么？\par

\Emph{埃兰}——不将心智与灵魂分开，而是承认在肉体之外，还取了一个\OriginalTerm{理性灵魂}{reasonable soul}。\par

\Emph{正统者}——那么，我们要就这一点与他辩论吗？\par

\Emph{埃兰}——当然。\par

\Emph{正统者}——但在这一标题下，我们断言\Person{马西昂}、\Person{瓦伦提诺}、\Person{摩尼}及其追随者所承认的是什么，所全然否认的又是什么？\par

\Emph{埃兰}——他们承认他们相信基督的\OriginalTerm{天主性}{Godhead}，但不接受祂\OriginalTerm{人性}{manhood}的教义。\par

\Emph{正统者}——因此我们将尽力说服他们也接受人性的教义，而不把\OriginalTerm{天主降生成人}{divine incarnation}称为单纯的\OriginalTerm{幻象}{appearance}。\par

\Emph{埃兰}——这样做很好。\par

\Emph{正统者}——因此我们要告诉他们，称基督不仅是天主，也是人，是正确的。\par

\Emph{埃兰}——当然。\par

\Emph{正统者}——而当我们自己逃避称基督为\Quote{人}时，又怎么可能说服别人这样称呼呢？他们不会听从我们的劝说，反而会责备我们同意他们。\par

\Emph{埃兰}——而我们既承认\OriginalTerm{天主圣言}{divine Word}取了\OriginalTerm{肉体}{flesh}和\OriginalTerm{理性灵魂}{reasonable soul}，又怎能同意他们呢？\par

\Emph{正统者}——如果我们承认事实，为何要回避这个词呢？\par

\Emph{埃兰}——应从祂更高贵的品质来称呼基督。\par

\Emph{正统者}——那么，就遵守这条规则吧。不要称祂为被钉十字架，也不要称祂为从死者中复活，诸如此类。\par

\Emph{埃兰}——但这些乃是\OriginalTerm{救恩}{salvation}受苦的名称。否认受苦就意味着否认\OriginalTerm{救恩}{salvation}。\par

\Emph{正统者}——而\Quote{人}这个名称是一个\OriginalTerm{性体}{nature}的名称。不发音这个词就是否认这个\OriginalTerm{性体}{nature}：否认\OriginalTerm{性体}{nature}就是否认受苦，而否认受苦就取消了\OriginalTerm{救恩}{salvation}。\par

\Emph{埃兰}——我认为承认\OriginalTerm{所取的性体}{assumed nature}是有益的；但称世界的\OriginalTerm{救主}{Saviour}为人，却是贬低了主的荣耀。\par

\Emph{正统者}——那么，你以为自己比\Person{伯多禄}和\Person{保禄}更智慧吗？甚至比\OriginalTerm{救主}{Saviour}本身还智慧？因为主对犹太人说过：\Quote{你们为什么想要杀我？我是一个把从父那里听到的真理告诉了你们的人。}而且祂常自称\OriginalTerm{人子}{Son of Man}。\par

而可敬的\Person{伯多禄}在向犹太民众讲道时说：\Quote{以色列人哪，请听这些话。\Person{纳匝肋耶稣}，一位由天主在你们中间证明的人。}而蒙福的\Person{保禄}在向\Place{亚略巴古}的首领传达\OriginalTerm{救恩}{salvation}的信息时，在许多话中也说了这样一段——\par

\Quote{天主在过去无知的时代，虽容忍了，但现在却命令各处的人都要悔改：因为祂已定了一个日子，要按正义藉着祂所指定的那个人来审判世界，并且赐给众人凭据，因叫祂从死者中复活了。}那么，凡回避使用主和祂的宗徒所指定并宣讲的名字的人，就自以为比这些伟大的导师还要智慧，是的，甚至比智慧本身的泉源还要智慧。\par

\Emph{埃兰}——他们给予这些教导是针对不信者的。如今世上大部分地方已经宣认了\OriginalTerm{信德}{faith}。\par

\Emph{正统者}——但我们中间仍有\OriginalTerm{犹太人}{Jews}、\OriginalTerm{外教人}{pagans}以及数不胜数的异端体系，对每一种我们都必须给予合适恰当的教导。但假设我们全都同心，请告诉我，称基督兼为天主和人有什么害处呢？难道我们没有在祂身上看到完美的\OriginalTerm{天主性}{Godhead}，以及同样毫无欠缺的\OriginalTerm{人性}{manhood}吗？\par

\Emph{埃兰}——这一点我们已一再承认。\par

\Emph{正统者}——那么，为何要否认我们一再承认的事呢？\par

\Emph{埃兰}——我认为称基督为\Quote{人}是不必要的，尤其在信徒与信徒交谈时。\par

\Emph{正统者}——你认为\OriginalTerm{神性宗徒}{divine Apostle}是一位信徒吗？\par

\Emph{埃兰}——是的：是所有信徒的导师。\par

\Emph{正统者}——而你认为\Person{弟茂德}配得上这样的称呼吗？\par

\Emph{埃兰}——是的：既作为宗徒的门徒，也作为其他人的导师。\par

\Emph{Orth.}—很好：那么请听那位万师之师给他非常成全的门徒写信所说的话：\BibleQuote{只有一个天主，也只有一个中保，在天主与人之间，就是人基督耶稣，他交出了自己，为众人作赎价。}停止你无聊的闲谈，不要再对天主的圣名妄加论断。而且，在这段经文中，\Quote{中保}这一名称本身就指示了\OriginalTerm{天主性}{Godhead}和\OriginalTerm{人性}{manhood}。他之所以被称为中保，是因为他并非仅作为天主而存在；因为，如果他没有任何属于我们本性的东西，他又如何能成为我们与天主之间的中保呢？但由于作为天主，他与天主结合，拥有同一本体；而作为人，与我们结合，因为从我们中他取了奴仆的形象，所以他被恰当地称为中保，在天主性（我说）和人性本性的合一中，将不同的性质在自己内统一起来。\par

\Emph{Eran.}—但梅瑟不也被称为\OriginalTerm{中保}{mediator}吗？尽管他只是一个人。\par

\Emph{Orth.}—他是实体的\OriginalTerm{预像}{type}；但预像并不具有实体的全部性质。因此，尽管梅瑟按本性不是天主，但为成就预像，他被称为神。因为祂说：\BibleQuote{看，我使你作了法郎的天主。}然后紧接着就给他指派了一位先知，如同给天主一样，因为祂说：\BibleQuote{你的兄弟亚郎，要作你的先知。}但实体按本性是天主，按本性是人。\par

\Emph{Eran.}—但谁会称一个不具有\OriginalTerm{原型}{archetype}显著特征的\OriginalTerm{预像}{type}为预像呢？\par

\Emph{Orth.}—看来，你不把皇帝的画像称为皇帝像。\par

\Emph{Eran.}—是的，我称它们为皇帝像。\par

\Emph{Orth.}—然而它们并不具有其\OriginalTerm{原型}{archetype}的全部特征。首先，它们没有生命和理性；其次，它们没有内在器官，即心、腹、肝及相邻部位。此外，它们呈现感觉器官的外观，却不执行任何功能，因为既不听、不说、不看；不能书写，不能行走，也不能进行任何其他人的活动；然而它们仍被称为御像。在这个意义上，梅瑟是中保，基督也是中保；但前者是作为\OriginalTerm{预像}{type}和\OriginalTerm{画像}{image}，后者是作为\OriginalTerm{实体}{reality}。但为使这一点从另一权威更清楚地向你显明，请回忆\WorkTitle{致希伯来人书}中论及默基瑟德的话。\par

\Emph{Eran.}—什么话？\par

\Emph{Orth.}—就是那位神圣的宗徒在比较\OriginalTerm{肋未司祭职}{Levitical priesthood}与基督司祭职时，在其他方面将默基瑟德比作主基督，并说主是按照\OriginalTerm{默基瑟德品位}{order of Melchisedec}作司祭的那些话。\par

\Emph{Eran.}—我想神圣宗徒的话如下：\BibleQuote{这位默基瑟德是撒冷王，是至高者天主的司祭；他迎接了亚巴郎从击杀诸王回来，祝福了他；亚巴郎也将一切所得的十分之一给了他。首先，他的名字翻译出来是\InnerQuote{正义之王}，其次他也是撒冷王，即\InnerQuote{和平之王}；他没有父亲，没有母亲，没有族谱，既无生之始，也无命之终；他却成了天主子相似，永远做司祭。}我猜想你说的是这段经文。\par

\Emph{Orth.}—是的，我说的正是这段。我得赞美你没有删节，而是引用了全部。请告诉我，这些特点中，每一个都自然地、真实地符合默基瑟德吗？\par

\Emph{Eran.}—神圣的宗徒既然断言了相符，谁还敢否认？\par

\Emph{Orth.}—那么你说所有这些都自然符合默基瑟德？\par

\Emph{Eran.}—是的。\par

\Emph{Orth.}—你说他是人，还是取了别的本性？\par

\Emph{Eran.}—是人。\par

\Emph{Orth.}—\OriginalTerm{受生}{begotten}还是\OriginalTerm{非受生}{unbegotten}？\par

\Emph{Eran.}—你问的问题非常荒谬。\par

\Emph{Orth.}—错在你公开反对真理。回答吧。\par

\Emph{Eran.}—只有一位\OriginalTerm{非受生}{unbegotten}者，就是天主父。\par

\Emph{Orth.}—那么，我们断言默基瑟德是\OriginalTerm{受生}{begotten}的？\par

\Emph{Eran.}—是的。\par

\Emph{Orth.}—但关于他的经文教导的正好相反。还记得你刚才引用的经文吗？\BibleQuote{他没有父亲，没有母亲，没有族谱，既无生之始，也无命之终。}那么，\Quote{无父无母}如何适用于他？说他既无存在之始也无终结，既然这一切都超越了人性？\par

\Emph{Eran.}—这些事确实越过了人性的界限。\par

\Emph{Orth.}—那么我们要说宗徒说谎了吗？\par

\Emph{Eran.}—绝对不要。\par

\Emph{Orth.}—那么，怎么既能证明宗徒的真实，又将\OriginalTerm{超自然}{supernatural}的事运用到默基瑟德身上呢？\par

\Emph{Eran.}—这段经文非常难解，需要很多解释。\par

\Emph{Orth.}—只要有人愿意留心思考，就不难领悟这些话的意思。说了\BibleQuote{他没有父亲，没有母亲，没有族谱，既无生之始，也无命之终}，神圣的宗徒接着又说\BibleQuote{他却成了天主子相似，永远做司祭。}这里他清楚地教导我们，在关于人性的事上，主基督是默基瑟德的\OriginalTerm{原型}{archetype}。他说默基瑟德\Quote{成了天主子相似}。现在，我们这样来考察这一点——你说主按肉体有父亲吗？\par

\Emph{Eran.}—当然没有。\par

\Emph{Orth.}—为什么？\par

\Emph{Eran.}—他是由圣\OriginalTerm{童贞玛利亚}{holy Virgin}独自所生。\par

\Emph{Orth.}—所以他被恰当地称为\OriginalTerm{无父}{without father}？\BibleQuote{他没有父亲。}\par

\Emph{Eran.}—对。\par

\Emph{Orth.}—你说按\OriginalTerm{天主性}{divine Nature}他有母亲吗？\par

\Emph{Eran.}—当然没有。\par

\Emph{Orth.}—因为他是\OriginalTerm{万世之前}{before the ages}由圣父独自\OriginalTerm{受生}{begotten}的？\par

\Emph{Eran.}—同意。\par

\Emph{Orth.}—然而，由于他由圣父所生的方式是\OriginalTerm{不可言喻}{ineffable}的，所以被称为\OriginalTerm{无族谱}{without descent}。先知说：\BibleQuote{谁能述说他的世代？}\par

\Emph{Eran.}—你说得对。\par

\Emph{Orth.}—这样，他既无生之始也无命之终就适合他了；因为他\OriginalTerm{无始}{without beginning}、\OriginalTerm{不可朽坏}{indestructible}，总而言之，是\OriginalTerm{永恒}{eternal}的，并与圣父\OriginalTerm{与圣父同等永恒}{coeternal with the Father}。\par

\Emph{Eran.}—这也是我的看法。但我们现在必须考虑这如何适用于可敬的默基瑟德。\par

\Emph{正统。}——作为图像和预像。正如我们刚刚指出的，图像并不具备原型的全部属性。因此，这些特性在本性上和事实上都属于救主；但人类起源的记述却将它们归于默基瑟德。因为在告诉我们圣祖亚巴郎的父亲、依撒格的父母、以及雅各伯和他的儿子们，并展示我们始祖的谱系之后，关于默基瑟德，经文既未记录他的父亲或母亲，也没有说他出自诺厄的任何儿子，为使他成为那真正无父、无母者的预像。这正是神圣宗徒要我们理解的，因为就在这段经文中，他接着说：\BibleQuote{但默基瑟德，虽不属于他们的世系，却从亚巴郎收取了什一之物，并祝福了那持有恩许的。}\par

\Emph{异端。}——那么，既然圣经未提及他的父母，他能被称为无父无母的吗？\par

\Emph{正统。}——如果他真是无父无母的，那他就不是图像，而是原型本身了。但由于这些特性并非出于本性，而是出于圣经的安排，他便展现出原型的预像。\par

\Emph{异端。}——预像必须具有原型的特征。\par

\Emph{正统。}——人被称为天主的肖像吗？\par

\Emph{异端。}——人不是天主的肖像，而是按照天主的肖像受造的。\par

\Emph{正统。}——那么请听宗徒的话。他说：\BibleQuote{男人当然不应蒙头，因为他是天主的肖像和光荣。}\par

\Emph{异端。}——姑且承认他是天主的肖像。\par

\Emph{正统。}——那么按你的说法，他必须清楚保存原型的特征，成为非受造的、非组合的、无限的。同样，他应当能从无中创造，用他的话语不劳而造，此外，还应无病、无忧、无怒、无罪，不朽、不坏，并具备原型的一切属性。\par

\Emph{异端。}——人并非在每一方面都是天主的肖像。\par

\Emph{正统。}——尽管在那些你承认的方面，他确实是肖像，但你将发现他与原型之间相隔甚远。\par

\Emph{异端。}——同意。\par

\Emph{正统。}——现在也请考虑这一点。神圣宗徒称圣子为圣父的肖像，因为他说：\BibleQuote{谁是不可见的天主的肖像？}\par

\Emph{异端。}——那么，圣子不具备圣父的一切属性吗？\par

\Emph{正统。}——祂不是圣父。祂不是无原因的。祂不是无生的。\par

\Emph{异端。}——如果祂是，祂就不是子了。\par

\Emph{正统。}——那么，我所说的不成立吗？图像不具备原型的一切属性？\par

\Emph{异端。}——对。\par

\Emph{正统。}——同样，神圣宗徒说默基瑟德相似天主子。\par

\Emph{异端。}——假设我们如你所说，承认他是无父、无母、无族谱的。但我们该如何理解他没有生之始、死之终呢？\par

\Emph{正统。}——神圣梅瑟在撰写古代谱系时，告诉我们亚当活了若干岁生了舍特，又活了若干岁后结束了生命。同样，他也写了舍特、哈诺客等其他人，但对于默基瑟德，他既未提及存在的开始，也未提及生命的终结。因此，就叙述而言，他没有生之始、死之终；但在真理和事实上，天主的独生子从未开始存在，也永无终结。\par

\Emph{异端。}——同意。\par

\Emph{正统。}——那么，就属于天主且真正神圣的方面而言，默基瑟德是主基督的预像；但就司祭职而言，它更属于人类而非天主，主基督是按照默基瑟德的品位作司祭的。因为默基瑟德是民众的大司祭，而主基督为全人类奉献了正义而神圣的救恩祭献。\par

\Emph{异端。}——我们在这件事上已花费许多言辞。\par

\Emph{正统。}——但还需要更多，如你所知，因为你说这是个难题。\par

\Emph{异端。}——让我们回到先前的问题。\par

\Emph{正统。}——什么问题？\par

\Emph{异端。}——当我主张基督不可被称为人，而只应被称为天主时，你本人除了许多其他见证外，还引用了宗徒在致弟茂德书中所用的名言——\BibleQuote{只有一位天主，在天主与人之间的中保也只有一位，就是那人基督耶稣，祂奉献了自己，为众人作赎价，这事在适当的时候得到了证明。}\par

\Emph{正统。}——我记得我们是从哪一点岔入这段离题的。当时我说中保的名称显示出救主的两种本性，而你说梅瑟被称为中保，尽管他只是一个人，并非天主和人。因此我不得不追述这些要点，以表明预像并不具备原型的一切属性。那么请告诉我，你是否承认救主也当被称为人？\par

\Emph{异端。}——我称祂为天主，因为祂是天主子。\par

\Emph{正统。}——如果你称祂为天主，因为你已知道祂是天主子，那么也称祂为人吧，因为祂常自称\Quote{人子}。\par

\Emph{异端。}——人的名号不适用于祂，如同天主的名号一样。\par

\Emph{正统。}——是因为并非真实属于祂，还是另有原因？\par

\Emph{异端。}——天主是祂本性的名号；人是降生成人的称谓。\par

\Emph{正统。}——但我们该看降生成人是真实的，还是某种想象和虚假的？\par

\Emph{异端。}——真实的。\par

\Emph{正统。}——那么，如果降生成人的恩宠是真实的，而我们所谓的降生成人，就是圣言成了人，那么人的名号就是真实的；因为祂取了人的本性后，被称为人。\par

\Emph{异端。}——在受难前，祂被称为人，但之后就不再如此称呼了。\par

\Person{Orth.}—— 但这是在受难与复活之后，神圣的宗徒在写给弟茂德的\WorkTitle{弟茂德书}中谈到救主基督是人；并在受难与复活之后写给格林多人的\WorkTitle{格林多前书}中高呼：\BibleQuote{因为死既是由于一人，死者的复活也是由于一人。}为了使其含义清楚，他又补充说：\BibleQuote{因为在亚当内众人都死了，照样，在基督内众人都要复活。}在受难与复活之后，神圣的伯多禄在向犹太人讲话时称祂为人。在祂被接升天之后，胜利的斯德望在石雨中向犹太人说：\BibleQuote{看，我见天开了，并见人子站在天主的右边。}难道我们以为自己比真理的光荣传报者更聪明吗？\par

\Person{Eran.}—— 我并不以为自己比圣师们更聪明，但我找不到使用这个称号的理由。\par

\Person{Orth.}—— 那么，如果否认主降生成人的那些人——我指的是马西昂派和摩尼教徒以及所有其他这样不正统的人——除非你提出这些和类似的证据，目的是表明主基督不但是天主，而且也是人，你怎能说服他们接受真理的教导呢？\par

\Person{Eran.}—— 也许有必要提出它们。\par

\Person{Orth.}—— 那么为什么不教导信徒教义的真实性呢？你忘记了宗徒的训诫吗？它命令我们要\BibleQuote{常准备好答复}。现在让我们从这一点来看问题：最好的将军是独自与敌人交战、用箭和标枪攻击、并尽力突破他们的阵线呢，还是他也武装他的士兵、部署他们、并鼓起他们的勇气去战斗？\par

\Person{Eran.}—— 他更应该做后一件事。\par

\Person{Orth.}—— 是的；因为将军的本分不是暴露自己的生命、亲自站在队列中、而让他的士兵酣睡，而是让他们在岗位上保持清醒准备战斗。\par

\Person{Eran.}—— 对。\par

\Person{Orth.}—— 神圣的保禄正是这样做的，因为他在写信给那些宣了信德的人时说：\BibleQuote{穿上天主的全副武装，为使你们能抵抗魔鬼的阴谋。}又说：\BibleQuote{所以要站稳，用真理作带，束紧你们的腰，}等等。请记住我们先前说过的话：医生补充自然所缺乏的。如果他发现冷过剩，他就提供热，如此等等；主正是这样做的。\par

\Person{Eran.}—— 你要在哪里显示主这样做过？\par

\Person{Orth.}—— 在圣福音中。\par

\Person{Eran.}—— 那么请展示给我看，履行你的诺言。\par

\Person{Orth.}—— 犹太人把我们的救主基督看作什么？\par

\Person{Eran.}—— 一个人。\par

\Person{Orth.}—— 他们完全不知道祂也是天主。\par

\Person{Eran.}—— 是的。\par

\Person{Orth.}—— 那么无知的人难道不必要学习吗？\par

\Person{Eran.}—— 同意。\par

\Person{Orth.}—— 那么请听祂对他们说的话：\BibleQuote{我从父那里显示了许多善事给你们；你们为了哪一件善事要拿石头打我？}当他们回答说：\BibleQuote{我们不是为了善事要拿石头打你，而是为了亵渎的话，并且因为你是一个人，却把自己当作天主。}祂又补充说：\BibleQuote{你们的法律上不是写着：\InnerQuote{我说过，你们是神}么？如果那些承受天主的话的人，经书是不能废弃的，尚且被称为神；那么，父所祝圣并派遣到世界上来的，你们却说：你说亵渎的话，因为我说了\InnerQuote{我是天主子}？我若不作我父的事，你们就不必信我……使你们知道父在我内，我在父内。}\par

\Person{Eran.}—— 在你刚才读的经文中，你已显示主向犹太人表明祂是天主而不是人。\par

\Person{Orth.}—— 是的，因为他们不需要学习他们所知道的事；他们知道祂是一个人，但他们不知道祂从起初就是天主。祂在处理法利塞人时也采取了同样的方式；因为当祂看到他们只把祂当作一个人来对待时，祂问他们：\BibleQuote{你们对基督有何看法？祂是谁的儿子？}当他们说：\BibleQuote{达味的。}祂继续说：\BibleQuote{那么，达味怎么还称祂为主，说：\InnerQuote{上主对我主说：你坐在我右边。}}然后祂继续论证：\BibleQuote{那么，如果达味称祂为主，祂怎能又是达味的儿子呢？}\par

\Person{Eran.}—— 你搬起了石头砸自己的脚，因为主明确地教导法利塞人不要称祂为\Quote{达味之子}，而要称祂为\Quote{达味之主}。因此，祂显然表明愿意被称为天主而不是人。\par

\Person{Orth.}—— 恐怕你没有留意神圣的教导。祂并没有拒绝\Quote{达味之子}这个称号，而是补充说祂也应被认为达味之主。这祂在以下的话中清楚地表明：\BibleQuote{如果达味称祂为主，祂怎能又是他的儿子呢？}祂没有说\Quote{如果祂是主，就不是儿子}，而是说\Quote{祂怎能又是他的儿子呢？}这等于说：从一方面看，祂是主；从另一方面看，祂是儿子。这些经文既清楚地显示了天主性，也显示了人性。\par

\Person{Eran.}—— 无须辩论。主清楚地教导说，祂不希望被称为达味之子。\par

\Person{Orth.}—— 那么祂就应该告诉那几位盲人、客纳罕妇人和群众不要称祂为达味之子；然而，盲人大声呼喊：\BibleQuote{达味之子，可怜我们吧！}客纳罕妇人说：\BibleQuote{主，达味之子，可怜我罢！我的女儿被鬼折磨得很苦。}群众：\BibleQuote{贺三纳归于达味之子！奉主名而来的，应受赞颂！}祂不仅不生气，反而称赞了他们的信德；因为祂使盲人摆脱了他们漫长疲惫的黑夜，赐予他们看见的能力；祂治愈了客纳罕妇人那疯狂错乱的女儿，赶出了邪恶的魔鬼；当司祭长和法利塞人对那些呼喊\Quote{贺三纳归于达味之子}的人感到恼怒时，祂不仅没有阻止他们呼喊，反而认可了他们的欢呼，因为祂说：\BibleQuote{我告诉你们：这些人若不作声，石头就要喊叫起来。}\par

\Emph{Eran.}——祂在复活前忍受这种称呼的方式，是为了迁就那些尚未正确相信之人的软弱。但复活之后，这些称呼便无必要了。\par

\Emph{Orth.}——我们应将蒙福的保禄置于何处？是完美者中，还是不完美者中？\par

\Emph{Eran.}——拿严肃的事开玩笑是不对的。\par

\Emph{Orth.}——轻视神圣启示的诵读是不对的。\par

\Emph{Eran.}——又有谁会如此不幸，竟至于藐视自己的救恩呢？\par

\Emph{Orth.}——回答我的问题，然后你便会知道自己的无知。\par

\Emph{Eran.}——什么问题？\par

\Emph{Orth.}——我们应将这位神圣的宗徒置于何处？\par

\Emph{Eran.}——显然属于最完美者，且是完美的教师之一。\par

\Emph{Orth.}——他何时开始教导？\par

\Emph{Eran.}——在救主升天、圣神降临以及得胜的斯德望被石头砸死之后。\par

\Emph{Orth.}——保禄在生命尽头，在写给门徒弟茂德的最后一封信中，仿佛将父辈的产业遗赠给他时，补充说：\Quote{你要记着：达味的后裔耶稣基督从死者中复活了，符合我所传的福音。}然后他提到了自己为福音所受的苦难，以此证明福音的真实性，说：\Quote{我为这福音受苦，甚至像作恶的人一样被捆锁。}\par

我可以轻易提出许多类似的见证，但我认为并无必要这样做。\par

\Emph{Eran.}——你曾应许证明主为需要者提供了欠缺的教导，并且你已经表明祂向法利塞人及其他犹太人谈论了祂自己的天主性。但关于血肉的教导，你尚未表明。\par

\Emph{Orth.}——谈论那在他们眼前的血肉本是多余，因为祂显然吃喝、劳苦、睡眠。此外，撇开受难前的诸多事件不谈，祂复活后向不信的门徒证明的并非祂的天主性，而是祂的人性；因为祂说：\Quote{你们看我的手和脚，是我自己。你们摸摸我，看看，因为鬼神没有骨肉，如同你们看见我有的。}\par

如今我履行了诺言，因为我们已证明：对不认识天主性的人，祂教导了天主性；对否认肉身复活的人，祂教导了肉身的复活。因此，停止争辩吧，承认救主的两个性体。\par

\Emph{Eran.}——在结合之前有两个，但结合之后便成了一个性体。\par

\Emph{Orth.}——你说结合是在何时发生的？\par

\Emph{Eran.}——我说是在成孕的那一刻。\par

\Emph{Orth.}——你是否否认天主圣言在成孕之前已存在？\par

\Emph{Eran.}——我说祂在万世之前已存在。\par

\Emph{Orth.}——那么血肉与祂共存吗？\par

\Emph{Eran.}——绝非如此。\par

\Emph{Orth.}——而是在天使问候之后，由圣神形成的？\par

\Emph{Eran.}——我是这样说的。\par

\Emph{Orth.}——因此，在结合之前，并非有两个性体，而只有一个。因为如果天主性预先存在，但人性并未共存，是在天使问候之后才形成，并且结合与形成同时发生，那么在结合之前只有一个性体，即那永远存在、在万世之前就已存在的性体。现在让我们再思考这一点：你认为成血肉或成为人，是否与结合有所不同？\par

\Emph{Eran.}——不。\par

\Emph{Orth.}——因为祂取了血肉，便成了血肉。\par

\Emph{Eran.}——显然。\par

\Emph{Orth.}——而结合与取血肉是同时发生的。\par

\Emph{Eran.}——我是这样说的。\par

\Emph{Orth.}——所以在成为人之前，只有一个性体。因为如果结合与成为人是同一的，并且祂通过取人性而成为人，而天主的形像取了奴仆的形像，那么在结合之前，天主性是一个。\par

\Emph{Eran.}——那么结合与成为人如何是同一的呢？\par

\Emph{Orth.}——刚才你承认这些术语之间没有区别。\par

\Emph{Eran.}——你用你的论证误导了我。\par

\Emph{Orth.}——那么，若你愿意，让我们再重述一遍同样的内容。\par

\Emph{Eran.}——我们最好这样做。\par

\Emph{Orth.}——按照此事的性质，降生成人与结合之间有区别吗？\par

\Emph{Eran.}——当然，有极大的区别。\par

\Emph{Orth.}——请充分说明这种区别的性质。\par

\Emph{Eran.}——甚至这些语词的含义就表明了区别，因为\Quote{降生成人}一词表明取了血肉，而\Quote{结合}一词则表明不同事物的组合。\par

\Emph{Orth.}——你是否认为降生成人先于结合？\par

\Emph{Eran.}——绝非如此。\par

\Emph{Orth.}——你说结合发生在成孕之中？\par

\Emph{Eran.}——是的。\par

\Emph{Orth.}——因此，如果取血肉与结合之间连一刹那的间隔都没有，且被取的本性并未先于取和结合而存在，那么降生成人与结合指的是同一件事。因此，在结合与降生成人之前，只有一个性体；而在降生成人之后，我们正确地谈论两个性体，即那取的与那被取的。\par

\Emph{Eran.}——我说基督出自两个性体，但我否认两个性体。\par

\Emph{Orth.}——那么请向我们说明，你如何理解\Quote{出自两个性体}这一表达？像镀金的银子？像电子（金银合金）的组合？像铅锡合金的焊料？\par

\Emph{Eran.}——我否认结合像这些中的任何一种；它是不可言喻的，超越所有理解。\par

\Emph{Orth.}——我也承认结合的方式无法理解。但我至少已从圣经得到教导：每个性体在结合后仍然完好无损。\par

\Emph{Eran.}——圣经哪里教导了这事？\par

\Emph{Orth.}——整部圣经充满了这一教导。\par

\Emph{Eran.}——请证明你所说的。\par

\Emph{Orth.}——难道你不承认每个性体的属性吗？\par

\Emph{Eran.}——不，也就是说，在结合之后不承认。\par

\Emph{Orth.}——那么让我们从圣经来学习这一点。\par

\Emph{Eran.}——我准备好服从圣经。\par

\Emph{Orth.}——那么，当你听到神圣的若望高呼\BibleQuote{在起初已有圣言，圣言与天主同在，圣言就是天主}和\BibleQuote{万物是藉着祂而造的}以及其余类似段落时，你肯定那肉体，还是那在万世之前由圣父所生、在起初就与天主同在、本性是天主、并创造了万物的天主圣言？\par

\Emph{Eran.}——我说这些属于天主圣言。但我并不将祂与成为一体的肉体分开。\par

\Emph{Orth.}——我们也不将肉体与天主圣言分开，也不将结合视为混淆。\par

\Emph{Eran.}——我承认结合后只有一个本性。\par

\Emph{Orth.}——福音作者们是在结合之前，还是在结合之后很长时间才写下福音？\par

\Emph{Eran.}——显然是在结合、诞生、奇迹、苦难、复活、升天以及圣神降临之后。\par

\Emph{Orth.}——那么请听若望说：\BibleQuote{在起初已有圣言，圣言与天主同在，圣言就是天主。祂在起初就与天主同在。万物是藉着祂造的，凡受造的，没有一样不是藉着祂造的}等等。也请听玛窦说：\BibleQuote{耶稣基督的族谱，达味之子，亚巴郎之子}等等。路加也将祂的家谱追溯到亚巴郎和达味。现在请将前者和后者的引文归于一个本性。你会发现这是不可能的，因为在起初存在与由亚巴郎而出，万物的创造与出自受造的祖先，是不一致的。\par

\Emph{Eran.}——你这样论证，是把独生子分成了两个位格。\par

\Emph{Orth.}——我认识并朝拜天主的一位圣子，主耶稣基督；但我已学会区分祂的神性与祂的人性。然而，你声称结合后只有一个本性，那么请使这与福音作者们的引言一致。\par

\Emph{Eran.}——你似乎认为这个命题很难，甚至是不可能的。我请求，让它简短而容易——只解答我们的问题。\par

\Emph{Orth.}——两种属性都属于主基督——在起初存在，以及按肉体从亚巴郎和达味所生。\par

\Emph{Eran.}——你曾定下法则：结合后不可论及一个本性。当心，在提及肉体时，你不要违反自己的法则。\par

\Emph{Orth.}——即使不提及肉体，也极易解释此问题，因为我把两者都应用于救主基督。\par

\Emph{Eran.}——我也断言这两种属性都属于主基督。\par

\Emph{Orth.}——是的，但你是在默观祂内的两个本性，并将各自的属性归于每一个。但如果基督是一个本性，怎么可能将互相矛盾的属性归于它？因为从亚巴郎和达味获得起源，尤其是达味之后许多代才出生，与在起初存在不一致。同样，出自受造之物与作为万物的创造者不一致；有人类的父亲与出自天主不一致。简言之，新与永恒不一致。\par

我们也这样来看这个问题。我们是否说天主圣言是宇宙的创造者？\par

\Emph{Eran.}——我们从圣经中学会了这样相信。\par

\Emph{Orth.}——那么，在创造天地之后多少天，亚当被造？\par

\Emph{Eran.}——第六天。\par

\Emph{Orth.}——从亚当到亚巴郎经过了多少代？\par

\Emph{Eran.}——我想是二十代。\par

\Emph{Orth.}——那么从亚巴郎到我们的救主基督，福音作者玛窦推算有多少代？\par

\Emph{Eran.}——四十二代。\par

\Emph{Orth.}——那么，如果主基督是一个本性，祂怎能同时是可见与不可见万物的创造者，又在如此多代之后由圣神在童贞女的胎中成形？祂又怎能同时是亚当的创造者又是亚当后裔的儿子？\par

\Emph{Eran.}——我已经说过，这两种属性都适合祂，因为祂是成为肉体的天主，因为我承认成为肉体的圣言只有一个本性。\par

\Emph{Orth.}——我的好先生，我们也不说天主圣言的两个本性成了肉体，因为我们知道天主圣言的本性是一个，但我们已学会，祂降生成人所取用的肉体是另一种本性，并且我认为你在此也同意我。现在告诉我：你说成肉身是如何发生的？\par

\Emph{Eran.}——我不知道如何发生，但我相信祂成了血肉。\par

\Emph{Orth.}——你不公平地以无知为借口，效法法利塞人的方式。因为他们眼见主询问的力量，又怀疑自己即将被驳倒，便回答：\BibleQuote{我们不知道。}但我十分公开地宣告：天主降生成人没有变化。因为如果祂藉任何变异或变化成了血肉，那么变化之后，祂的名字和作为中一切神圣的都不适合祂了。\par

\Emph{Eran.}——我们一再同意，天主圣言是不变的。\par

\Emph{Orth.}——祂藉取用肉体而成了血肉。\par

\Emph{Eran.}——是的。\par

\Emph{Orth.}——成了血肉的天主圣言之本性，异于那肉体，天主圣言的本性因取用这肉体而成了血肉并成为人。\par

\Emph{Eran.}——同意。\par

\Emph{Orth.}——那么祂是否变成了肉体？\par

\Emph{Eran.}——当然不是。\par

\Person{正统者}——如果祂成为血肉，不是由于变化，而是由于取了血肉，并且如你方才所说，前者和后者两种属性都适合祂这位成为血肉的天主，那么，两性并未混淆，而是保持完整无损。只要我们这样持守，我们也会看到福音作者们的和谐一致；因为一位宣告独生子——主基督——的天主性属性，另一位则陈述祂的人性属性。同样，主基督亲自教导我们，时而称自己为天主子，时而称自己为人子；时而尊崇祂的母亲，如同尊崇生祂者，时而又作为她的主责备她。有时祂不责备那些称祂为达味之子的人，有时又教导那些无知的人，祂不仅是达味之子，也是达味之主。祂称\Place{纳匝肋}和\Place{葛法翁}为祂的家乡，又喊道：\BibleQuote{在亚伯拉罕出现以前，我就是。}你会看到圣经中充满了类似的章节，它们所指的不是一个性体，而是两个。\par

\Person{厄兰涅斯特}——谁在基督内默观两个性体，就把独生子分成了两个儿子。\par

\Person{正统者}——是的；而谁说了保禄是由灵魂与身体组成的，就把一个保禄变成了两个保禄。\par

\Person{厄兰涅斯特}——这个类比不成立。\par

\Person{正统者}——我知道它不成立；因为在这里，结合是同龄、受造、同爲仆役的部分之间的自然结合；但在主基督的情形中，一切都是出于善意、爱人之情和恩宠。这里也一样，虽然结合是自然的，但各性体的特有属性仍保持完整无损。\par

\Person{厄兰涅斯特}——如果各性体的特有属性保持区别，灵魂怎能与身体一同渴求食物呢？\par

\Person{正统者}——灵魂并不渴求食物。它既是不朽的，怎能渴求呢？但身体从灵魂获得生命力，感到需要，并渴望接受所缺乏的。所以劳累后它渴望休息，醒来后渴望睡眠，其他欲望也是如此。所以一旦身体分解，由于不再有生命力，它甚至不再渴求所缺乏的，停止接受，便朽坏了。\par

\Person{厄兰涅斯特}——你看，口渴、饥饿和类似的欲望属于灵魂。\par

\Person{正统者}——如果这些属于灵魂，那么灵魂即使在脱离身体之后，也会遭受饥饿、口渴和类似的需要。\par

\Person{厄兰涅斯特}——那么你说什么是灵魂所特有的？\par

\Person{正统者}——理性、自主、不朽、无形可见。\par

\Person{厄兰涅斯特}——身体呢？\par

\Person{正统者}——复合、可见、可朽。\par

\Person{厄兰涅斯特}——我们说人是由这些组成的吗？\par

\Person{正统者}——是的。\par

\Person{厄兰涅斯特}——那么我们将人定义为有死有理性的生物。\par

\Person{正统者}——同意。\par

\Person{厄兰涅斯特}——我们从这两种属性出发给他命名。\par

\Person{正统者}——是的。\par

\Person{厄兰涅斯特}——既然如此，在这种情况下我们不作区分，而称同一个人既是理性的又是有死的；同样，在基督的情形中，我们也应当这样做，将天主性及人性都归于祂。\par

\Person{正统者}——这正是我们的论点，虽然你没有准确表达。你看，当我们讨论人的灵魂时，我们只提及适合其活动和本性的事吗？\par

\Person{厄兰涅斯特}——只此而已。\par

\Person{正统者}——当我们的讨论关乎身体时，我们不也只提及适合它的事吗？\par

\Person{厄兰涅斯特}——正是如此。\par

\Person{正统者}——但是，当我们的论述涉及整个存有时，我们便毫无困难地提出两组属性，因为身体和灵魂的属性都适用于人。\par

\Person{厄兰涅斯特}——毫无疑问。\par

\Person{正统者}——嗯，正应如此谈论基督：当讨论祂的性体时，各归其所有，确认某些属于天主性，某些属于人性。但当讨论位格时，我们必须将各性体的特有属性视为共有，并将两组属性都应用于救主，称同一存有既是天主又是人，既是天主子又是人子——既是达味之子又是达味之主，既是亚伯拉罕的后裔又是亚伯拉罕的创造者，等等。\par

\Person{厄兰涅斯特}——基督的位格是唯一的，天主性及人性都可以归于祂，你说得很对，我接受这一信德定义；但你真正的立场，即讨论性体时必须各归其属性，在我看来似乎消解了联合。正因如此，我拒绝接受这些论点及类似的论点。\par

\Person{正统者}——然而，当我们探讨灵魂与身体时，你认为这些术语的区分甚为精妙，并立即表示赞同。那么，为何你拒绝在主基督的天主性与人性上接受同样的规则呢？你甚至反对将基督的天主性与人性比作灵魂与身体吗？因此，你既然承认灵魂与身体有不混淆的结合，你敢说基督的天主性与人性经历了混合与混淆吗？\par

\Person{厄兰涅斯特}——我承认基督的天主性，是的，祂的肉身，远较灵魂与身体尊贵；但联合之后，我确言一个性体。\par

\Person{正统者}——但如今，既然主张灵魂与身体结合绝不遭受混淆，却否认宇宙之主的天主性有能力保持其本性不混淆，或使祂所摄取的人性保持在适当范围内，这难道不是不虔敬且骇人听闻的吗？我说，将分明之物混合，将分离之物搅在一起，难道不是不虔敬的吗？一个性体之观念提供了这种混淆的根据。\par

\Person{厄兰涅斯特}——我同样急切地避免\Quote{混淆}一词，但我害怕断言两个性体，以免陷入两个儿子的二元论。\par

正统者: — 我同样急于躲避两难困境的任何一个方面，既是不敬的混淆，也是不敬的区分；因为对我来说，将独一圣子分裂为二，以及否认两性之双重性，同样是不圣洁的想法。但是现在，请以真理之名告诉我。倘若亚略或欧诺米一派的人，在与你辩论时，试图贬低圣子，描述祂比圣父更小、更低，借助他们所有熟悉的论据和从圣经引用的经文，即\BibleQuote{父啊，如果可能，就让这杯离开我罢}和另一处\BibleQuote{现在我心神烦乱}以及其他类似段落，你将如何处理他的异议？你如何能表明圣子在这些表达中丝毫没有丧失尊严，也不是出于另一种实体，而是出于圣父的实体所生？\par

厄兰: — 我应该说，圣经按\OriginalTerm{神学}{theology}运用一些术语，按\OriginalTerm{经世}{œconomy}运用另一些，将属于\OriginalTerm{经世}{œconomy}的用于属于\OriginalTerm{神学}{theology}的，是错误的。\par

正统者: — 但你的对手会反驳说，即使在旧约中，圣经也以经世的方式说了许多事，例如，\BibleQuote{亚当听见了上主天主散步的声音}和\BibleQuote{我现在要下去，看看他们是否全照那达到我的呼声做了；如果没有，我会知道}以及\BibleQuote{现在我知道你敬畏天主}等等。\par

厄兰: — 对此我可能会回答，各种\OriginalTerm{经世}{œconomies}之间有重大区别。在旧约中有言语的\OriginalTerm{经世}{œconomy}；在新约中有行为的\OriginalTerm{经世}{œconomy}。\par

正统者: — 那么你的对手会问，是什么行为？\par

厄兰: — 他会立刻听到降生成人的行为。因为天主子成为人后，在言语和行为上，时而显示肉身，时而显示天主性：正如在所引用的段落中，祂显示了肉身和灵魂的软弱，即恐惧的感觉。\par

正统者: — 但如果他接着说，\Quote{祂没有取灵魂，只取了身体；因为天主性代替灵魂与身体结合，行使了灵魂的一切功能}，你用什么论证来应对他的异议？\par

厄兰: — 我能从圣经中提出证据，表明天主圣言不仅取了肉身，也取了灵魂。\par

正统者: — 我们在圣经中可以找到什么证据呢？\par

厄兰: — 你没有听过主说\BibleQuote{我有权放下它，也有权再取回它……我自愿放下它，为再取回它}吗？又说\BibleQuote{现在我心神烦乱}。又说\BibleQuote{我的心灵忧闷得要死}，以及达味的话，经伯多禄解释为\BibleQuote{祂的灵魂没有被撇在阴间，祂的肉身也没有见到腐朽}。这些和类似的段落明确指出，天主圣言不仅取了身体，也取了灵魂。\par

正统者: — 你引用这个见证非常恰当和合适，但你的对手可能会回答说，甚至在降生之前，天主就对犹太人说过，\BibleQuote{守斋、圣日和节期，我的灵魂憎恨}。然后他可能会进一步辩称，如同在旧约中提到灵魂，虽然祂没有灵魂，同样在新约中也是如此。\par

厄兰: — 但他会再次被告知，圣经在谈论天主时，甚至提及身体的部位，如\BibleQuote{侧耳倾听}、\BibleQuote{睁眼观看}、\BibleQuote{上主亲口说了}、\BibleQuote{你的手造了我、塑造了我}以及无数其他段落。\par

如果我们在降生之后被禁止将灵魂理解为灵魂，那么同样禁止将身体理解为身体。这样，\OriginalTerm{经世}{œconomy}的伟大奥迹将被发现只是想象；我们将与这些幻想的发明者\Person{玛尔西雍}、\Person{瓦伦提努}和\Person{摩尼}毫无区别。\par

正统者: — 但如果\Person{阿波利纳里}的一位追随者突然介入我们的讨论，问说\Quote{最卓越的先生，你说基督取了什么灵魂？}你会如何回答？\par

厄兰: — 我首先要说，我只知道人的一种灵魂；然后我会回答，\Quote{但如果你算有两种灵魂，一种理性的，另一种非理性的，我说所取的灵魂是理性的。你的似乎是非理性的，因为你认为我们的救恩不完整。}\par

正统者: — 但假设他要求你所说的话的证据？\par

厄兰: — 我非常容易就能给出。我将引用福音作者的圣言：\BibleQuote{孩童耶稣渐渐长大，心神强健，天主的恩宠常在他身上}以及\BibleQuote{耶稣在智慧和身量上，并在天主和人前的恩爱上，日渐增长}。我应该说，这些与天主性无关，因为身体在身量上增长，而智慧方面是灵魂——不是非理性的，而是理性的。所以，天主圣言取了理性的灵魂。\par

正统者: — 善哉！你勇敢地突破了敌人三重方阵；但那个联合，以及著名的混合与混淆，不仅在两方面，而是在三方面，你都拆散了；你不仅指出了天主性与人性的区别，而且以两种方式区分了人性，指出灵魂是一回事，身体是另一回事，因此按照我们的论点，我们的救主耶稣基督不再是两性，而是三性可以理解。\par

厄兰: — 是的；因为你不是说过，除了身体的本性之外，还有另一种灵魂的实体吗？\par

正统者: — 是的。\par

厄兰: — 那么，这个论点在你看来如何荒谬呢？\par

正统者: — 因为当你反对两性时，你却承认了三性。\par

厄兰: — 与对手的争辩迫使我们如此，因为除了从圣经中引证这些点的证据，还能以其他什么方式与那些否认取肉身、或取灵魂、或取理性的人辩论呢？又怎能驳斥那些急于贬低独生子天主性的人，除非指出圣经有时\OriginalTerm{按神学}{theologically}说话，有时\OriginalTerm{按经世}{œconomically}说话。\par

\Emph{正统者}— 你现在说的对。这正是我，不，正是所有完整持守宗徒规矩的人所说的。你自己已成为我们教义的支持者。\par

\Emph{厄兰}— 我拒绝承认两个儿子，又怎能支持你的呢？\par

\Emph{正统者}— 你什么时候听见过我们肯定两个儿子？\par

\Emph{厄兰}— 主张两个本性的人就是在主张两个儿子。\par

\Emph{正统者}— 那么你主张三个儿子，因为你说过三个本性。\par

\Emph{厄兰}— 别无他法能应对对手的论证。\par

\Emph{正统者}— 听听我们也说同样的话；因为你我都面对同样的对手。\par

\Emph{厄兰}— 但我不主张结合后有两个本性。\par

\Emph{正统者}— 然而刚才你在结合产生许多后代后仍用了同样的话。不过请向我们解释，你在什么意义上主张结合后只有一个本性。你是说由两者产生的一个本性，还是一个本性在另一个毁灭后存留？\par

\Emph{厄兰}— 我坚持天主性存留，而人性被它吞没。\par

\Emph{正统者}— 这一切都是外邦人的神话和摩尼教徒的蠢话。我甚至羞于提及这些东西。希腊人有他们的神祇吞没，摩尼教徒则写过光明之女。但我们拒绝这种既荒谬又不敬的教导，因为一个绝对、单纯、包含宇宙、不可接近、无限的本性，怎能吸收它所取的本性呢？\par

\Emph{厄兰}— 如同大海接纳一滴蜂蜜，因为那滴蜂蜜一与海水混合便立刻消失。\par

\Emph{正统者}— 海与滴在数量上不同，尽管性质相同；一个最大，一个最小；一个甜，一个苦；但在其他方面你会发现非常密切的关系。两者的本性都是潮湿、液态、流动的。两者都是受造的。两者都无生命，但都被称为物体。因此，这些同源的本性发生混合，并且一个被另一个湮没，并无荒谬之处。相反，在我们面前的情形中，差异是无限的，大到找不到任何实物的比喻。不过，我将努力向你指出几个物质混合而不混淆、且保持完好的例子。\par

\Emph{厄兰}— 这世上谁曾听说过不混合的混合物？\par

\Emph{正统者}— 我将努力让你承认这一点。\par

\Emph{厄兰}— 倘若你将要提出的证明为真，我们不会反对真理。\par

\Emph{正统者}— 那么就回答吧，根据论证在你看来是否合理，可表示反对或赞同。\par

\Emph{厄兰}— 我会回答。\par

\Emph{正统者}— 在你看来，光线升起时是否充满整个大气，除了那些被关在洞穴中的人可能依旧缺少它？\par

\Emph{厄兰}— 是的。\par

\Emph{正统者}— 那么全部光线是否弥散在整个大气中？\par

\Emph{厄兰}— 到目前为止我同意你。\par

\Emph{正统者}— 混合物不是也弥散在其所支配的一切之中吗？\par

\Emph{厄兰}— 当然。\par

\Emph{正统者}— 但是，这个被光照的大气，我们不是把它看作光并称之为光吗？\par

\Emph{厄兰}— 正是如此。\par

\Emph{正统者}— 然而，当光线存在时，我们有时会感受到潮湿和干燥；常常感受到热和冷。\par

\Emph{厄兰}— 是的。\par

\Emph{正统者}— 光线离开后，大气便独自存留。\par

\Emph{厄兰}— 对。\par

\Emph{正统者}— 再考虑这个例子。当铁与火接触时，它被烧红。\par

\Emph{厄兰}— 当然。\par

\Emph{正统者}— 火是否弥散在整个铁质中？\par

\Emph{厄兰}— 怎么？\par

\Emph{正统者}— 那么，完全的联合以及普遍弥散的混合，为什么没有改变铁的本性？\par

\Emph{厄兰}— 但它完全改变了它。现在不再被视为铁，而是被视为火，并且确实具有火的主动属性。\par

\Emph{正统者}— 但铁匠不是称它为铁，把它放在砧上，用锤子敲打它吗？\par

\Emph{厄兰}— 毫无疑问。\par

\Emph{正统者}— 那么铁的本性并没有因与火接触而受损。如果在自然物体中可以找到不混淆混合的例子，那么对于那既不知腐朽也不知改变的本性，竟持有混淆和毁灭所取本性的观念，就完全是愚蠢的，尤其当这个本性被取来是为了给人类带来祝福时。\par

\Emph{厄兰}— 我所主张的不是所取本性的毁灭，而是它转变为天主性的实体。\par

\Emph{正统者}— 那么人类就不再像以前那样受限了？\par

\Emph{厄兰}— 不。\par

\Emph{正统者}— 它何时经历这种改变？\par

\Emph{厄兰}— 在完全结合之后。\par

\Emph{正统者}— 你为这指定什么时间？\par

\Emph{厄兰}— 我一再说过，是受孕的那一刻。\par

\Emph{正统者}— 然而受孕后，他在母腹中是个未出生的婴儿；出生后，他是个婴儿，被称为婴儿，受到牧羊人的朝拜，同样地，他成了孩子，天使也这样称呼他。你承认这一切吗？或者你认为我在编造神话？\par

\Emph{厄兰}— 这在\WorkTitle{神圣福音}的记载中有教导，无可辩驳。\par

\Emph{正统者}— 现在让我们考察接下来的事。我们承认，不是吗，主受了割损？\par

\Emph{厄兰}— 是的。\par

\Emph{正统者}— 那是谁的割损？是肉身的，还是天主性的？\par

\Emph{厄兰}— 肉身的。\par

\Emph{正统者}— 那么，智慧与身量的成长和增加是谁的呢？\par

\Emph{厄兰}— 这当然不适用于天主性。\par

\Emph{正统者}— 饥饿和口渴也不适用？\par

\Emph{厄兰}— 不。\par

\Emph{正统者}— 走路、疲倦、睡觉也不适用？\par

\Emph{厄兰}— 不。\par

\Emph{正统者}— 如果结合发生在受孕之时，而这一切发生在受孕和出生之后，那么结合之后，人性并没有失去其本性。\par

\Emph{厄兰}— 我没有准确表达我的意思。是在从死者中复活之后，肉身才经历了转变为天主性的变化。\par

\Emph{正统者}— 那么复活之后，一切显示其本性的东西都不存留在它里面了？\par

\Emph{厄兰}— 如果存留，那神圣的改变就没有发生。\par

\Emph{正统者}— 那么他怎能将自己手和脚给不信的门徒看？\par

\Emph{Eran.}——正如祂在门关上时进来一样。\par

\Emph{Orth.}——但祂在门关上时进来，正如祂从母腹中出来，童贞女的锁和闩未开启，也正如祂在海面上行走。那么，按照你的论点，连那时也没有发生本性的变化吗？\par

\Emph{Eran.}——\Person{主}向\Person{宗徒}们显示祂的双手，如同祂与\Person{雅各伯}角力一样。\par

\Emph{Orth.}——不；\Person{主}不允许我们按这个意思理解。门徒们以为看见了幽灵，但\Person{主}打消了这个念头，并显示了肉身的本性，因为祂说：\BibleQuote{你们为什么惊慌？为什么心里起疑念？看我的手和脚，就是我自己；你们摸我看看，因为幽灵没有肉和骨，如同你们看见我有的。}请注意用语的确切性。祂不是说\Quote{不是肉和骨}，而是说\Quote{没有肉和骨}，为要指出拥有者与被拥有者的本性是分开和区别的。同样，那取者和被取者是分开和区别的，而\Person{基督}在两者中被视为一体。因此拥有部分与拥有部分完全不同；然而并不将在这两者中被思想的那个人分成两个位格。事实上，当门徒们还在疑惑时，\Person{主}要求食物，拿起并吃了，并非只是在表象中消费食物，也不是满足身体的需要。\par

\Emph{Eran.}——但必须接受两者之一：要么祂因需要而吃，要么不需要而看似吃，并非真正吃食物。\par

\Emph{Orth.}——祂如今成为不朽的身体不需要食物。关于复活的人，\Person{主}说：\BibleQuote{他们不娶也不嫁，却像天使一样。}然而宗徒们作证祂吃了食物，因为真福\Person{路加}在\BibleBook{}{宗徒大事录}的序言中说：\BibleQuote{祂与宗徒们聚集在一起，\Person{主}命令他们不要离开\Place{耶路撒冷}。}而极其神圣的\Person{伯多禄}更明确地说：\BibleQuote{祂从死者中复活后，曾与祂一同吃喝的人。}因为吃是现世生活之人的特征，所以\Person{主}必须通过吃和喝向那些不承认肉身真实复活的人证明复活。祂在\Person{拉匝禄}和\Person{雅依洛}的女儿的事上采用了同样的方法。因为祂复活了后者后，命令给她一些东西吃，祂让\Person{拉匝禄}与祂同席，从而显示了复活的真实性。\par

\Emph{Eran.}——如果我们承认\Person{主}真的吃了，那么让我们承认复活后所有的人都吃食物。\par

\Emph{Orth.}——救主通过某种\OriginalTerm{经纶}{œconomy}所做的事不是自然的规则和法则。这可由祂通过经纶做了其他事来证明，而这些事决不会发生在复活的人身上。\par

\Emph{Eran.}——你是什么意思？\par

\Emph{Orth.}——复活的人的身体不会成为不朽和不可朽坏的吗？\par

\Emph{Eran.}——神圣的\Person{保禄}如此教导我们。他说：\BibleQuote{它播种在可朽坏中，复活在不朽坏中；它播种在羞辱中，复活在光荣中；它播种在软弱中，复活在能力中；它播种了属生灵的身体，复活了属神身体。}\par

\Emph{Orth.}——但是那复活所有人身体的\Person{主}，使身体完整无缺（因为在复活的人中没有肢体跛脚、眼睛失明），却在祂自己的身体上留下了钉痕和肋旁的伤口，这事的见证者既是\Person{主}自己，也是\Person{多默}之手。\par

\Emph{Eran.}——不错。\par

\Emph{Orth.}——那么，如果复活后\Person{主}既吃了食物，又向门徒显示了祂的手和脚以及其中的钉痕，并显示了带有伤口印记的肋旁，对他们说：\BibleQuote{你们摸我看看，因为幽灵没有肉和骨，如同你们看见我有的。}那么，祂复活后，祂身体的本性得以保存，并没有变成另一种实体。\par

\Emph{Eran.}——那么复活后它是可朽坏和可受苦的？\par

\Emph{Orth.}——绝不；它是不朽坏、不受伤、不死的。\par

\Emph{Eran.}——如果它是不朽坏、不受伤、不死的，它已变成了另一种本性。\par

\Emph{Orth.}——因此所有人的身体都会变成另一种实体，因为所有的人都将是不朽坏不死的。或者你没有听过\Person{宗徒}的话：\BibleQuote{因为这可朽坏的必须穿上不朽坏，这可死的必须穿上不死？}\par

\Emph{Eran.}——我听过。\par

\Emph{Orth.}——因此本性存留，但它的可朽坏变成了不朽坏，它的可死变成了不死。但让我们这样看问题：我们称一个有病和一个健康的身体同样为身体。\par

\Emph{Eran.}——毫无疑问。\par

\Emph{Orth.}——为什么？\par

\Emph{Eran.}——因为两者都分享同一实体。\par

\Emph{Orth.}——然而我们在他们身上看到非常大的差别，因为一个是完整、完美、未受伤的；另一个要么失去了一只眼睛，要么断了一条腿，要么遭受了其他痛苦。\par

\Emph{Eran.}——但健康和疾病都属于同一本性。\par

\Emph{Orth.}——所以身体被称为\OriginalTerm{实体}{substance}；疾病和健康被称为\OriginalTerm{偶性}{accident}。\par

\Emph{Eran.}——当然。因为这些东西是身体的偶性，并又停止存在。\par

\Emph{Orth.}——同样，可朽坏和死亡必须被称为\OriginalTerm{偶性}{accidents}，而不是\OriginalTerm{实体}{substances}，因为它们也是偶性并停止存在。\par

\Emph{Eran.}——对。\par

\Emph{Orth.}——所以\Person{主}的身体不朽坏、不受伤、不死的复活了，被天上的万有崇拜，但它仍然是一个身体，拥有它原有的限制。\par

\Emph{Eran.}——在这一点上你似乎说得对，但在它被接升天后，我想你不会否认它已变成了\OriginalTerm{天主性}{Godhead}的本性吧？\par

\Emph{Orth.}— 我不会仅凭人的论证就这样说，因为我并不如此轻率，敢于说出任何圣经沉默不言的事。但我听见了神圣的保禄呼喊：\BibleQuote{天主已定了一日，要藉着祂所指定的那一位，按正义审判天下，祂已使祂从死者中复活，给众人做了凭据。}我也从圣天使那里得知，祂要怎样来，就像门徒看见祂升天一样。现在他们看见祂的本性不是无限的。因为我听见了主的话：\BibleQuote{你们要看见人子乘着天上的云彩降临。}我承认，被人看见的事物是有限的，因为无限的本性是看不见的。此外，坐在光荣的宝座上，把羔羊放在右边，把山羊放在左边，也表明是有限的。\par

\Emph{Eran.}— 那么，即使在降生成人之前，祂也不是无限的，因为先知看见祂被色辣芬围绕。\par

\Emph{Orth.}— 先知并没有看见天主的本体，而是看见了按他能力所呈现的某种显现。然而，在复活之后，全世界都将看见审判者那可见的本性。\par

\Emph{Eran.}— 你曾答应不引无证据的论证，但你却引入了适合我们的论证。\par

\Emph{Orth.}— 这些事我是从圣经学来的。我听见了匝加利亚先知的话：\BibleQuote{他们要仰望他们所刺透的那一位。}除非钉死祂的人认出他们所钉死的那一本性，否则事件如何应验预言？我又听见了胜利的殉道者斯德望的呼喊：\BibleQuote{看，我看见天开了，人子站在天主的右边。}他看见了可见的，而非不可见的本性。\par

\Emph{Eran.}— 这些事是这样记载的，但我认为你无法证明，在升天之后，那身体仍被默感作者称为身体。\par

\Emph{Orth.}— 刚才所说的话已经十分清楚地表明了那是身体，因为被看见的就是身体。但我仍要向你指出，即使在升天之后，主的身体仍被称为身体。请听宗徒的教导：\BibleQuote{我们的家乡原是在天上，我们等待主耶稣基督救主从那里降来，祂将改变我们卑贱的身体，使之相似祂光荣的身体。}它没有变成另一种本性，而是仍然作为身体，却充满神圣的光荣，发出光芒。圣人们身体的形状也要相似它。但如果它变成了另一种本性，他们的身体也同样要改变，因为他们要相似它。但如果圣人们的身体保存其本性的特征，那么主的身体也同样保持其本性不变。\par

\Emph{Eran.}— 那么，圣人们的身体会与主的身体相等吗？\par

\Emph{Orth.}— 在不朽和永生方面，他们也将分享。此外，在光荣上他们也将参与，正如宗徒所说：\BibleQuote{如果我们和祂一同受苦，也必和祂一同受光荣。}在量上可发现巨大的差异，如同太阳与星星之间的差异，更恰当地说，如同主人与仆人、施光者与受光者之间的差异。然而，祂将自己的名分赐给了祂的仆人，祂是光，也称祂的圣徒为光，因为祂说：\BibleQuote{你们是世界的光。}祂被称为\OriginalTerm{正义的太阳}{Sun of Righteousness}，论及祂的仆人时祂说：\BibleQuote{那时，义人要在他们父的国里发光如同太阳。}因此，圣人们的身体是按质，而不是按量相似主的身体。现在我已清楚地告诉了你你所要求的。此外，如果你同意，让我们从另一个角度来看这件事。\par

\Emph{Eran.}— 正如谚语所说，应当\Quote{翻遍每一块石头}来寻找真理，尤其当涉及神圣教义时。\par

\Emph{Orth.}— 那么请告诉我：那些由司祭举行礼仪时奉献给天主的\OriginalTerm{奥迹象征}{mystic symbols}，是象征什么的？\par

\Emph{Eran.}— 主的身体和血。\par

\Emph{Orth.}— 是真实的身体还是不是？\par

\Emph{Eran.}— 真实的。\par

\Emph{Orth.}— 很好。因为必须有\OriginalTerm{原型}{archetype}的图像。所以画家模仿自然，描绘可见物体的形象。\par

\Emph{Eran.}— 对。\par

\Emph{Orth.}— 那么，如果神圣奥迹是真实身体的\OriginalTerm{反型}{antitypes}，那么即使现在，主的身体仍是身体，没有转变为\OriginalTerm{神性}{Godhead}本性，而是充满神圣光荣。\par

\Emph{Eran.}— 你恰当地引入了神圣奥迹的主题，因为由此我将能向你显示主的身体转变为另一种本性。现在回答我的问题。\par

\Emph{Orth.}— 我会回答。\par

\Emph{Eran.}— 在\OriginalTerm{司祭呼求}{priestly invocation}之前，所奉献的礼物你称为什么？\par

\Emph{Orth.}— 公开说出是不合适的；也许有一些未进教者在场。\par

\Emph{Eran.}— 请用隐晦的方式回答。\par

\Emph{Orth.}— 某种谷类食物。\par

\Emph{Eran.}— 另一个象征我们如何称呼？\par

\Emph{Orth.}— 这也是一个常见的名称，指一种饮料。\par

\Emph{Eran.}— \OriginalTerm{祝圣}{consecration}之后，这些你又如何称呼？\par

\Emph{Orth.}— 基督的身体和基督的血。\par

\Emph{Eran.}— 你相信你领受了基督的身体和血吗？\par

\Emph{Orth.}— 我相信。\par

\Emph{Eran.}— 那么，正如主身体和血的象征在司祭呼求之前是一种东西，在呼求之后被改变成为另一种东西；同样，主的身体在升天之后被转变为\OriginalTerm{神性本体}{divine substance}。\par

\Emph{Orth.}——你已自投罗网。因为即使在祝圣之后，奥迹象征并未失去其本性；它们仍保持原有的实体、形状和形态，仍同先前一样可看见、可触摸。然而它们被视为所变成的，并被信为如此，并被敬拜为所信之物。那么，将图像与原形相比，你就会看出相似之处，因为类型必须与事实相似。因为那身体保持其原有的形态、形状和界限，总而言之，身体的实体；但在复活之后，它成为不朽的，超越朽坏；它成为配得坐在右边之位；它被每一个受造物敬拜，因为被称为主的本性之体。\par

\Emph{Eran.}——是的，而且奥迹象征改变了它原先的名称；它不再被称为它以前的名字，而是被称为身体。所以，事实本身必须被称为天主，而不是身体。\par

\Emph{Orth.}——你似乎不知道——因为祂不仅被称为身体，甚至被称为生命之粮。所以主亲自使用这个名字，而那一身体我们称为神性的身体、生命之赐予者、属主和上主，教导说它并非人人共有，而是属于我们的主耶稣基督，祂是天主也是人。因为\BibleQuote{耶稣基督昨天、今天、直到永远，常是一样的。}\par

\Emph{Eran.}——关于这一点，你说得很多，但我追随那些在教会中古代发光照耀的圣人；那么，如果你能的话，请向我指出，在他们的著作中，在结合之后区分本性。\par

\Emph{Orth.}——我来给你读他们的作品，我相信你会惊讶于他们在对抗不虔敬的异端者的斗争中，在他们的著作中插入的无数关于区分的提及。现在请听那些我已经引用过证言的人，他们公开而明确地谈论这些要点。\par

\Emph{圣依纳爵，安提约基雅主教，殉道者的证言：——}\par

从\WorkTitle{致斯米纳人书}：\Quote{我承认并相信祂在复活后仍以肉身存在：当祂来到与伯多禄在一起的人那里时，祂对他们说：\InnerQuote{你们摸摸看，摸摸我，知道我不是无体的鬼怪。}他们就立刻摸祂并相信了。}\par

同一作者，同一书信：——\par

\Quote{祂复活后，与他们一同吃喝，因为祂有肉身，虽然在灵性上与父合而为一。}\par

\Emph{圣依勒内，里昂古代主教的证言：——}\par

从他所著的\WorkTitle{驳斥异端}第三卷（第二十章）。\par

\Quote{正如我们以前所说，祂将人与天主结合。因为如果人没有战胜人的仇敌，敌人就不会被正确击败；另外，如果天主没有赐予救恩的恩赐，我们也不能安全地拥有它。如果人没有与天主结合，他就不能分享不朽。因为天主与人的中保，藉着祂与两者的亲密关系，有必要使两者进入友谊与和谐，使人亲近天主，并使天主为人所认识。}\par

同一作者，同一著作的第三卷（第十八章）：——\par

\Quote{所以在他的书信中又说：\InnerQuote{凡信耶稣是基督的，是由天主而生的。}承认同一位耶稣基督，天上的门为祂敞开，因为祂取了肉身。祂要在祂曾受苦的同一肉身中降临，揭示父的光荣。}\par

同一作者，第四卷（第七章）：——\par

\Quote{正如依撒意亚说：\InnerQuote{祂要使雅各伯的后代扎根。以色列将开花发芽，使世界充满果实。}这样，祂的果实散布到全世界，那些从前结出好果实的人（因为从他们中产生了在肉身中的基督和宗徒们）被遗弃和移除。如今他们不再适合结出果实。}\par

同一作者，同一卷（第五十九章）：——\par

\Quote{祂也审判厄俾翁派的人。除非是天主在地上实行了他们的救恩，他们如何能得救？或者，除非天主来到人中间，人如何能来到天主面前？}\par

同一作者，同一卷（第六十四章）：——\par

\Quote{那些宣讲厄玛奴耳出自童贞女的人，表明了天主圣言与祂受造物的结合。}\par

同一作者，同一著作（第五卷第一章）：——\par

\Quote{现在这些事不是看似发生，而是在本质真理中发生。因为如果祂看似是人，而祂并非人，那么圣神就不是继续祂的真实所是；因为圣神是不可见的；在祂里面也没有真理，因为祂不是祂所显现的样子。我们以前说过，亚巴郎和其他先知在预言中看见祂，预言将要发生的事会成为实际所见。那么如果现在祂也显现为那样的性格，而实际上祂并非所显现的样子，那么人就会得到一种预言性的异象，我们仍须等待另一个来临，在那来临中祂将真正是祂现在在预言中所显现的。现在我们已经证明，说祂只是看似显现，与说祂从玛利亚什么都没取有，这两者之间没有区别，因为如果祂没有在自身中更新亚当的旧创造，祂就根本没有真正拥有救赎我们的血肉和血液。所以瓦伦提努的派别如此教导是徒劳的，为要抛弃肉身的生命。}\par

\Emph{圣希玻里，主教和殉道者，从他的著作\WorkTitle{论塔冷通的分配}的证言：——}\par

\Quote{有人可能会说，这些人和那些持不同观点的人是邻居，因为他们犯了同样的错误，因为他们要么承认基督在世时仅是凡人，否认祂天主性的塔冷通，要么承认祂是天主，另一方面却否认祂是人，说祂用不真实的显现欺骗了观看者的眼睛；并且祂不是作为人而穿戴了人性，而只是一种虚幻的假象，正如马尔强、瓦伦提努和诺斯替派所教导的，他们将圣言从肉身中撕扯开，拒绝了那唯一的塔冷通——降生成人。}\par

同一作者，从他写给某位王后的书信：——\par

\Quote{祂称祂为\InnerQuote{长眠者中初熟的果实}，因为祂是\InnerQuote{从死者中首生的}，而祂复活后，为表明复活的身体与受死的身体相同，当门徒们怀疑时，祂叫多默来，说：\InnerQuote{你伸手摸摸看，因为神魂没有血肉，如同你们看见我有的。}}\par

同一位，取自其论厄耳卡纳和亚纳的讲道：—\par

\Quote{因此，一年中的三个季节预表了救主自己，以应验关于祂的奥秘。在逾越节，祂显明自己是注定被祭献的羔羊，并显示真正的逾越节，如宗徒所说：\BibleQuote{基督、天主、我们的逾越节，已被祭献，为了我们。}在五旬节，祂首先升入天堂，向天主呈上人作为礼物，以宣示天国。}\par

同一位，取自其论大圣咏的著作：—\par

\Quote{祂从最深的地狱中，将首先由泥土形成、因死亡束缚而迷失的人拉上来；祂从上面降下，抬起了那在下的人；祂成了死者的传福音者、灵魂的赎世者、墓葬之人的复活者。祂就是在自己内成为被征服之人的扶助者，如同首生的圣言一样的人；祂在童贞女内探访了首造的亚当；祂在母胎中追寻了属神的；祂以永生者之身寻访了因悖逆而死的；祂以属天者之身召叫属地的到上界，尊贵者意欲藉自己的顺服使奴隶得自由；祂使那化为尘土、成为蛇肉的人变得如金刚石般坚硬；祂使那挂在木头上的人成为胜过那曾征服祂的主，并藉着这木头得胜。}\par

同一位，取自同一著作：—\par

\Quote{那些如今在肉身中不认得天主子的人，有一天当祂带着荣耀作为审判者降临时，必会认识祂，尽管现在祂在一个卑微的身体中受苦。}\par

同一位，取自同一著作：—\par

\Quote{再者，宗徒们第三天来到坟墓时，没有找到耶稣的身体，正如以色列子民上山，找不到梅瑟的坟墓。}\par

同一位，取自其圣咏第二篇释义：—\par

\Quote{祂来到世界时，显现为天主和人。祂的人性易于察觉，因为祂饥饿、疲倦、劳累口渴、恐惧逃跑、祈祷忧伤、枕着枕头入睡、祈求苦杯离开祂、痛苦挣扎出汗、被天使坚强、被犹达斯出卖、被盖法侮辱、被黑落德藐视、被比拉多鞭打、被士兵戏弄、被犹太人钉在十字架上、大声喊叫将灵魂交托给父、断气时垂下头、被长矛刺透肋旁、裹在细麻布里、安放在坟墓中，第三天由父复活。祂的神性同样明显：被天使朝拜、被牧童瞻仰、被西默盎期待、被亚纳证实、被贤士寻觅、被星星指认、在婚宴上将水变成酒、斥责被狂风翻腾的海、在海面上行走、使生来瞎眼的人看见、使死了四天的拉匝禄复活、行了许多各样的奇迹、赦免罪过、并将权能赐给门徒。}\par

同一位，取自其圣咏第二十四篇著作：—\par

\Quote{祂来到天门，天使与祂同行，诸天的门关闭着。因为祂尚未升天。现在，肉身首次升天，出现在天上的掌权者面前。于是，圣言从那些奔跑在主和救主前面的天使那里向掌权者发出：\InnerQuote{王侯啊，举起门来，永恒的城门，请被举起，荣耀的君王要进来。}}\par

\Emph{圣欧斯塔修斯，安提约基雅主教与宣信者的见证。}\par

取自其《论圣咏标题》的著作：—\par

\Quote{祂预言自己要坐在神圣的宝座上，表明因那常居于祂内的天主，祂已被安置在与圣神相同的宝座上。}\par

同一位，取自其论灵魂的著作：—\par

\Quote{在祂受难之前，祂每次都预言了自己身体的死亡，说祂会被交给大司祭的父亲，并宣告十字架的胜利。受难之后，当祂第三天从死者中复活时，门徒们对祂的复活心存疑虑，祂就以其真实的身体显现给他们，承认祂有完整的肉和骨头，让他们观看祂被刺的肋旁，并给他们看钉痕。}\par

同一位，取自其讲道\InnerQuote{上主在祂道路的起始就创造了祂}：—\par

\Quote{保禄没有说\InnerQuote{同化于天主子}，而是说\InnerQuote{同化于祂儿子的肖像}，为指出圣子与其肖像之间的区别。因为圣子带有祂父卓越性的神圣标记，是祂父的肖像；既然相似者由相似者所生，后裔便显现为父母的真实肖像。但祂所穿戴的人性却是圣子的肖像，如同不同颜色的肖像被画在蜡上，有些是手工制成，有些则藉自然和相似。此外，真理的法律本身也宣告这一点：无体的智慧之神并不同化于有身体的人，而是那由圣神成为人的明确肖像，带着与其余人相同数目的肢体，并穿着相似的形状。}\par

同一位，取自同一著作：—\par

\Quote{祂说身体同化于人的身体，这在祂致斐理伯人的书信中更清楚地教导：\BibleQuote{我们的家乡在天上，我们也等待救主，主耶稣基督，祂要改变我们卑贱的身体，使之相似祂光荣的身体。}如果祂藉着改变人卑贱身体的形状使之相似祂自己的身体，那么对手的虚假教训就被证明是完全无价值的。}\par

同一位，取自同一著作：—\par

\Quote{但正如因由童贞玛利亚所生而说祂由女人成为人，同样，祂被描述为属于法律之下，因祂有时遵行法律的规定，例如当祂的父母热心催促祂受割损，祂尚是八天大的婴孩时，如\Person{路加}圣史所记载，之后\BibleQuote{他们带祂到主面前}，\BibleQuote{献上取洁的祭品}，\BibleQuote{照主的法律所吩咐的，献上一对斑鸠或两只雏鸽。}因此，按法律为祂献上取洁的礼物，且祂在第八天受了割损，宗徒很恰当地写道祂就这样被置于法律之下。当然，圣言并不受法律约束（如我们诽谤性的对手所臆测），因为祂本身就是法律；天主也无需取洁的牺牲，因祂一口气就能洁净并圣化一切。但祂从童贞女取了人的肢体，成为法律的臣民，并按长子的仪式受了取洁，并非因为祂自身需要这样的仪式，而是为救赎那些卖身于诅咒刑罚之下、受法律奴役的人。}\par

\Emph{圣亚大纳削，亚历山大主教的证言。}\par

摘自其\WorkTitle{第二篇反异端讲道集}：—\par

\Quote{我们本不能从罪和诅咒中被救赎，除非圣言所穿戴的血肉按本性是人的血肉，因为对于不属于我们的事物，我们毫无共通之处；同样，人也不可能成为天主，除非那成了血肉的圣言按本性是由圣父所生，且确实是并真正属于圣父的。这结合的性质是这样的：使本性的天主可与本性的人结合，如此，人的救恩与\OriginalTerm{神化}{deification}才能稳妥。因此，凡否认祂是由圣父本性所生，并承认祂是圣父的\OriginalTerm{本体}{substance}之子的，也否认祂从童贞玛利亚取了真实的人的血肉。}\par

摘自其\WorkTitle{致厄皮克提督书}：—\par

\Quote{如果因救主的身体源自玛利亚，且圣经描述其为人的身体，他们就幻想以\OriginalTerm{四位一体}{quaternity}取代三一体，仿佛身体带来某种增加，那他们就大错特错；他们将受造物与造物主等同，并认为天主性是可以被增添的。他们未能看到：圣言成为血肉并非为了给天主性增添什么，而是为使血肉能复活。圣言从玛利亚出来，并非为\OriginalTerm{壮大}{aggrandisement}圣言，而是为救赎人类。他们怎能认为那被圣言赎出并\OriginalTerm{赋予生命}{quickened}的身体能对赋予它生命的圣言在天主性方面增添什么呢？}\par

摘自同一书信：—\par

\Quote{应告诉他们，如果圣言曾是受造物，那么受造物就不会取一个身体来赋予它生命。因为受造物自身需要救恩，又能从另一个受造物得到什么帮助呢？但圣言，祂自己是\OriginalTerm{造物主}{Creator}，成了受造物的\OriginalTerm{创造者}{maker}，因此在这时期的圆满中，祂将\OriginalTerm{受造物}{creature}与自己结合，以便再一次作为造物主更新它，并能重新创造它。}\par

摘自较长的\WorkTitle{De Fide}：—\par

\Quote{我们也关于\InnerQuote{坐在我的右边}这话加这一点，即它们是就主的身体说的。因为若\InnerQuote{主说：我岂不充满天地？}如\Person{耶肋米亚}所说，且天主包含一切，不被任何事物所包含，那么祂坐在怎样的宝座上呢？因此，那身体就是祂对之所说\BibleQuote{坐在我的右边}的，魔鬼及其邪恶势力、犹太人和外邦人都是它的仇敌。通过这身体，祂也成为并被称为大司祭和宗徒，借着祂赐给我们的\OriginalTerm{奥迹}{mystery}，祂说：\BibleQuote{这是我的身体，为你们}，以及\BibleQuote{我的血，新约的血（非旧约的），为你们而倾流。}}如今天主性既无身体也无血；但祂从玛利亚所取的人性却是它们的缘由，宗徒们曾论及祂说：\BibleQuote{纳匝肋人耶稣，是天主在你们中间所证明的人。}\par

摘自其\WorkTitle{反亚略派论}：—\par

\Quote{当他说：\BibleQuote{因此，天主又极高地举扬了祂，赐给祂那超越一切名号的名字}，他指的是\OriginalTerm{身体的殿}{temple of the body}，而非天主性；因为至上者并未被高举，而是至上者的血肉被高举，且祂将那超越一切名号的名字赐给了至上者的血肉。圣言也不是作为恩惠而领受天主的称号，而是祂的血肉与祂一样被视为天主。}\par

摘自同一著作：—\par

\Quote{当耶稣说：\BibleQuote{圣神还没有赐下，因为耶稣还没有受光荣}，他说祂的血肉尚未受光荣；因为光荣的主不受光荣，而是血肉本身随祂一同升到天堂时，从主的荣耀中接受荣耀；因此他说，义子的\OriginalTerm{圣神}{spirit of adoption}尚未降到人中，因为从人中所取的\OriginalTerm{初果}{first fruits}尚未升到天堂。因此，凡圣经说圣子领受并被光荣的，都是因为祂的人性，而非祂的\OriginalTerm{天主性}{Godhead}。}\par

摘自同一著作：—\par

\Quote{因此，祂是真天主，无论是在祂成为人之前，还是祂成为天主与人的中保之后；耶稣基督，在精神上与圣父合一，在血肉中与我们合一，祂在天主与人之间做中保，祂不仅是人，也是天主。}\par

\Emph{圣盎博罗削，米兰主教的证言。}\par

摘自其\WorkTitle{信仰诠释}：—\par

我们宣认，我们的主耶稣基督，天主的独生子，在万世之前，无始地由圣父所生；在末世，同一基督由圣童贞玛利亚取得血肉，取了完美的人性，即有理性的灵魂和身体，就祂的天主性而论，与圣父同一性体；就祂的人性而论，与我们同一性体。因为两个完美性体的结合是以不可言喻的方式实现的。因此，我们承认一个基督、一个儿子、我们的主耶稣基督；我们知道，祂就祂的天主性而论，与祂自己的圣父同永，也因这同一性体，祂是万物的创造者。祂在圣童贞同意之后——当她向天使说：\Quote{我是主的婢女，愿照你的话成就于我吧}——屈尊以不可言喻的方式，用她为自己建造了一座殿宇，并通过她的受孕将这殿宇与自身结合；祂并非取用与自己性体同永、从天上带来的身体，并与自己结合，而是从我们性体的质料，即从童贞女而来。天主圣言并未变成血肉；祂的显现并非不真实；祂恒常不变地持守自己的性体，取了我们本性的初果，并将之与自身结合。天主圣言并非从童贞女取得起始，而是与自己的圣父同永；祂以无限的慈爱，屈尊将我们本性的初果与自身结合，未经历任何混合，但在两种性体中显现为同一的一位，正如经上所载：\BibleQuote{你们拆毁这座圣殿，三天之内，我要把它重建起来。}因为基督的神性，就祂所取之我的性体而言，被摧毁了；同一基督就神性而论（祂亦是万物的创造者）复活了那被摧毁的圣殿。在祂屈尊从受孕那一刻起与自身结合之后，祂从未离开过自己的殿宇，也因其对人类的不可言喻的爱而无法离开。\par

同一个基督既是可受苦难的，又是不受苦难的；就祂的人性而论是可受苦难的，就祂的天主性而论是不受苦难的。\Quote{看，看，是我，我未发生变化}——当天主圣言复活了自己的殿宇，并在这殿宇内成就了我们本性的复活与更新时，祂向门徒显示了这本性，并说：\BibleQuote{你们摸摸我，看看，因为鬼神没有骨肉，你们看我是有的。}祂不是说\InnerQuote{是}而是说\InnerQuote{有}。祂如此说，同时指向拥有者和被拥有者，为使你们明白所发生的并非混合、并非改变、并非变化，而是结合。为此，祂也显示了钉痕和枪伤，并在门徒面前进食，以各种方式使他们确信，我们本性的复活已在祂内得以更新；再者，由于祂按照祂天主性的真福性体——不变、不受苦难、不死——生活，一无所缺，祂以许可的方式允许一切可感知的事物被带到祂自己的殿宇，并凭自己的力量使之复活，藉着祂自己的殿宇使我们的本性得以完美更新。\par

\Quote{因此，那些主张基督是纯粹的人，或天主圣言是可受苦的，或变成血肉，或祂所拥有的身体是同体的，或祂从天上带来，或是不真实的；或主张天主圣言是可死的，需要从圣父那里领受复活，或祂所取的身体没有灵魂，或没有理智的人性，或基督的两个性体因混乱和混合而成为一个性体；那些否认我们的主耶稣基督是两个不相混淆的性体、却是一个位格，正如祂是一个基督、一个儿子的人——所有这些，公教会和宗徒教会都予以谴责。}\par

同一作者：——\par

\Quote{既然所有人的肉身在基督内或曾在基督内遭受苦难，怎能认为它与天主性同一性体呢？因为如果圣言与那出自泥土本质的肉身是同一性体，那么圣言与祂所取之完美的灵魂也是同一性体，因为圣言与天主是同一天性，既按照圣父的言语，也按照圣子自己的宣认，祂说：\BibleQuote{我与父原是一体。}如此，就必须认为圣父与身体是同一性体。你们为何还生亚略人的气？他们称圣子是受造物，而你们自己却宣称圣父与祂的受造物是同一性体？}\par

同一作者致格拉提安皇帝书：——\par

\Quote{让我们保持天主性与肉身之间的区分。天主的一个儿子在两者中发言，因为两个性体都在祂内。同一个基督发言，但并非总是同样，而是有时以不同的方式。观察祂何时表达神圣的荣耀，何时表达人的情感。作为天主，祂讲论天主的事，因为祂是圣言；作为人，祂谦卑地说话，因为祂在我的性体内交谈。}\par

同一主题，同一书卷：——\par

\Quote{至于我们读到荣耀的主被钉十字架的经文，我们不可认为祂是在自己的荣耀中被钉。但既然祂既是天主又是人——就祂的天主性而言是天主，就祂所取之肉身而言是人——主耶稣基督，荣耀的主，被称为被钉十字架。因为祂分有两者：即人的性与神的性。在人性中，祂经历了苦难，为要使那受苦者能被毫无区别地称为荣耀的主和人子。正如经上所载：\BibleQuote{那从天上降下来的。}}\par

同样，同一作者：——\par

\Quote{那么，让那关于言辞的虚妄问题静寂吧，正如经上所载：天主的国不在于\InnerQuote{巧言妙语}，而在于\InnerQuote{圣神之明证}。因为有一个天主子在两方面发言，因为两个性体都在祂内；但虽然祂亲自发言，祂并非总是同样发言；因为你在祂身上有时看到天主的荣耀，有时看到人的情感。作为天主，祂讲论天主的事，因为是圣言；作为人，祂讲论人的事，因为在这性体内祂说话。}\par

同一作者，出自\WorkTitle{论主的降生成人驳阿波利拿里派}一书：——\par

\Quote{但我们正在驳斥这些人的时候，又有一群人起来，断言基督的身体与祂的天主性是同一性体。地狱竟呕吐出如此可怕的亵渎？真的，亚略派更可容忍，因为他们的不信——由于这些人——反而得到加强，以致他们更强烈地否认圣父、圣子及圣神是同一实体，因为他们至少曾努力维护主的天主性与祂的肉身是同一性体。}\par

同一位（出自同一章）：—\par

\Quote{他屡次告诉我，他坚持尼西亚大公会议的信条，但在那次审查中，我们的教父们规定：天主的圣言——而非肉身——与圣父同一实体，他们承认圣言出自圣父的实体，但肉身则出自童贞玛利亚。那么，他们为何向我们举出尼西亚大公会议的名称，而实际上却在引进我们的祖先从未想过的新奇事物呢？}\par

同一位驳斥阿波利纳里：—\par

\Quote{你不可允许身体与天主性在本性上平等。即使你相信基督的身体是真实的，并把它带到祭台上进行转化，却未能区分身体与天主性的本性，我们会对你说：\InnerQuote{你若献得对，却区分得不对，你就犯了罪，安静吧。}要区分什么是自然属于我们的，什么是圣言所独有的。因为我本没有祂本有的，祂本没有我本有的，但祂取了我本有的，为使我们分享祂本有的。祂领受这个，不是为混淆，而是为成全。}\par

同一位，稍后：—\par

\Quote{愿那些说圣言的本性已转变为身体的本性的人，不要再这样说，免得按同样的解释，圣言的本性似乎已转变为罪的腐败。因为摄取者与被摄取者有区别。能力降临于童贞玛利亚，正如天使对她说：\BibleQuote{至高者的能力要庇荫你}。但所生的出自童贞玛利亚的身体，因此降生是神圣的，而受孕是属人的。所以，肉身与天主性的本性不能相同。}\par

\Emph{圣巴西略，凯撒勒雅主教的见证。}\par

出自其 \WorkTitle{感恩节讲道词}：—\par

\Quote{因此，当祂为祂的朋友哭泣时，祂显示出对人性的参与，并使我们脱离两个极端，既不容许我们在痛苦中变得过于软弱，也不容许我们对痛苦麻木不仁。正如主在消化了固体食物后感到饥饿，在身体的水分耗尽时感到口渴，在旅途使祂的神经和肌腱紧张时感到疲倦，并不是因为祂的天主性被劳苦所压，而是因为祂的身体表现出其本性的惯常症状。同样，当祂允许自己哭泣时，祂让肉身顺其自然而行。}\par

出自同一位的 \WorkTitle{驳斥尤诺米}：—\par

\Quote{我说，具有天主的形态，与具有天主的实体具有相同的效力；因为正如取了奴仆的形态，表明我们的主具有人性的实体，同样，说祂具有天主的形态，便将天主实体的特有属性归于祂。}\par

\Emph{圣额我略，纳齐安格主教的见证。}\par

出自其讲道 \WorkTitle{论新主日}：—\par

\Quote{你要相信祂将在祂荣耀的来临时再次降临，审判生者死者，不再是有肉身的，但并非没有身体。}\par

\Quote{为使祂被刺透祂的人看见，并作为天主而无粗重之质。}\par

同一位出自其 \WorkTitle{致克莱多尼书信}：—\par

\Quote{天主和人乃是两个性体，如同灵魂和身体是两个；但并非是两个子，也并非是两个男人，尽管保禄如此谈论外在的人和内在的人。总之，救主的存在的起源有两种，因为可见的与不可见的、无时间的与有时间的是有分别的，但祂并不是两个存在。天主不许。}\par

同一位出自同一篇 \WorkTitle{致克莱多尼的论述}：—\par

\Quote{若有人说，肉身如今已被弃置，天主性是赤裸无身体的，并且它不是、也不会与所摄取者一同来临，愿他被剥夺看见来临之荣耀的恩赐！因为身体现在在哪里，岂非与摄取它者同在？它肯定没有像摩尼教徒所编造的那样，被圣子吞没，以便通过羞辱而得尊荣；它也没有像声音、香气的流痕或无形闪电的闪光那样被倾倒、消散在空气中。复活后可以被触摸的能力在哪里？那一日它将被刺透祂的人看见。因为天主性本身是不可见的。}\par

同一位出自 \WorkTitle{关于圣子的第二篇讲道}：—\par

\Quote{作为圣言，祂既不服从也不不服从，因为这些性质属于那些受管辖和地位较低者；前者属于较温顺者，后者属于当受惩罚者。但在奴仆的形态中，祂迁就祂的同仆，穿上不属于祂的形态，在祂自己身上承担了我的全部和我的一切，为在祂自己里面耗尽并消灭那卑劣的部分，如同蜡被火耗尽，或地上的雾气被太阳驱散。}\par

同一位出自其 \WorkTitle{主显节讲道}：—\par

\Quote{既然祂从童贞玛利亚出来，带着两个彼此对立的因素——肉身和灵魂——其中一个被摄入天主，另一个则显示出天主性的恩宠。}\par

同一位稍后：—\par

\Quote{祂被派遣，但作为人。因为祂的本性是双重的，无疑从此以后祂按人身体的习惯而疲倦、饥饿、口渴、痛苦、流泪。}\par

同一位出自 \WorkTitle{关于圣子的第二篇讲道}：—\par

\Quote{祂被称为天主，不是圣言的天主，而是可见受造界的天主，因为那绝对是天主的，怎能成为祂的天主？同样，祂被称为圣父，不是可见受造界的圣父，而是圣言的圣父。因为祂具有双重本性。因此，一个绝对属于两者，另一个则不绝对。因为祂绝对是我们的天主，但不绝对是我们的圣父。正是这种名称的连结导致了异端者的错误。一个证明在于，当本性在思想中被区分时，名称也有区分。请听\Person{保禄}的话：\InnerQuote{我们的主耶稣基督的天主，荣耀的圣父}——对基督而言，祂是天主，对荣耀而言，祂是圣父；若两者为一，这不是由于本性，而是由于连结。还有什么比这更清楚的呢？第五，要说祂接受了生命、权柄、万国的产业、管辖一切血肉的权柄、荣耀、门徒或你愿意的任何事物；这一切都属于人性。}\par

同一位作者，同一著作：—\par

\Quote{\InnerQuote{因为天主只有一个，在天主与人之间的中保也只有一个，就是那人基督耶稣。}作为人，祂仍为我的救恩代求，因为祂保留了所取的肉身，直到祂藉降生成人的能力使我成为天主——尽管祂不再按肉身，即按肉体的情欲被认识，尽管祂没有罪。}\par

同一位作者，同一著作：—\par

\Quote{这不是很明显吗：祂作为天主知道，但祂说作为人无知？也就是说，如果有人将外表与理智感知的对象区分开。因为产生这种意见的原因在于，圣子的称谓是绝对而无关联的，并未附加祂是谁的圣子；因此，为了给这种无知赋予最虔诚的意义，我们相信它属于人性，而非神性。}\par

\Emph{\Person{圣额我略}，\Place{尼撒}主教，的证言。}\par

出自他的\WorkTitle{要理讲授}：—\par

\Quote{谁说天主性的无限被肉体的局限所理解，如同被某种容器？}\par

同一位作者，同一著作：—\par

\Quote{但如果人的灵魂因其本性的必然与身体混合，处处有权威，又何必断言天主性受肉身本性的限制？}\par

同一位作者，同一著作：—\par

\Quote{既然我们承认某种天主性相对于人性的合一与接近，又有什么阻止我们在这种接近中保留天主性的理智，相信天主性即使在人身内也超越一切限制？}\par

同一位作者，出自他的\WorkTitle{驳欧诺弥}著作：—\par

\Quote{\Person{玛利亚}之子与弟兄交谈，但独生子没有弟兄，因为独生子的名号如何在弟兄中保存？同一位基督，祂说\InnerQuote{天主是神}，却对门徒说\InnerQuote{摸我}，以表明人性可触摸，天主性不可触摸；祂说\InnerQuote{我去}表明从一处到另一处的移动，而那包含万有、\InnerQuote{藉着祂}，如宗徒所说，\InnerQuote{万有被创造，藉着祂万有得以存立}的祂，在一切现存之物中，没有任何在祂之外和超越祂自己，能与祂构成运动或移动的关系。}\par

同一位作者，同一著作：—\par

\Quote{\InnerQuote{祂被高举在天主的右边。}那么谁被高举？是卑微的还是至高者？卑微的若不是人性，又是什么？至高者除了天主性，又是什么？但天主是至高者，无需高举，因此宗徒说人性被高举，即被高举\InnerQuote{成为主和基督}。因此，宗徒用\InnerQuote{使祂成为}这个词，并非指主永恒的存在，而是指在天主右边发生的卑微向崇高的转变。藉着这个词，他宣告了虔敬的奥迹，因为当他说\InnerQuote{被高举在天主的右边}时，他明确揭示了奥迹中不可言说的救恩计划：那创造万有的天主的右手，就是万有藉以被造、没有祂万有皆不存的主，藉着结合，将与之合一的人性提升到自己的崇高之中。}\par

\Emph{\Person{圣安斐洛西}，\Place{依科尼雍}主教，的证言。}\par

出自他关于\Quote{我的圣父比我大}的讲道：—\par

\Quote{从此区分本性：天主的本性和人的本性。因为祂不是因离开天主而成为人，也不是因从人增长和进展而成为天主。}\par

出自他关于\Quote{圣子什么也不能做}的讲道：—\par

\Quote{因为在复活后，主显示了两者——既显示身体不是这种本性，也显示身体复活，因为记得历史。在苦难和复活之后，门徒们聚集在一起，\InnerQuote{当门关闭时，主站在他们中间}。在苦难之前，祂从未这样做过。那么，基督难道不能在很久以前就这样做吗？因为\InnerQuote{天主凡事都能}。但祂在苦难之前没有这样做，免得你以为降生成人是不真实或外表，并以为基督的肉身是属神的，或从天上降下，是与我们肉身不同的实体。一些人发明了所有这些理论，以为由此尊崇主，却忘记他们藉着感恩反而亵渎了自己，并控告真理说谎：因为我不说这谎言完全荒谬。因为如果祂取了另一个身体，那与我的需要救恩的身体何干？如果祂从天上带来肉身，那与我从地上出生的肉身何干？}\par

同一位作者，同一著作：—\par

\Quote{因此，主不是在苦难之前，而是在苦难之后，当门都关着的时候，站在门徒中间，为叫你知道，你的自然身体在播种之后被\InnerQuote{复活为属神身体}（\OriginalTerm{属神身体}{\GreekText{σῶμα} \GreekText{πνευματικόν}}），并且使你不要以为被复活的身体是另一个身体。当多默在复活后怀疑时，主向他显示钉痕，向他显示枪伤。难道主在复活后没有能力医治自己吗？即使在复活前祂已医治了所有人。但通过显示钉痕，祂表明这正是这个身体；而通过门关着时进来，祂表明它并不具有同样性质；这同一个身体是为了完成降生成人的工作，复活那已成尸体的身体，但却是一个改变了的身体，使它不再陷入腐朽，也不再受死亡的辖制。}\par

\Emph{亚历山大里亚主教真福德奥斐洛的证词。}\par

从《驳奥力振》一书中：——\par

\Quote{祂所取我们的相似性并没有转变为天主性的本质，祂的天主性也没有转变为我们的相似性。因为祂保持祂从起初就是的天主，并且这样保持着，在祂内保存我们的\OriginalTerm{实质}{\GreekText{ὑπόστασις}}。}\par

同上，出自同篇论文：——\par

\Quote{但是你持续不断地用你的亵渎之言攻击天主子，并使用这些话：\InnerQuote{如同圣子和圣父是一体，同样圣子所取的灵魂与圣子本身也是一体。}你不知道圣子和圣父因着同一\OriginalTerm{本体}{\GreekText{οὐσία}}和同一\OriginalTerm{天主性}{\GreekText{Θεότης}}而是一体；但灵魂与圣子各自具有不同的本体和不同的本性。因为如果圣子的灵魂与圣子本身是一体，正如圣父与圣子是一体，那么圣父与灵魂也会是一体，而圣子的灵魂将有一天说：\BibleQuote{看见了我，就是看见了父。}但这并非如此；天主不许。因为圣子和圣父是一体，因为在本性上没有区别，但灵魂与圣子在自然和本体上都是有区别的，因为那在本质上与我们同一本体的灵魂是由祂造的。因为如果灵魂与圣子是一体，如同圣父与圣子是一体，正如奥力振所认为的，那么灵魂将与圣子同是\BibleQuote{天主荣耀的光辉和祂本体的真像}。但这是不可能的；不可能圣子和灵魂像祂与圣父那样是一体。而当奥力振再次攻击自己时，他会做什么？因为他写道：\InnerQuote{灵魂痛苦和\InnerQuote{极其忧伤}绝不可能是\InnerQuote{首生者}（\OriginalTerm{首生者}{\GreekText{πρωτότοκος} \GreekText{πάσης} \GreekText{κτίσεως}}）。}因为天主圣言作为比灵魂更强有力的，圣子自己说：\BibleQuote{我有权舍掉，也有权再取回来。}如果圣子比自己的灵魂更强有力，正如所公认的，那么祂的灵魂怎能与天主同等，且具有天主的形状呢？因为我们说：\BibleQuote{祂贬抑自己，取了奴仆的形状。}在祂不敬的狂悖中，奥力振超越所有其他异端者，正如我们所表明的，因为如果圣言具有天主的形状且与天主同等，而他又如此胆敢写说救主的灵魂具有天主的形状且与天主同等，那么同等者如何能更伟大，当天性低下的证明那超越者的优越性呢？}\par

\Emph{君士坦丁堡主教圣金口若望的证词。}\par

来自《在大堂中的讲道》：——\par

\Quote{你的主将人高举到天上，而你甚至不肯给他分享市场的一角。但为什么我说\InnerQuote{到天上}呢？祂将人安坐在王座上。你却将他从城中驱逐出去。}\par

同上，论\BibleBook{圣咏}{第四十二篇}的开端：——\par

\Quote{直到今日，保禄没有停止说：\BibleQuote{我们作基督的使者，好像是天主借着我们劝你们；我们替基督求你们：与天主和好吧！}祂不仅站在这地上，而且取了你们本性的初果，坐在\BibleQuote{远超一切执政的、掌权的、有能的、主治的，和一切有名的，不但是今世的，连来世的也超过了。}有什么能与此荣耀相比？我们族类的初果，曾如此冒犯并如此蒙羞，竟坐得如此高，享受如此大的荣耀。}\par

同上，论语言的分裂：——\par

\Quote{请思量，看到我们的本性乘驾在革鲁宾之上，天上所有军旅都围绕它，这是何等的事。也要思想保禄的智慧，他搜寻了多少词语以陈明基督对人的爱；因为他不仅仅说恩宠，也不仅仅说丰盛，而是说\BibleQuote{祂在天主的仁慈中对我们所施的恩宠的极丰富的丰盛。}}\par

同上，出自他的《教理论》，主题是：基督在谦卑中所说的话和所做的事，并非出于软弱，而是出于不同分施的差异：——\par

\Quote{复活后，当祂看到祂的门徒不信时，祂不避开向他显示伤口和钉痕，并让他将手放在伤痕上，说：\BibleQuote{你们看看，摸摸，因为鬼神没有肉和骨。}祂不从一开始就取了成年的身体，而甘愿受孕、出生、吃奶，并在地上生活如此之久，其原因正是为借着时间的悠长和所有其他情况，为这事提供保证。}\par

同上，驳斥那些声称魔鬼统治人事的人：——\par

\Quote{没有比人更无价值的，也没有比人更成为珍贵的。他是理性受造物中的最后部分，但脚成了头，并通过初果被带到君王宝座上。正如一位高尚宽厚的人，看见一个遭遇船难的可怜人，从波浪中只救出赤身，便张开双臂欢迎他，给他穿上光彩夺目的长袍，并提升他到最高的尊荣；同样，天主也如此对待我们的本性。人失去了他所拥有的一切：他的自由、他与天主的交往、他在乐园中的居所、他的无痛苦生活，于是他像一个从船难中完全赤裸的人出来，但天主接纳了他，立即给他穿上衣服，牵着他的手，一步一步地引导他，将他带到天堂。}\par

同上，出自同篇著作：——\par

\Quote{但是天主使所得大于所失，并将我们的本性高举至王座。因此保禄高呼说：\Quote{祂使我们一同复活，并使我们在天上坐在祂的右边。}}\par

同一位作者在其反对犹太人的第三篇讲道中：—\par

\Quote{祂敞开了天；化敌为友；将他们引入天上；使我们的本性坐在宝座的右边；赐给我们无数其他的恩惠。}\par

同一位作者在其论升天的讲道中：—\par

\Quote{祂将我们的本性提高到了这样的距离和高度。看它曾低落在何处，又攀升到了何处。不可能降到比人降落的更低之处；也不可能升到比祂所高举的更高之处。}\par

同一位作者在其《厄弗所书》释义中：—\par

\Quote{按照祂在自己内所预定的善意，也就是祂所热切渴望的，祂仿佛在努力告诉我们这奥秘。这奥秘是什么？就是祂愿意将人高举于高天；正如实际所成就的那样。}\par

同一位作者出自同一释义：—\par

\Quote{\Quote{我们主耶稣基督的天主}指的是祂，而非天主圣言。}\par

同一位作者出自同一释义：—\par

\Quote{\Quote{当我们死在过犯中的时候，祂使我们与基督一同活过来。}基督再次立于中间，这作为是奇妙的。如果初果活着，我们也活着。祂使祂和我们都活过来。你看见这一切都是按照肉身说的吗？}\par

同一位作者出自《若望福音》释义：—\par

\Quote{为什么祂加上\Quote{并居住在我们中间}？仿佛祂说：不要从\Quote{成了}这句话中想象任何荒谬之事。因为我并未提到那不变本性中的任何变化，而是说到帐幕和居住。那支搭帐幕的与帐幕不是同一的，而是一物在另一物中支搭帐幕；否则就没有帐幕了。没有任何物居住在自己里面。我所说的是本质的区别。因为藉着结合与连接，天主圣言和肉身在不混淆、不毁灭本质的情况下成为一体，而是藉着不可言喻、无法形容的结合。}\par

同一位作者出自《玛窦福音》释义：—\par

\Quote{正如一个人站在彼此分离的二者之间，伸出双手将它们连接起来，同样地，祂也将旧与新的、天主性与人性、祂的与我们的连接起来。}\par

同一位作者出自《基督升天论》：—\par

\Quote{因为当两个对手准备战斗时，有第三者插入其间，立即制止争斗，结束他们的敌意，基督也是这样。作为天主，祂曾发怒，但我们轻视祂的义怒，背弃了我们慈爱的主的面容。于是基督投身其间，使两种本性和好相爱，并亲自承担了父加于我们身上的惩罚的重担。}\par

同一位作者出自同一著作：—\par

\Quote{看，祂将我们本性的初果献给父，父自己也悦纳这礼物，既因献祭者的崇高尊位，也因祭品的无瑕。祂用自己的双手接住它，在自己的宝座上安置了一把椅子；不仅如此，更让它坐在自己的右边。那么，让我们认清，那听到\Quote{坐在我的右边}的是谁，以及那被天主说\Quote{你是尘土，仍要归于尘土}的是何种本性。}\par

同一位作者稍后之处：—\par

\Quote{我不知用何论证，说何言语；那原本腐朽、毫无价值、被视为万物中最卑下的本性，竟战胜了一切，征服了世界。今日它被尊为超越万有；今日它获得了自古以来众天使所渴望的；今日总领天使得以观看那长久以来所期盼的，他们注视我们的本性，在王的宝座上，在祂不朽的荣耀中闪耀。}\par

\Emph{圣弗拉维亚诺，安提约基雅主教，的见证。}\par

出自《路加福音》释义：—\par

\Quote{主在我们众人身上写下祂圣德的明像，并以各种方式向我们本性指示救恩之路。祂赐给我们许多明确的证据，既有关于祂肉体降临的，也有关于祂藉肉体行事的天主性的。因为祂愿意使我们确信祂的两种本性。}\par

同一位作者论《主显节》：—\par

\Quote{\Quote{谁能述说上主的大能，或彰显祂的一切赞美？}谁能用言语表达祂对我们恩惠的伟大？人的本性结合于天主性，而两种本性各自独立。}\par

\Emph{济利禄，耶路撒冷主教，的见证。}\par

出自其第四篇《论十端》的教理讲道。\par

论童贞诞生：—\par

\Quote{你要相信，天主的独生子，因我们的罪，从天降下到地上，取了与我们相同情感的人性，并由圣童贞和圣神所生。这降生成人不是在表象和虚假中，而是在真实中成就的。祂并非仅仅经过童贞如经过一条通道，而是真实地从她取了肉身。祂真地像我们一样吃，并从童贞真实地吃奶。因为如果降生成人是虚假的，那么我们的救恩就是幻觉。基督是双重的——可见的人，不可见的天主。祂作为人真实地吃，与我们一样，因为祂所穿上的肉身与我们有着相同的情感；祂作为天主用五个饼喂养了五千人。祂作为人真实地死了；祂作为天主在第四天复活了死人。祂作为人在船中睡觉；祂作为天主在水面上行走。}\par

\Emph{安提约古，托勒迈斯主教，的见证}：—\par

\Quote{不要混淆本性，你就会生动地理解降生成人。}\par

\Emph{圣依拉略，主教与精修者，的见证}，在其第九卷\Quote{论信德}中：\par

\Quote{不认识耶稣基督是真天主又是真人的人，实际上不认识自己的生命，因为我们若否认基督耶稣或天主圣神，或我们自身肉身的肉体，便遭受同样的危险。\BibleQuote{凡在人前承认我的，我在我天上的父前也必承认他；凡在人前否认我的，我在我天上的父前也必否认他。} 这些话是成了血肉的圣言所说的；这些话是基督耶稣——人而天主、荣耀之主——所教导的，祂成了中保，以祂在天主与人之间中保的奥迹本身，为教会的救恩而服务。为了这个目的，祂由联合的两性成为一体，借着两种性体，祂是同一而相同的，但在两者中祂都不缺欠，免得祂因生为人而不再是天主，或因为仍然是天主而不是人。因此，真正的信仰在人间的福乐就在于宣讲天主和人，宣认圣言和血肉，承认天主也是人，并且不忽略血肉也是圣言。}\par

同一位作者，同书：—\par

\Quote{所以，独生天主由童贞女降生为人，在时期一满，祂亲自预定要成就人迈向天主的进步，祂在福音的全部言语中遵守了这种秩序，以便教导人相信祂是天主之子，并提醒我们宣讲祂是人子。作为人，祂总是按照人的方式说话行事，但从不以同样的言语方式谈论两者，除非有意同时指明天主和人。但是，异端者由此得到借口，将简单无知的人诱入陷阱：他们将我们的主按照人性所说的话，错误地断言是在祂天主性软弱时所说的，并且由于同一说话者说了所有的话，他们便主张祂所说的一切都是关于祂自己的。我们甚至也不否认祂现存的一切言语都出自祂的本性。但既然承认基督既是人又是天主；承认祂作为人时并非那时才成为天主；承认祂作为人时也是天主；承认在主的位格中取人性之后，圣言既是人又是天主，那么必然的结果是：祂言语的奥迹与祂诞生的奥迹是同一而相同的。每当你在祂身上根据场合需要区分人性和天主性时，也要努力将天主的言语与人的言语分开。每当你承认天主和人时，就要辨别天主的话和人的话。当谈及天主和人，以及完全的人与完全的天主时，要谨慎考虑场合。如果说了什么来指明特定场合所适合的事，就将言语应用于该场合。必须区分：在人性存在之前的天主，人而天主，以及人性与天主性结合后完全的人与完全的天主。因此，要小心不要在言语和行为中混淆降生成人的奥迹。因为根据本性种类的性质，在言语方式上必然存在区别：在人性诞生之前，按照奥迹仍趋近死亡时，以及再次成为永恒时。\InnerQuote{因为祂在诞生、受难和死亡中，按照我们的本性行事，但祂仍凭自己本性的能力成就了这一切。}}\par

同一位作者，同书：—\par

\Quote{那么你看到，这样承认天主和人，以致死亡归属于人，肉身的复活归属于天主；因为请思考天主的本性和复活的能力，并在死亡中认识关于人的\OriginalTerm{救恩经济}{œconomy}。既然死亡和复活都在各自的本性中实现了，我请求你记住：一个基督耶稣，祂出自两者。我简短地向你说明了这些要点，为使我们记得，在我们的主耶稣基督内有两个本性，\BibleQuote{祂虽具有天主的形体，却取了奴仆的形体。}}\par

\Emph{至圣主教\Person{奥斯定}的证言。}\par

出自他致\Person{沃鲁西亚诺}的书信。第三书：\par

\Quote{但现在祂显现为天主与人之间的中保，在祂位格的统一中结合了两性，将惯常的与不惯常的、不惯常的与惯常的结合起来。}\par

同一位作者，出自他\WorkTitle{若望福音释义}：—\par

\Quote{那么，异端者啊？既然基督也是人，祂就作为人说话；而你却诽谤天主？祂在自己内提升人性，而你竟敢贬低祂的天主性。}\par

同一位作者，出自他\WorkTitle{论信仰释义}一书：—\par

\Quote{相信是我们的本分，知道是祂的事；因此，让天主圣言自己在接受了属于人的一切之后，成为人；让人在祂的取纳和接受了属于天主的一切之后，不再是别的，而是天主。不可因为据说祂\OriginalTerm{降生成人}{incarnate}并混合，就以为祂的\OriginalTerm{本体}{substance}减少了。天主知道祂混合时并无自然的腐败，祂是真实地混合。祂也知道祂如此接纳在自己内，以至于没有给自己增添什么，正如祂知道祂倾注了整个自己，以至于没有遭受减少。那么，不要按照我们软弱的智力，并根据经验和感觉的教导形成猜测，以为天主和人像受造物和同等物混合那样混合，并且从这种圣言和肉身的混淆中仿佛形成了一个身体。天主禁止我们如此相信，免得我们以为按照混淆之物的方式，两个本性被带进一个\OriginalTerm{位格}{hypostasis}。因为这种说法意味着对两部分的毁灭；但基督自己，包含而不被包含，察验我们而自身超越察验，充满而不被充满，处处同时是完整的自己，并渗透宇宙，因祂倾注自己的力量，因怜悯感动，与人的本性混合，而人的本性并未与神的混合。}\par

\Emph{加巴拉主教\Person{色味利亚诺}的证言。}\par

出自\WorkTitle{基督的诞生}：—\par

\Quote{奥迹，真是属天而又在人间——是可见而又不被看见的奥迹，因为基督降生后正是如此：属天而又在人间；持有而又不被拥有；可见而又不可见；就天主性的本质而言属于天上，就人性的本质而言属于人间；在肉身上可见，在精神上不可见；就身体而言可以被持有，就圣言而言不可被把握。}\par

\Emph{阿提库斯，君士坦丁堡主教，的见证。}\par

摘自他致欧普西基乌斯的信：\par

\Quote{那么，至智者应当如何行动？通过所取肉身的媒介，以及天主圣言与由玛利亚所生之人的结合，祂由两种本性而成，以致基督由两者合而为一，既在天主性中构成，又在其不受苦难的本性之尊严中永存，但在肉身中接近死亡，同时向肉身显示同类的本性如何藉着死亡而蔑视死亡，并藉着祂的死亡确认了新盟约的正义。}\par

\Emph{济利禄，亚历山大主教，的见证。}\par

摘自他致聂斯多略的信：\par

\Quote{在真实合一中结合在一起的本性是各别的，并且出自两者有一位天主和圣子，但本性的差异并未因合一而消除。}\par

同一位作者摘自他反对东方派的信：\par

\Quote{有两性之结合，因此我们承认一位基督、一位圣子、一位主。依照这无混淆合一的领悟，我们承认童贞玛利亚为天主之母，因为天主圣言成了血肉并成了人，并且从怀孕起就将从她所取的殿宇与祂自己结合。}\par

同一位作者：\par

\Quote{只有一位主耶稣基督，即使认出我们所说不可言喻的合一由之而成的两性之区别。}\par

同一位作者：\par

\Quote{因此，如我所说，在赞美降生成人的方式时，我们看到两性在不可分离的合一中无混淆、无分割地走到一起，因为肉是肉，而非任何天主性，尽管它成了天主的肉；同样，圣言是天主，而非肉，尽管祂按照\OriginalTerm{经济}{œconomy}使肉成为自己的。}\par

同一位作者摘自他对《希伯来书》的诠释：\par

\Quote{因为虽然在合一中相聚的各性彼此视为不同且不平等——我是指肉和天主——然而由两者所成的圣子是一位。}\par

同一位作者摘自他对同一书信的诠释：\par

\Quote{然而，即使独生天主圣言被说成与肉\OriginalTerm{位格}{hypostasis}地合一，我们也否认各性之间有任何混淆，并宣称每一性保持其本来面目。}\par

同一位作者摘自他的注释集：\par

\Quote{圣父的圣言，由童贞女所生，被称为人，尽管祂本性是天主，因分享了与我们一样的血肉，因为祂就这样在地上被人看见，没有放弃自己的本性，却按我们人性的理性取了圆满的人性。}\par

同一位作者论降生成人（评注第13章）：\par

\Quote{那么，在降生成人之前，有一位真天主，而在人性中，祂仍然是祂曾是的、现在是、将来也是的那一位；主耶稣基督不可被分开为单独的人和单独的天主，而是认出各性的区别并保持它们彼此无混淆，我们宣称有同一的基督耶稣。}\par

同一位作者在其他注释之后：\par

\Quote{有一种清晰的认识，即一物寓于另一物中，即天主性寓于人性中，没有任何混合、混淆或变成不是它的东西。因为说寓于另一物的，并不变成它所寓之物，而是被视为一物在另一物中。但在圣言的本性和人性的本性中，区别只给我们指出本性的差异，因为从两者认出是一位基督。因此祂说圣言\InnerQuote{在我们中间搭帐棚}，谨慎地保持无混淆，因为祂承认一位独生圣子，成了血肉并成了人。}\par

现在，我亲爱的先生，你已听到了世界的大光；你已看到了他们教导的光芒，并得到了精确的指导，不仅是关于降生之后，而且是关于成就救恩的苦难、复活和升天之后，他们如何表明天主性和人性的合一是不混淆的。\par

\Emph{厄然。}——我原以为他们在合一之后没有分别各性，但我发现了无数的区别。\par

\Emph{正统。}——对那些信仰的高贵捍卫者，你竟敢动舌，实在是疯狂和鲁莽。但我将为你引述阿波利纳里的话，好让你知道他也断言合一是不混淆的。现在请听他的话。\par

\Emph{阿波利纳里的见证。}\par

摘自他的概要：\par

\Quote{在天主方面和身体方面之间有合一。一方面是可敬的创造主，祂是永恒的智慧和能力；这些属于天主性。另一方面是玛利亚的儿子，在末期出生，敬拜天主，在智慧上长进，在能力上坚强；这些属于身体。为罪和咒诅的苦难已经来到，不会消逝，也不会变成无实体。}\par

再往后一点：\par

\Quote{人就无理性的身体而言与无理性的动物是\OriginalTerm{同体}{consubstantial}；就他们有理性的程度而言，他们属于另一种实体。同样，天主按肉身与人是同体的，就祂是圣言又是人的程度而言，祂属于另一种实体。}\par

在另一处他说：\par

\Quote{混杂之物的性质相互混合而不消灭。因此，有些成分得以从混合物中分离出来，如酒从水中分离；且不与身体混合，亦非如身体与身体之混合，而是混合也保存着未混合者，以至于根据各次需要，神性之能力或独立行动或结合行动，正如主禁食之时，神性因结合而相应超出需要之上，饥渴受到抑制；但当它不再以超于需要之状态对抗渴求时，饥渴便发生，以败坏魔鬼。但如果身体的混合未受改变，神性的混合岂不更如此？}\par

在另一处，他说：——\par

\Quote{若与铁之混合使铁本身成为火而不改变其性质，同样，天主与身体的结合也不意味着身体的改变，即使身体将其神性的能力延伸至可达到的范围。}\par

对此他立即补充说：——\par

\Quote{若一人既有灵魂又有身体，两者保持合一，那么，基督既有神性和身体，岂不更能使两者稳固而不混淆？}\par

又稍后：——\par

\Quote{因为人性按其所能分享神性的能力，但大小之差分明如最小与最大。人是天主的仆人，但天主不是人的仆人，甚至不是祂自己的仆人。人是天主的受造物，但天主不是人的受造物，甚至不是祂自己的受造物。}\par

又：——\par

\Quote{若有人将\BibleQuote{子看见圣父所作的事，祂也照样作}\BibleRef{若 5:19}这句话视为就神性而非就肉体而言，则成为肉身的祂便与未成为肉身的圣父有所区分，从而区分了两种神性的能力。但并无区分。因此，祂不是就神性而言说的。}\par

他又说：——\par

\Quote{正如人并非无理性的受造物，因为理性与无理性接触；同样，救主也非受造物，因为受造物与不受造的天主接触。}\par

对此他也补充说：——\par

\Quote{那不可见的与可见的身体结合而因此被看见，仍保持不可见，且保持不混杂，因为它不被身体所限定；身体在其本有尺度内，按照被赋予生命的方式接受与天主的结合，而被赋予生命的并不是赋予生命者。}\par

又稍后他说：——\par

\Quote{若灵魂与身体的混合，虽从起初便结合，并未使灵魂因身体而可见，也未将其改变为身体的其他属性，以致能被割裂或减少；那么，身体以灵魂为生，且人的身体仍保持其本性，何况天主（祂与身体并非同性）与身体结合而不受改变，则在基督身上，混合同样不改变身体，使之不再是身体。}\par

再往后他又说：——\par

\Quote{凡宣认灵魂和身体由圣经构成一体者，在断言圣言与身体的这种结合是一种改变时，便自相矛盾——这样的改变即使在灵魂的情况下也未曾见到。}\par

再听清他的呼喊：——\par

\Quote{若否认主的肉体常住者为不敬，那么完全拒绝接受祂降生成人者更为不敬。}\par

并且在他关于降生成人的小册子中写道：——\par

\Quote{\InnerQuote{坐在我的右边}这句话是对人说的，因为作为天主的圣言从地上升天之后，祂不常坐在荣耀的宝座上；但这话是对那如今作为人被高举至天上荣耀者说的，如宗徒们说：\BibleQuote{\Person{达味}并没有升到天上，但他自己说：主对我主说：你坐在我右边}\BibleRef{咏 110:1}；\BibleRef{宗 2:34-35}。这命令是人的，为坐立设定了起始；但与天主同坐则是神性的尊荣，\BibleQuote{成千上万事奉祂，万万千千侍立在祂面前}\BibleRef{达 7:10}。}\par

又稍后：——\par

\Quote{祂并非作为天主将仇敌置于脚下，而是作为人；但被看见的天主和人乃是同一。\Person{保禄}也教训我们，\BibleQuote{直到我使你的仇敌作你的脚凳}\BibleRef{咏 110:1}是对人说的，当然按照祂的神性，将胜利归于自己，\BibleQuote{按祂能征服万有归服自己的大能}\BibleRef{斐 3:21}。看，神性和人性在一位格内不可分离地存在着。}\par

又：——\par

\Quote{\BibleQuote{圣父，请祢在祢自己内光荣我，赐给我在世界存在之前，我在祢面前所拥有的光荣}\BibleRef{若 17:5}。祂作为人使用\InnerQuote{光荣}一词，但作为天主启示出祂在万世之前已拥有此光荣。}\par

又：——\par

\Quote{但我们不要因以为敬拜天主子是羞辱——即使在祂人性的样式中——而感到羞愧；要如同尊敬一位穿着破衣却带有王权荣耀的君王，尤其看见祂所披的外袍本身也被光荣，正如天主的身体和世界救主的身体，它是永生之种、神圣作为的工具、一切邪恶的毁灭者、死亡的击杀者和复活之元首；尽管它从人获得本性，却从天主获得生命，从天上获得能力与神圣之德。}\par

又：——\par

\Quote{因此我们敬拜身体如同圣言；我们领受身体如同领受圣神。}\par

如今已清楚向你们表明，那位最先引入本性混合的作者，公开使用了二者间的区分论据；他将身体称为外衣、受造物和工具；他甚至称之为仆人，这是我们中无人敢做的。他还说身体被配得坐在右边，并使用了其他许多被你们虚妄异端所拒绝的表述。\par

\Person{埃兰}：——但是，为何那位最先引入混合的人又将如此巨大的区分插入他的论据中呢？\par

\Person{正统}：——真理的力量迫使即使激烈反抗她的人也同意她所说的；但若你愿意，现在让我们开始讨论主的无受苦性。\par

\Person{埃兰}：你知道乐师通常让琴弦休息，通过转动弦轴放松它们。如果完全没有理性和灵魂的事物都需要一些娱乐，那么我们这些拥有理性和灵魂的人，如果按我们的能力分配劳作，就不会有任何荒谬之处。那么我们就把它推迟到明天吧。\par

\Person{正统}：神圣的\Person{达味}嘱咐我们日夜留意神圣的圣言；但就照你说的，我们把剩余主题的探究留到明天吧。\par

\SourceRecord{Translated by Blomfield Jackson.；Nicene and Post-Nicene Fathers, Second Series；Vol. 3.；Edited by Philip Schaff and Henry Wace.；Buffalo, NY: Christian Literature Publishing Co.,；1892.；<http://www.newadvent.org/fathers/27032.htm>.}

% structure: p0001 source=h1 target=section reason=page-title

\section{第三篇对话}

\Emph{\Person{正统派}与\Person{杂拌者}。}\par

\Emph{正}——在我们先前的讨论中，我们已经证明了圣言是\OriginalTerm{不变的}{immutable}，祂\OriginalTerm{降生成人}{incarnate}并非因转变为血肉，而是因取了完美人性。圣经和教会的导师、世界之光已清楚教导我们：在结合之后，祂仍然保持原样，未混合、\OriginalTerm{不能受苦的}{impassible}、未改变、不受限制；并且祂保全了祂所取的本性，毫无损缺。因此，接下来我们的话题是祂的\OriginalTerm{苦难}{passion}，这将是一个极有益处的话题，因为从那里为我们带来了救恩之水。\par

\Emph{杂}——我也认为这番话是有益的。但我不同意我们先前的方法，而是提议由我来提问。\par

\Emph{正}——我会回答，不反对方法的改变。拥有真理者，不仅在提问时，在被问时也因真理的威力而得到支持。请问你想要的吧。\par

\Emph{杂}——按照你的看法，谁受了苦难？\par

\Emph{正}——我们的主耶稣基督。\par

\Emph{杂}——那么一个人给了我们救恩。\par

\Emph{正}——不；难道我们曾承认我们的主耶稣基督仅仅是个人吗？\par

\Emph{杂}——现在请你定义你认为基督是什么。\par

\Emph{正}——永生天主的降生成人的圣子。\par

\Emph{杂}——天主子是天主吗？\par

\Emph{正}——是天主，与生祂的天主有同一性体。\par

\Emph{杂}——那么天主受了苦难。\par

\Emph{正}——如果祂没有身体就被钉在十字架上，请把苦难归于神性；但如果祂取了肉体而成为人，为什么你要把能受苦的免除苦难，反而使不能受苦的受苦难呢？\par

\Emph{杂}——但祂取肉体的原因，正是为了使不能受苦的借着能受苦的而受苦难。\par

\Emph{正}——你说不能受苦，却把苦难加给祂。\par

\Emph{杂}——我说过祂取了肉体是为了受苦。\par

\Emph{正}——如果祂本有能受苦的本性，祂本可以无需肉体而受苦；那么肉体就成为多余的了。\par

\Emph{杂}——神性是不死的，肉体的本性是可死的，所以不死者与可死者结合，为的是通过它尝到死亡。\par

\Emph{正}——那本然不死者，即使与可死者结合，也不会经历死亡；这一点很容易看出。\par

\Emph{杂}——请证明它，并解决这个难题。\par

\Emph{正}——你认为人的灵魂是不死的，还是可死的？\par

\Emph{杂}——不死的。\par

\Emph{正}——身体是不死的还是可死的？\par

\Emph{杂}——毫无疑问是可死的。\par

\Emph{正}——我们说人是由这些本性组成的吗？\par

\Emph{杂}——是的。\par

\Emph{正}——那么不死者与可死者结合了？\par

\Emph{杂}——对。\par

\Emph{正}——但是当连接或结合结束时，可死者服从死亡的律法，而灵魂仍保持不死，尽管罪引入了死亡；或者你不认为死亡是一种刑罚吗？\par

\Emph{杂}——圣经是这样教导的。因为我们得知，当天主禁止亚当吃知善恶树的果子时，祂附加了说：\BibleQuote{你吃的那一日，你必定要死。}\par

\Emph{正}——那么死亡是对犯罪者的惩罚？\par

\Emph{杂}——同意。\par

\Emph{正}——那么为什么，当灵魂和身体都一同犯了罪时，只有身体单独承受死亡的惩罚？\par

\Emph{杂}——是身体用邪恶的眼睛看那棵树，伸出手，摘了禁果。是嘴用牙齿咬，嚼碎，然后食道把它送进肚子，肚子消化它，把它交给肝；肝把所接受的变成血，再传给空静脉，静脉传给邻近部位，它们再传给其余部分；这样禁果的偷窃就遍布全身。所以身体单独承受罪的惩罚是很恰当的。\par

\Emph{正}——你给我们作了一番关于食物本性、经过的所有部分以及在同化于身体之前所受变化的生理学论述。但是有一件事你不愿意注意，就是身体离开灵魂就不能进行你提到的任何一个过程。当身体被夺去灵魂这个伴侣时，它就变得无气息、无声、无动作；眼睛看不见对或错；没有声音传到耳朵；手不能动；脚不能走；身体就像没有音乐的乐器。那么你怎么能说只有身体犯了罪，而身体没有灵魂甚至连一口气都不能吸呢？\par

\Emph{杂}——身体确实从灵魂接受生命，它给灵魂提供有罪的占有的惩罚。\par

\Emph{正}——怎么回事，以什么方式？\par

\Emph{杂}——通过眼睛，它使灵魂错看；通过耳朵，使灵魂听无益的声音；通过舌头，说有害的话；并通过所有其他部分行恶。\par

\Emph{正}——那么我想我们可以说：聋子有福了；失去视力、被剥夺其他官能的人有福了，因为如此残废之人的灵魂与身体的邪恶无分无关。但是，最明智的先生，为什么你只提到那些应受谴责的身体功能，却没有说值得称赞的呢？可以用爱的眼睛、仁慈的眼睛去看；可以擦去忏悔的眼泪；可以听天主的神谕；可以向穷人侧耳；可以用舌头赞美造物主；可以给邻舍好的教导；可以用行善的手；总之，可以用身体的各部分完全获得良善。\par

\Emph{杂}——这一切都是真的。\par

\Emph{正}——因此，对律法的遵守和违反是灵魂和身体共有的。\par

\Emph{杂}——是的。\par

\Emph{正}——在我看来，灵魂在两者中都起主导作用，因为它在身体行动之前先运用理性。\par

\Emph{杂}——你这样说是什么意思？\par

\Emph{正}——首先，心智仿佛画出美德或恶的草图，然后用适当的材料和颜色赋予其中之一以形式，并通过身体的各部分作为工具来使用。\par

\Emph{杂}——似乎如此。\par

\Emph{正教派}——如果灵魂与身体一同犯罪，不，甚至是在罪中领先，因为约束和指导动物部分的责任交托给了灵魂，那么为何灵魂既然分享了罪，却不也分享惩罚呢？\par

\Emph{厄蓝}——但不朽的灵魂怎么可能分享死亡呢？\par

\Emph{正教派}——然而，既然分享了过犯，理当分享惩罚。\par

\Emph{厄蓝}——是的，理当如此。\par

\Emph{正教派}——但事实并非如此。\par

\Emph{厄蓝}——当然不是。\par

\Emph{正教派}——至少在来世，它将与身体一同被投入\OriginalTerm{地狱}{Gehenna}。\par

\Emph{厄蓝}——所以祂说：\BibleQuote{你们不要害怕那杀害肉身而不能杀害灵魂的；但更要害怕那能把灵魂和肉身都投入地狱中的。}\BibleRef{玛 10:28}\par

\Emph{正教派}——因此，在今世，灵魂作为不朽的，逃脱了死亡；在来世，它将受惩罚，不是经历死亡，而是在生命中以受折磨来承受惩罚。\par

\Emph{厄蓝}——神圣的圣经是这样说的。\par

\Emph{正教派}——那么，不朽的本性经历死亡是不可能的。\par

\Emph{厄蓝}——看来是这样。\par

\Emph{正教派}——那么，你怎么说天主圣言尝了死亡呢？因为如果被造的不朽者显然不能变成可朽的，那么那无始无终、永远不朽、创造可朽和不朽本性的创造者，又怎么可能分受死亡呢？\par

\Emph{厄蓝}——我们也知道祂的本性是不朽的，但我们是说祂在肉身内分受了死亡。\par

\Emph{正教派}——但我们已清楚证明，那本性不朽的，无论如何不可能分受死亡；因为就连与身体一同受造、结合、并分担罪的灵魂，也不与身体同受死亡，仅仅因其本性不朽。但我们再从另一个观点来看这同一立场。\par

\Emph{厄蓝}——我们当然应该不遗余力地寻求真理。\par

\Emph{正教派}——那么让我们这样来探讨这个问题。我们是否主张，在德性与恶习上，有些是教师，有些是学生？\par

\Emph{厄蓝}——是的。\par

\Emph{正教派}——我们是否说，德性的教师应得更大的赏报？\par

\Emph{厄蓝}——当然。\par

\Emph{正教派}——同样，恶习的教师应得双倍甚至三倍的惩罚？\par

\Emph{厄蓝}——正确。\par

\Emph{正教派}——那么，我们要把魔鬼放在哪一类呢？是教师还是学生？\par

\Emph{厄蓝}——教师中的教师，因为他本身就是一切不义的父和教师。\par

\Emph{正教派}——那么，谁成了他的第一批门徒？\par

\Emph{厄蓝}——亚当和厄娃。\par

\Emph{正教派}——谁接受了死刑判决？\par

\Emph{厄蓝}——亚当和他所有的后裔。\par

\Emph{正教派}——那么门徒因学到的坏教训而受罚，而我们刚刚宣称应得双倍甚至三倍惩罚的教师，却逃脱了惩罚？\par

\Emph{厄蓝}——看来如此。\par

\Emph{正教派}——虽然事情是这样，我们却都承认并宣称审判者是公义的。\par

\Emph{厄蓝}——当然。\par

\Emph{正教派}——但既然是公义的，祂为什么没有向他追究恶教的账呢？\par

\Emph{厄蓝}——祂为他准备了地狱不灭的火，因为祂说：\BibleQuote{你们可咒骂的，离开我，到那给魔鬼和他的使者预备的永火里去吧！}\BibleRef{玛 25:41}而他之所以在此世没有与门徒一同死亡，是因为他具有不朽的本性。\par

\Emph{正教派}——那么，即使最大的罪人，如果具有不朽的本性，也不能遭受死亡。\par

\Emph{厄蓝}——同意。\par

\Emph{正教派}——那么，既然不义的发明者和教师本身，因其本性不朽而没有遭受死亡，你难道不害怕说那不朽与正义之源分受了死亡吗？\par

\Emph{厄蓝}——如果我们说过祂被动地承受了苦难，那么你的指控也许有点道理。但如果我们所宣讲的苦难是自愿的，死亡是甘愿的，那么你不但不该指控我们，反而应当赞美祂对人的无限仁爱。因为祂受苦是因为祂愿意受苦，祂分受死亡是因为祂愿意。\par

\Emph{正教派}——你似乎完全不了解神性，因为主天主所愿的，无一与祂的本性相悖；祂能做祂愿意的一切，而祂所愿的，都适合并符合祂自己的本性。\par

\Emph{厄蓝}——我们学到，在天主一切都是可能的。\par

\Emph{正教派}——你这样不明确地表达，甚至包含了属魔鬼的事，因为说\Quote{一切}，不仅包括了善，也包括了它的反面。\par

\Emph{厄蓝}——但高贵的约伯不是笼统地说：\BibleQuote{我知道你万事都能做，你的旨意不能拦阻}\BibleRef{约 42:2}吗？\par

\Emph{正教派}——如果你先读一下这位义人前面说的话，你就会从另一段话看出这段话的含义，因为他这样说：\BibleQuote{求你记念，你曾造我如抟泥，你还要使我归于尘土吗？你不是将我倒出像奶，使我凝结如同奶饼吗？你以皮和肉给我穿上，用骨与筋把我联络起来，你将生命和慈爱赐给我。}\BibleRef{约 10:9-12}\par

然后他又说：——\par

\BibleQuote{我心中存有这事，就知道你万事都能做，在你没有难成的事。}\BibleRef{约 42:2}那么，他不正是把这些事都归诸于那不朽的本性、宇宙的天主吗？\par

\Emph{厄蓝}——对于全能的天主，没有不可能的事。\par

\Emph{正教派}——那么，按照你的定义，全能的天主也能犯罪吗？\par

\Emph{厄蓝}——绝无可能。\par

\Emph{正教派}——为什么？\par

\Emph{厄蓝}——因为祂不愿意。\par

\Emph{正教派}——为什么祂不愿意？\par

\Emph{厄蓝}——因为罪与祂的本性不相符。\par

\Emph{正教派}——那么，有许多事祂不能做，因为过犯有许多种。\par

\Emph{厄蓝}——这类事，天主既不愿也不能做。\par

\Emph{正教派}——那些违反神性的事也一样。\par

\Emph{厄蓝}——什么事？\par

\Emph{正教派}——例如，我们学到天主是理智的和真正的光。\par

\Emph{厄蓝}——正确。\par

\Emph{正教派}——我们不能称祂为黑暗，也不能说祂愿意变成黑暗或能变成黑暗。\par

\Emph{厄蓝}——绝无可能。\par

\Emph{正教派}——同样，神圣的圣经称祂的本性是不可见的。\par

\Person{Eran}——是的。\par

\Person{Orth}——我们绝不能言语谓祂能为可见。\par

\Person{Eran}——不能，确实。\par

\Person{Orth}——也不能被理解。\par

\Person{Eran}——不能；因为祂并非如此。\par

\Person{Orth}——不能；因为祂不可理解，且全然不可接近。\par

\Person{Eran}——你说得对。\par

\Person{Orth}——而那自有者绝不能变为不存在。\par

\Person{Eran}——不可有此念！\par

\Person{Orth}——圣父也不能变为圣子。\par

\Person{Eran}——不可能。\par

\Person{Orth}——未生者也不能变为受生者。\par

\Person{Eran}——祂怎能如此。\par

\Person{Orth}——圣父绝不能变为圣子吗？\par

\Person{Eran}——绝不。\par

\Person{Orth}——圣神也绝不能变为圣子或圣父。\par

\Person{Eran}——这一切都不可能。\par

\Person{Orth}——我们还会发现许多其他同类的事，同样不可能，因为永恒者不会成为时间的，非受造者不会成为受造和被造的，无限者不会成为有限的，等等。\par

\Person{Eran}——这些没有一样是可能的。\par

\Person{Orth}——所以我们发现许多事对全能的天主是不可能的。\par

\Person{Eran}——真的。\par

\Person{Orth}——但在这些方面任何一项无能，都不是软弱的证明，而是无限大能的证明；而能有能力则肯定是非大能而是无能的证明。\par

\Person{Eran}——你为何这样说？\par

\Person{Orth}——因为这些事每一样都宣告天主不变且不更改的特性。因为善不能变为恶，显明了善的无限；那公义者绝不能变为不义，那真实者绝不能变为说谎者，表明真理和公义中的稳定与力量。因此，真光绝不能变为黑暗；自有者绝不能变为不存在，因为存在是永久的，光本性不变。所以，在考察所有其他事例后，你会发现\Quote{不能}是最高大能的宣告。这类事在天主是不可能的，神圣的宗徒也已感知并设定，因为他在致希伯来人书中说：\BibleQuote{藉这两件不可更改的事，天主不可能说谎，我们才得到强有力的安慰。}他表明这无能不是软弱，而是真正的大能，因为他声称祂是如此真实，以致在祂内甚至不可能有任何谎言。所以，真理的能力通过它的无能得以表明。而写给真福弟茂德时，宗徒补充说：\BibleQuote{可信的话，若我们与祂同死，也必与祂同活；若我们受苦，也必与祂一同为王；若我们否认祂，祂也必否认我们；若我们不信，祂仍是信实的，祂不能否认自己。}那么，\InnerQuote{祂不能}这句话又表明无限的大能，因为即使所有人都否认祂，祂说天主就是祂自己，且只能按祂自己的本性存在，因为祂的存在是不可毁灭的。这就是\InnerQuote{祂不能否认自己}这句话的意思。因此，不可能变得更坏证明了无限的能力。\par

\Person{Eran}——这完全真实，且与神圣的话一致。\par

\Person{Orth}——既已承认在天主那里有许多事不可能——即一切与天主性相抵触之事——那么，为何你忽略所有属于天主性的其他特性——良善、公义、真理、不可见性、不可理解性、无限和永恒，以及我们断言为天主所固有的其余属性——却唯独主张祂的不死性和不能受苦性是可变化的，在其中承认变化的可能性，并赋予天主一个表明软弱的能力呢？\par

\Person{Eran}——我们从圣经学到了这一点。神圣的若望呼喊：\BibleQuote{天主如此爱了世界，竟赐下祂的独生子，}而神圣的保禄说：\BibleQuote{如果我们还是仇敌时，藉祂圣子的死与天主和好了，那么，和好之后，更要藉祂的生命得救。}\par

\Person{Orth}——当然这一切都是真的，因为这些是神圣的谕言，但请记住我们常宣认的事。\par

\Person{Eran}——什么事？\par

\Person{Orth}——我们宣认了：天主圣言、天主子，并非没有身体而显现，而是取了完美的人性。\par

\Person{Eran}——是的；我们宣认了这事。\par

\Person{Orth}——祂被称为人子，因为祂取了身体和人的灵魂。\par

\Person{Eran}——真的。\par

\Person{Orth}——因此主耶稣基督确实是我们天主；因为这两性中，一性从永远就是祂的，另一性是祂所取的。\par

\Person{Eran}——毫无疑问。\par

\Person{Orth}——那么，作为人祂受难，作为天主祂保持不能受苦。\par

\Person{Eran}——那么圣经怎么说天主子受难呢？\par

\Person{Orth}——因为受苦的身体是祂的身体。但让我们这样看这事：当我们听到圣经说\BibleQuote{当依撒格年老时，他的眼睛昏花以致不能看见，}我们的心被带往何处，依托于什么，是依撒格的灵魂还是他的身体？\par

\Person{Eran}——当然是他身体。\par

\Person{Orth}——那么我们猜想他的灵魂也分享了失明的状态？\par

\Person{Eran}——当然不。\par

\Person{Orth}——我们断言只有他的身体失去了视觉？\par

\Person{Eran}——是的。\par

\Person{Orth}——再者，当我们听到阿玛齐雅对先知亚毛斯说\BibleQuote{哦，先见者，你逃往犹大地去吧，}以及撒乌耳询问：\BibleQuote{请你告诉我，先见者的家在哪里，}我们不作任何身体性的理解。\par

\Person{Eran}——当然不。\par

\Person{Orth}——然而所用的词语表明视觉器官的健康。\par

\Person{Eran}——真的。\par

\Person{Orth}——然而我们知道圣神的能力赐予更纯净的灵魂时，激发起先知之恩，使它们甚至看见隐藏的事物，因而它们被称为先见者和见证者。\par

\Person{Eran}——你所说的是真的。\par

\Person{Orth}——让我们也考虑这一点。\par

\Person{Eran}——什么？\par

\Person{Orth}——当我们听到神圣福音使者叙述他们如何将一位患瘫痪病的人带到天主面前，躺在一张床上时，我们说是灵魂的部分瘫痪还是身体的部分瘫痪？\par

\Person{Eran}——显然是身体。\par

\Emph{正统}：现在我们读到希伯来书，看到宗徒所说的\BibleQuote{所以，你们要把下垂的手、发软的膝挺起来，为你们的脚修直道路，免得瘸子偏离正道，反而使他得到痊愈}时，我们是否认为天主的宗徒这些话说的是身体的部分？\par

\Emph{厄兰}：不是。\par

\Emph{正统}：我们是否应该说，他是在消除灵魂的软弱和病弱，激发门徒们的勇气？\par

\Emph{厄兰}：显然如此。\par

\Emph{正统}：但我们在圣经中并没有发现这些区分；他在描述依撒格的眼瞎时，并没有提到身体，而是说依撒格完全瞎了；在描述先知们为\Quote{先见者}和\Quote{观望者}时，也没有说他们的灵魂看见了、观望了隐藏的事，而是提到了他们本人。\par

\Emph{厄兰}：是的，是这样。\par

\Emph{正统}：他也没有指出瘫痪者的身体瘫痪了，而是称那人为瘫痪者。\par

\Emph{厄兰}：没错。\par

\Emph{正统}：甚至天主的宗徒也没有特别提到灵魂，尽管他的目的是要坚固和激发它们。\par

\Emph{厄兰}：对，他没有。\par

\Emph{正统}：但当我们探究这些话的含义时，我们就理解哪些属于灵魂，哪些属于身体。\par

\Emph{厄兰}：这是很自然的，因为天主造了我们是有理性的。\par

\Emph{正统}：那么，让我们对我们的造主和救主也运用这种理性能力，让我们认出哪些属于祂的天主性，哪些属于祂的人性。\par

\Emph{厄兰}：但这样做，我们会破坏至高的结合。\par

\Emph{正统}：在依撒格、先知、瘫痪者以及其他人的例子中，我们这样做并没有破坏灵魂与身体的本性结合；我们甚至没有将灵魂与它们各自的身体分开，而是仅凭理性区分了哪些属于灵魂，哪些属于身体。那么，既然我们在灵魂和身体的事上采取这种做法，却拒绝在我们的救主身上如此行，并把本性区别如此巨大的（其差别不像灵魂与身体那样，而是如同暂时与永恒、造物主与受造物之间有天壤之别）本质混淆起来，这难道不是荒谬的吗？\par

\Emph{厄兰}：圣经说天主子受了苦难。\par

\Emph{正统}：我们并不否认是祂受了苦难，但受圣经的教导，我们知道天主性的本质是不能受苦的。我们知道了不能受苦和受苦、人性和天主性，因此我们把苦难归于能受苦的身体，并宣认那不能受苦的本质没有受任何苦难。\par

\Emph{厄兰}：那么，是一个身体为我们赢得了救恩。\par

\Emph{正统}：是的，但不是一个普通人的身体，而是我们的主耶稣基督、天主的独生子的身体。如果你认为这个身体微不足道且无足轻重，那么你怎么能认为它的预像是敬拜的对象和救恩的工具呢？而且，如果预像是可敬和尊贵的，那么原型又怎么会是可鄙和微不足道的呢？\par

\Emph{厄兰}：我并不认为身体无足轻重，但我反对将它从天主性中分割开来。\par

\Emph{正统}：我的好先生，我们并不分割结合，而是注重各本质的特性，而且我相信你很快就会持同样的看法。\par

\Emph{厄兰}：你说得像个先知。\par

\Emph{正统}：不，不是像先知，而是因为知道真理的能力。但现在回答我这个问题：当你听到主说\BibleQuote{我与父原是一体}和\BibleQuote{看见我的就是看见了父}时，你认为这是指肉体还是指天主性？\par

\Emph{厄兰}：肉体与父怎么可能同为一体呢？\par

\Emph{正统}：那么，这些经文指示的是天主性？\par

\Emph{厄兰}：是的。\par

\Emph{正统}：同样，经文说：\BibleQuote{在起初已有圣言，圣言就是天主}等等。\par

\Emph{厄兰}：同意。\par

\Emph{正统}：再者，当圣经说：\BibleQuote{耶稣因旅途劳顿，就坐在井旁}，这劳顿是应归于天主性还是身体？\par

\Emph{厄兰}：我不能忍受分割那联合的。\par

\Emph{正统}：那么，看来你把劳顿归给了神圣本质？\par

\Emph{厄兰}：我是这样认为。\par

\Emph{正统}：但这样你就直接违背了先知的宣告：\BibleQuote{他不疲倦，也不困乏；他的智慧无法测度。他赐给疲倦者力量，给无力者增添能力。}稍后又说：\BibleQuote{但那些仰望上主的人必获得新力量，必像鹰一样振翅高飞；他们奔跑而不困乏，行走而不疲倦。}现在，那赐予他人不困乏和无需要的恩赐者，自己又怎会受饥渴呢？\par

\Emph{厄兰}：我一遍又一遍地说过，天主是不能受苦的，且没有任何需要，但降生成人后，祂变得能够受苦了。\par

\Emph{正统}：但祂这样做，是让祂的天主性承受苦难呢，还是允许那能受苦的本质经历它本性的苦难，并通过这苦难宣告那被看见的并非虚影，而是真实地取自人性？但现在让我们这样看问题：我们说天主性的本质是无限的。\par

\Emph{厄兰}：对。\par

\Emph{正统}：无限的本质不受任何限制。\par

\Emph{厄兰}：当然不是。\par

\Emph{正统}：因此它不需要移动，因为它在各处。\par

\Emph{厄兰}：是的。\par

\Emph{正统}：不需要移动的也就不需要旅行。\par

\Emph{厄兰}：那很明显。\par

\Emph{正统}：不旅行的就不疲倦。\par

\Emph{厄兰}：不。\par

\Emph{正统}：那么，那无限、不需要移动的天主性本质是不会疲倦的。\par

\Emph{厄兰}：但圣经说耶稣疲倦了，而耶稣是天主；\BibleQuote{我们的主耶稣基督，万物都是藉着祂而有的。}\par

\Emph{正统}：但圣经的确切说法是耶稣\Quote{疲倦了}而不是\Quote{常疲倦}。我们必须思考如何能把这两者都应用于同一位。\par

\Person{Eran}：那麼，請設法指出這一點來，因為你總是在強迫我們接受名詞的區別。\par

\Person{Orth.}：我認為即使一個蠻夷之人也能輕易作出這區別。既然承認兩種不相像的本性聯合，基督的位格因著這聯合而兼有兩者；各本性本身的屬性歸於各本性：無界限的不受累乏歸於那不受限制的，而會轉移、旅行的本性則受累乏。因為旅行是腳的職能；肌肉則因過度運動而緊繃。\par

\Person{Eran}：對於這些是身體的受苦，沒有爭議。\par

\Person{Orth.}：那麼，我所做的預言——你曾譏笑它——已經應驗了；因為你看：你已向我們展示了什麼屬於人性，什麼屬於神性。\par

\Person{Eran}：但我並沒有將一位兒子分為兩位。\par

\Person{Orth.}：我的朋友，我們也沒有這樣做；而是留意本性的區別，我們思考什麼適宜於神性，什麼屬於身體。\par

\Person{Eran}：這區別並非聖經的教導；聖經說天主之子死了。所以宗徒說：\BibleQuote{因為我們作仇敵的時候，曾藉著祂兒子的死，得以與天主和好。}\newline{}他又說主從死人中復活，因為\BibleQuote{天主叫主從死人中復活。}\par

\Person{Orth.}：當聖經說\BibleQuote{虔誠的人把斯德望送去埋葬，為他大聲哀哭}時，有人會說他的靈魂與身體一同被埋葬嗎？\par

\Person{Eran}：當然不會。\par

\Person{Orth.}：當你聽先祖雅各伯說\BibleQuote{請將我與我的祖先同葬}時，你認為這是指身體還是靈魂？\par

\Person{Eran}：毫無疑問是指身體。\par

\Person{Orth.}：現在請讀下文。\par

\Person{Eran}：\BibleQuote{在那裡葬了亞巴郎和他的妻子撒辣，在那裡葬了依撒格和他的妻子黎貝加，我也在那裡葬了肋阿。}\par

\Person{Orth.}：剛才你所讀的經文中，聖經沒有提及身體，但就所用詞句而言，也指靈魂。然而我們作出適當的區別，說先祖們的靈魂是不死不滅的，只有他們的身體被葬在雙重洞穴裡。\par

\Person{Eran}：對。\par

\Person{Orth.}：而當我們讀\BibleBook{宗徒}{大事錄}中黑落德用刀劍殺了若望的兄弟雅各伯時，我們不會認為他的靈魂死了。\par

\Person{Eran}：不，我們怎能這樣想？我們記得主的警告：\BibleQuote{不要害怕那些殺害身體而不能殺害靈魂的。}\par

\Person{Orth.}：但對普通人而言，避免靈魂與身體的必然聯繫，且在經文提及死亡和埋葬時，心思中將靈魂與身體分開，只將它們與身體聯繫，同時因信賴主的教導而認為靈魂不朽；然後當你聽聞天主之子的苦難時，卻採取完全不同的做法——這難道不令你覺得不敬和荒唐嗎？你絕口不提苦難所屬的身體，反倒將不能受苦、不變、不朽的神性描繪成可朽和能受苦的？而你明知如果天主聖言的本性是可受苦的，那麼取用身體便是多餘的。\par

\Person{Eran}：我們從聖經中學到，天主之子受了苦。\par

\Person{Orth.}：但神聖的宗徒解釋了苦難，並指明什麼本性受了苦。\par

\Person{Eran}：請立刻向我展示這一點，把問題弄清楚。\par

\Person{Orth.}：你是否不熟悉\BibleBook{}{致希伯來人書}中那段話，神聖的保祿說：\BibleQuote{所以，祂稱他們為弟兄，也不以為恥，說：\InnerQuote{我要將你的名傳與我的弟兄，在教會中我要歌頌你。}又說：\InnerQuote{看哪，我與天主所給我的兒女。}}\par

\Person{Eran}：是的，我知道這段，但它並沒有給出你所應許的。\par

\Person{Orth.}：不，甚至這些已暗示了我應許要展示的。弟兄一詞表示親屬關係，而親屬關係是由於取用了本性，取用則公開宣告了神性的不能受苦。但為更清楚地理解這一點，請讀下文。\par

\Person{Eran}：\BibleQuote{兒女既同有血肉之體，祂也照樣親自成了血肉之體，特要藉著死，敗壞那掌死權的……並要釋放那些一生因怕死而為奴僕的人。}\par

\Person{Orth.}：我想這無需解釋；它清楚地教導了救恩計劃的奧蹟。\par

\Person{Eran}：我在這裡看不到你所應許的證明。\par

\Person{Orth.}：然而神聖的宗徒清楚地教導：造物主憐憫這本性——它不僅被死亡殘酷地攫取，而且一生為死亡之奴——藉著一個身體為我們的身體實現了復活，並藉著一個可朽的身體消除了死亡的統治；因為祂自己的本性是不朽的，祂義正言辭地願意藉著取用屬於死亡者的初果來終止死亡的統治，同時使這初果（即身體）無可指責、無罪；一方面祂允許死亡對它下手，以滿足其貪婪，另一方面，因著對這身體所行的不義，祂制止了對其餘所有人所篡奪的不義統治。祂使這些被不義吞沒的初果復活，並將帶領整個種族跟隨它們。\par

將這解釋與宗徒的話並列，你就會理解神性的不能受苦。\par

\Person{Eran}：所讀的經文中沒有神聖不能受苦的證明。\par

\Person{Orth.}：不然：神聖宗徒的話——祂使兒女同有血肉之體的原因，是要藉著死敗壞那掌死權的——豈不是清楚表明神性的不能受苦和肉身的能受苦，並且因為神性不能受苦，祂才取了那能受苦的本性，藉著它敗壞了魔鬼的權勢？\par

\Person{Eran}：祂如何藉著肉身上敗壞魔鬼的權勢和死亡的統治？\par

\Person{Orth.}：魔鬼在最初奴役人性时，使用了什么武器？\par

\Person{Eran.}：它俘虏那曾被立为乐园公民者的手段，是罪。\par

\Person{Orth.}：天主为违犯诫命规定了什么惩罚？\par

\Person{Eran.}：死亡。\par

\Person{Orth.}：那么，罪是死亡之母，魔鬼是死亡之父。\par

\Person{Eran.}：正是。\par

\Person{Orth.}：因此，一场战争是由罪向人性发动的。罪诱使那些服从它的人沦为奴隶，把他们带到自己卑劣的父亲那里，并交给它那极其苦涩的后代。\par

\Person{Eran.}：这很明显。\par

\Person{Orth.}：所以，造物主为了摧毁这两种权势，合情合理地取了那正遭受战争攻击的本性，并使之远离一切罪，既使它从魔鬼的统治下获得自由，又借它摧毁了魔鬼的权势。因为死亡是罪人的惩罚，而死亡不义地、违背神圣法律地抓住了主无罪的肉身，祂首先复活了那被非法扣留的，然后应许给那些被正义监禁的人释放。\par

\Person{Eran.}：但是，你认为那被非法扣留者的复活，何以能由那些曾被正义交付给死亡的身体共同分享呢？\par

\Person{Orth.}：你认为，当亚当违犯诫命时，他的后代跟随先祖，这又何以是正义的呢？\par

\Person{Eran.}：虽然后代没有参与那著名的违犯，但他们犯了别的罪，并因此招致了死亡。\par

\Person{Orth.}：然而，并非只有罪人，义人、圣祖、先知、宗徒以及许多在各类美德上发光的人，也都落入了死亡的网罗。\par

\Person{Eran.}：是的；因为出自必死父母的家族，怎能保持不死呢？亚当在违犯和神圣判决之后，并处于死亡权势之下，才认识了他的妻子，被称为父亲；他自身成为可朽的，便成了可朽者的父亲；因此，所有领受了可朽本性的人合理地跟随他们的先祖。\par

\Person{Orth.}：你很好地说明了我们有份于死亡的原因。关于复活，也必须承认同样的道理，因为药必须对症。当人类的元首被定罪时，整个种族都与他一同被定罪；同样，当救主摧毁了诅咒时，人性获得了自由；正如分享亚当本性的人随他下到阴府，同样，全体人性将在主基督的复活中与他一同分享新生命。\par

\Person{Eran.}：教会的法令不仅应以断然的方式给出，还应以论证的方式给出。请告诉我，这些道理在圣经中是如何教导的。\par

\Person{Orth.}：请听宗徒写给罗马人、并通过他们教导全人类的话：\BibleQuote{如果因一人的过犯，众人都死了，那么天主的恩宠和那因一人耶稣基督的恩宠所赐予的恩惠，更要丰富地洋溢到众人身上。这恩赐不同于那因一人犯罪而来的判决；因为审判是由一人而定罪，恩赐却是由许多过犯而成义。如果因一人的过犯，死亡就因一人而作了王，那么那些领受了丰盛恩宠和正义恩赐的人，更要因一人耶稣基督而在生命中作王。}又说：\BibleQuote{因此，就如因一人的过犯，审判归于众人而定罪；同样，因一人的正义，恩赐归于众人而成义得生命。正如因一人的悖逆，众人成了罪人；同样，因一人的服从，众人将要成为义人。}当他在给格林多人介绍关于复活的论证时，他简短地向他们揭示了{\OriginalTerm{救恩计划}{\GreekText{οἰκονομία}}}的奥秘，说：\BibleQuote{但是，基督已从死者中复活，成为安眠者中的初果。因为死亡既因一人而来，死者的复活也因一人而来。就如在亚当内众人都死了，同样，在基督内众人也都要复活。}这样，我从神圣的圣言中给你带来了证据。现在，看看属于亚当的与属于基督的对比：疾病与药物，伤口与膏药，罪与正义的财富，诅咒与祝福，定罪与释放，违犯与遵守，死亡与生命，地狱与天国，亚当与基督，人与那人。然而，主基督不仅是人，也是永恒的天主，但神圣的宗徒是从他所取的本性称呼他的，因为他正是用这个本性与亚当比较。成义、战斗、胜利、死亡、复活，全都属于这个人性；正是这个本性，我们与他共享；在这个本性中，那些事先在天国公民生活中操练的人，将与他一同为王。我谈论这个本性，并非分开神性，而是指属于人性的部分。\par

\Person{Eran.}：你对此进行了冗长的讨论，并用圣经的见证加强了你的论点。但是，如果受难确实属于肉身，那么当宗徒歌颂天主对人的爱而呼喊时：\BibleQuote{祂没有怜惜自己的儿子，反而为我们众人把祂交出来了}，他说被交出来的是哪个儿子呢？\par

\Person{Orth.}：你要谨慎你的言辞。天主子只有一位，因此祂被称为独生子。\par

\Person{Eran.}：那么，如果天主子只有一位，神圣的宗徒称祂为自己的儿子。\par

\Person{Orth.}：是的。\par

\Person{Eran.}：然后他说祂被交出来了。\par

\Person{Orth.}：是的，但不是没有身体，这一点我们已再三同意。\par

\Person{Eran.}：已经再三同意祂取了身体和灵魂。\par

\Person{Orth.}：因此，宗徒所说的是有关身体的事。\par

\Person{Eran.}：但神圣的宗徒明确说：\BibleQuote{祂没有怜惜自己的儿子。}\par

\Person{Orth.}：那么，当你听见天主对亚巴郎说：\BibleQuote{因为你没有留住你的儿子，你的独生子，}你是否认为依撒格被杀了呢？\par

\Person{Eran.}：当然不是。\par

\Person{Orth.}：然而天主说：\BibleQuote{你没有留住，}而万有的天主是真实的。\par

\Emph{Eran.}——\Quote{你沒有留住}一語指的是亞巴郎的準備，因為他準備好祭獻那少年，但天主阻止了。\par

\Emph{Orth.}——好；在亞巴郎的故事中，你不滿足於字面，而是展開它並使意義清晰。用完全相同的方式，審視宗徒話語的意義。那麼你將看到，被保留的不是神性本性，而是被釘在十字架上的肉體。甚至在預象中，真理也容易察覺。你認為亞巴郎的祭獻是為世界所作的奉獻的預象嗎？\par

\Emph{Eran.}——一點也不，我也不能將在教會中修辭性地說出的話作為信仰的規則。\par

\Emph{Orth.}——你理應跟隨教會的導師，但既然你不當地反對他們，聽聽救主親自對猶太人所說的話：\BibleQuote{你們的父親亞巴郎曾歡欣喜樂地期望看見我的日子，他看見了，就歡樂起來。}注意主稱祂的苦難為\Quote{一個日子}。\par

\Emph{Eran.}——我接受主的見證，並不懷疑預象。\par

\Emph{Orth.}——現在比較預象和實體，你將看到即使在預象中，神性的不受苦性。在前者和後者中，都有一位父親；在前者和後者中，都有一位愛子，各自背負祭獻的材料。一個背負木柴，另一個背負十字架在肩上。據說，山頂因雙方的祭獻而尊榮。此外，日夜的數目與隨後的復活相對應，因為依撒格被他父親願意的心靈殺死後，在好施的天主命令此事發生的第三天，經由愛人者之聲，他復活到新生命。一隻公綿羊被看見困在灌木叢中，提供了十字架的形象，並代替少年被宰殺。那麼，如果這是實體的預象，而在預象中獨生子沒有受祭，而是一隻公綿羊被取代並放在祭台上完成了奉獻的奧蹟，那麼為什麼在實體中你猶豫將苦難歸於肉體，並宣告神性的不受苦性？\par

\Emph{Eran.}——在你對這個預象的觀察中，你描述依撒格因天主的命令而再次活過來。因此，如果我們將實體配合預象，宣告天主聖言受苦並重新活過來，沒有什麼不當。\par

\Emph{Orth.}——我一次又一次地說過，預象不可能在每一方面都與原型的實體相匹配，這在當前情況下也容易理解。依撒格和公綿羊，就其本性的差異而言，適合形象，但就其分離位格的劃分而言，則不再適合。我們宣講神性和人性如此緊密的結合，以至於理解一個不可分的位格，並承認同一位既是天主又是人，可見又不可見，有限又無限，並將所有表示神性和人性的屬性應用於這一位格之一。既然公綿羊，一個無理性的存在，沒有天主肖像的恩賜，不可能預表復活，這兩者就分擔了救恩奧蹟的預象，一個提供死亡的形象，另一個提供復活。我們在梅瑟的祭獻中發現完全相同的事，因為在其中也可以看到預先勾勒的救恩苦難的預象。\par

\Emph{Eran.}——什麼梅瑟的祭獻預示實體？\par

\Emph{Orth.}——可以說，全部舊約都是新約的預象。正是為此，神聖的宗徒清楚地說：\BibleQuote{法律擁有未來美物的陰影}，又說\BibleQuote{現在這一切事發生在他們身上，作為預象}。原型的形象非常清楚地由在埃及被宰殺的羔羊，以及在營外焚燒的紅母牛所展示，而且宗徒在致希伯來人書中也提到，他寫道：\BibleQuote{因此耶穌為用自己的血聖化人民，也在城門外受苦。}\par

但暫時不再多說這個。然而，我將提到那祭獻，其中有兩隻公山羊被獻上，一隻被宰殺，另一隻被釋放。在這兩隻公山羊中，有救主兩性本性的預先形象：在被釋放的那隻中，是不受苦的神性；在被宰殺的那隻中，是可受苦的人性。\par

\Emph{Eran.}——你不認為將主比作山羊是不敬嗎？\par

\Emph{Orth.}——你認為什麼是更應避免和憎恨的對象，蛇還是山羊？\par

\Emph{Eran.}——蛇顯然是可憎的，因為它傷害接近它的人，並且常常傷害不傷害它的人。另一方面，按照法律，山羊屬於潔淨可食用的動物之列。\par

\Emph{Orth.}——現在聽主將救恩的苦難比作銅蛇。祂說：\BibleQuote{正如梅瑟在曠野裡高舉了蛇，人子也同樣要被高舉起來，使一切信祂的人不致滅亡，反而獲得永生。}如果銅蛇是被釘十字架的救主的預象，那麼我們將救恩的苦難與公山羊的祭獻相比，犯了什麼不當？\par

\Emph{Eran.}——因為若望稱主為\Quote{羔羊}，依撒意亞也稱祂為\Quote{羔羊}和\Quote{綿羊}。\par

\Emph{Orth.}——但蒙福的保祿稱祂為\Quote{罪}和\Quote{詛咒}。因此，作為詛咒，祂應驗了被詛咒之蛇的預象；作為罪，祂說明了公山羊祭獻的圖像，因為在法律中，為罪獻上的是公山羊，而不是羔羊。所以主在福音中將義人比作羔羊，卻將罪人比作山羊羔；既然祂不僅為義人，也為罪人而受難，祂適當地通過羔羊和公山羊預示自己的奉獻。\par

\Emph{Eran.}——但兩隻公山羊的預象引導我們想到兩個位格。\par

\Emph{Orth.}— 人性之\OriginalTerm{可受苦性}{passibility}与天主性之\OriginalTerm{不受苦难性}{impassibility}不可能同时由一只山羊预表。被宰杀的那只无法显示活着的本性。因此取了两只，以解释两种本性。从另一祭献也能学到同样的教训。\par

\Emph{Eran.}— 从哪个祭献？\par

\Emph{Orth.}— 从那律法者命人献两只洁净鸟的祭献——一只宰杀，另一只蘸于被杀之鸟的血中后放走。这里我们也看到了天主性及人性的\OriginalTerm{预像}{type}——被宰杀的人性与那承受苦难的天主性。\par

\Emph{Eran.}— 你给了我们许多预像，但我反对隐晦的比喻。\par

\Emph{Orth.}— 但神圣的宗徒称这些叙述为\OriginalTerm{预像}{type}。哈加尔被称为旧盟约的预像；撒辣被比作天上的耶路撒冷；依市玛耳是以色列的预像，依撒格则是新子民的预像。所以你必须控告圣神那响亮的号角，因为它为我们众人提供了隐晦的比喻。\par

\Emph{Eran.}— 即使你提出任何数量的论据，也永远无法说服我分割苦难。我听到了天使对玛利亚和她的同伴所说的话：\Quote{你们来看安放主的地方。}\par

\Emph{Orth.}— 这完全符合我们通常的习惯：我们以整体之名称呼部分。当我们进入安葬圣宗徒或先知或殉道者的教堂时，我们时常问：\Quote{谁躺在圣龛中？}那些能告知我们的人回答说，也许是宗徒多默，或洗者若翰，或首位殉道者斯德望，或任何其他圣人，提及他们的名字，尽管可能只有少量\OriginalTerm{圣髑}{relics}在此。但听见这些灵魂与身体共有的名字的人，不会想象灵魂也被关在箱子里；人人都知道箱中只含有身体或身体的小部分。圣天使正是以同样的方式说话，他用人的名字来描述身体。\par

\Emph{Eran.}— 但你如何证明天使对妇女们说的是关于主的身体？\par

\Emph{Orth.}— 首先，坟墓本身足以解决这个问题，因为坟墓既不交付灵魂，也不交付其本性是\OriginalTerm{不受局限的}{uncircumscribed}的天主性；坟墓是为身体而造的。此外，圣经清楚地教导了这点，因为圣玛窦如此叙述该事件：\Quote{到了晚上，有一个亚利马太的财主，名叫若瑟，他也是耶稣的门徒。他去见比拉多，求耶稣的身体，比拉多就吩咐将身体交给他。若瑟取了身体，用洁净的细麻布包裹，安放在自己的新坟墓里，就是他凿在岩石里的；他又把一块大石头滚到墓门口，就离开了。}请看，他屡次提到身体，为堵住亵渎天主性之人的口。三倍有福的马尔谷也遵循同一方式，我也要引用他的叙述：\Quote{到了晚上，因为是预备日，即安息日的前一日，亚利马太的若瑟，一位尊贵的议员，也是等待天主国的，大胆地进见比拉多，求耶稣的身体。比拉多诧异耶稣已经死了，便叫百夫长来，问他耶稣是否已死。既从百夫长得知实情，就把身体赐给若瑟。若瑟买了细麻布，把耶稣取下来，用细麻布裹好，安放在坟墓里。}等等。请赞叹术语的和谐，以及身体一词如何一贯而连续地出现。杰出的路加也以同样的方式叙述若瑟求身体，并领受后按礼仪处理。神圣的若望告诉我们更多：\Quote{亚利马太的若瑟，因为是耶稣的门徒，只因怕犹太人，就暗中去求比拉多，要把耶稣的身体领去。比拉多允准，他就来把耶稣的身体领去了。又有从前夜间来见耶稣的尼苛德摩，带着约一百斤的没药和沉香调和的香料。他们取了耶稣的身体，用香料和细麻布裹好，照犹太人殡葬的习俗。在耶稣被钉十字架的地方有一个园子，园子里有一座新坟墓，是从来没有葬过人的。只因是犹太人的预备日，那坟墓近便，他们就把耶稣安放在那里。}请观察身体被提及多少次；福音作者如何显示被钉在十字架上的是身体，是若瑟向比拉多求的身体，是从木架上取下的身体，是用没药和沉香和细麻布裹好的身体，然后以人的名字称呼它；记载耶稣被安放在坟墓中。因此天使说：\Quote{你们来看安放主的地方。}以整体之名称呼部分；我们也常这样做。我们在此地说：某人葬在这里；而不是某人的身体。每个有理智的人都知道我们指的是身体，这种说话方式在圣经中是惯用的。我们读到亚郎死了，他们把他葬在曷尔山上。撒慕尔死了，他们把他葬在辣玛，还有许多类似的例子。神圣的宗徒在论及主的死亡时也遵循同一用法。他写道：\Quote{我当日传给你们的，原是我所领受的，就是基督照圣经所说，为我们的罪死了，而且埋葬了，又照圣经所说，第三天复活了。}等等。\par

\Emph{Eran.}— 在刚才所读的段落中，宗徒并未提到身体，而是提到我们众人的救主基督。你引用了反对你自身的证据，用自己的武器伤了自己。\par

\Emph{Orth.}——你似乎很快就忘记了我反复向你证明的那段长篇论述，即身体是以人名称呼的。这正是如今神圣的宗徒所行的，而从这段经文本身便可轻易证明。现在我们就来看一看。为什么神圣的作者要这样写信给格林多人呢？\par

\Emph{Eran.}——他们被一些人欺骗，相信没有复活。世界的导师得知此事后，就向他们提出了关于身体复活的论证。\par

\Emph{Orth.}——那么，当他想要证明身体的复活时，为什么他要引入主的复活呢？\par

\Emph{Eran.}——因为这足以证明我们所有人的复活。\par

\Emph{Orth.}——祂的死与其他人的死有什么相似之处，以致于能通过祂的复活来证明所有人的复活呢？\par

\Emph{Eran.}——天主的独生子降生成人、受苦受死的原因，乃是为了摧毁死亡。因此，在复活之后，祂以自己的复活宣讲了所有人的复活。\par

\Emph{Orth.}——但是，听到天主的复活，谁会相信所有人的复活会与它完全一样呢？本性的差异不允许我们相信复活的论证。祂是天主，而他们是人，天主与人之间的差距是不可估量的。他们是有死的，且受制于死亡，如同草和花一般。祂则是全能的。\par

\Emph{Eran.}——但在降生成人之后，天主圣言拥有一个身体，并借此证明了祂与人的相似。\par

\Emph{Orth.}——是的；正因如此，受苦、死亡和复活都属于身体；为证明这一点，神圣的宗徒在另一处向所有人应许了生命的更新；对于那些相信救主复活、却又视所有人的普遍复活为神话的人，他大声宣告：\BibleQuote{既然基督被传报说祂从死者中复活了，怎么在你们中间有人说没有死人的复活呢？如果没有死人的复活，基督也就没有复活；如果基督没有复活……你们的信德便是虚妄的，你们仍在你们的罪恶中。}他从过去印证将来，又从被怀疑的事驳斥那被相信的事，因为他说：如果那一件对你们看来不可能，那么这一件便是虚假的；如果那一件看来真实可信，那么这一件也同样可以视为真实，因为此处所宣讲的正是身体的复活，而这身体被称为那些人的初果。在多次论证之后，他直接肯定了这身体的复活：\BibleQuote{但是基督确实从死者中复活了，成为了那些已睡之人的初果，因为死亡既由一人而来，死人的复活也由一人而来；正如在亚当内众人都死了，照样在基督内众人都要复活。}他不仅确认了复活的论证，还揭示了\OriginalTerm{奥迹}{\GreekText{οἰκονομία}}的奥秘。他称基督为\Quote{人}，为的是证明补救与疾病相称。\par

\Emph{Eran.}——那么基督只是一个人吗？\par

\Emph{Orth.}——千万不要这样说。相反，我们一再公开承认祂不仅是人，而且是永生的天主。但祂是作为人受苦，而非作为天主。神圣的宗徒清楚地教导我们这一点，当他说\BibleQuote{因为死亡既由一人而来，死人的复活也由一人而来。}在写给得撒洛尼人的书信中，他借着我们救主的复活加强了关于普遍复活的论证，经文说\BibleQuote{因为我们若信耶稣死了，也复活了，同样，那些在耶稣内睡了的人，天主也要把他们与耶稣一同带来。}\par

\Emph{Eran.}——宗徒借着主的复活来证明普遍的复活，很明显，在此情况下死去并复活的是一个身体。因为除非两者的实体之间有某种关联，否则他绝不会试图借此来证明普遍的复活。我决不会同意将苦难仅仅归于人性。在我看来，说\Quote{天主圣言在肉体中死了}似乎更合乎我的观点。\par

\Emph{Orth.}——我们已多次指出，那在本性上不朽的，绝不可能死去。那么，如果祂死了，祂就不是不朽的；而这句话中蕴含的亵渎是何等的危险！\par

\Emph{Eran.}——祂按本性是不朽的，但祂成了人，并受了苦。\par

\Emph{Orth.}——因此，祂经历了变化，因为否则祂作为不朽者怎能受制于死亡呢？但我们已经同意，天主圣三的实体是不变的。既然祂具有超越变化的本性，祂绝无分于死亡。\par

\Emph{Eran.}——神圣的伯多禄说\BibleQuote{基督在肉身内为我们受了苦。}\par

\Emph{Orth.}——这与我们所说的相符，因为我们从圣经中学会了教义的规则。\par

\Emph{Eran.}——那么，你怎能否认天主圣言在肉体中受苦呢？\par

\Emph{Orth.}——因为我们没有在圣经中找到这种说法。\par

\Emph{Eran.}——但我刚才引用了伟大的伯多禄的话。\par

\Emph{Orth.}——你似乎忽略了术语的区别。\par

\Emph{Eran.}——什么术语？难道你不认为主基督就是天主圣言吗？\par

\Emph{Orth.}——在我们主和救主的情况下，\Quote{基督}一词指的是降生成人的圣言、厄玛奴耳——天主与我们同在——既是天主也是人；而\Quote{天主圣言}一词单独使用时，指的是那在世界之前、超越时间、无形体的单纯本性。因此，通过圣宗徒发言的圣神，从未将苦难或死亡归于此名。\par

\Emph{Eran.}——如果苦难归于基督，而天主圣言在成为人之后被称为基督，那么我认为，说\Quote{天主圣言在肉体中受苦}的人绝无不合理之处。\par

\Emph{Orth.}——这种尝试极其危险和鲁莽。但让我们这样来看这个问题：圣经是否说天主圣言出自天主和圣父？\par

\Emph{Eran.}——是的。\par

\Emph{Orth.}——它同样描述圣神是出自天主的吗？\par

\Emph{Eran.}——同意。\par

\Emph{Orth.}——但它称天主圣言为独生子。\par

\Emph{Eran.}——是的。\par

\Emph{Orth.}——它从未这样称呼圣神。\par

\Emph{Eran.}——没有。\par

\Emph{Orth.}——然而圣神也是从圣父和天主那里获得祂的位格存在。\par

\Emph{Eran.}——确实。\par

\Emph{正统派}——那么，我们承认圣子和圣神皆出自天主圣父；但你敢称圣神为子吗？\par

\Emph{异端者}——当然不敢。\par

\Emph{正统派}——为什么？\par

\Emph{异端者}——因为我在圣经中找不到这个名词。\par

\Emph{正统派}——或是受生？\par

\Emph{异端者}——不。\par

\Emph{正统派}——为何？\par

\Emph{异端者}——因为我在圣经中同样没有学到这个。\par

\Emph{正统派}——但既非受生又非受造的，当以何名称为恰当？\par

\Emph{异端者}——我们称其为非受造且非受生。\par

\Emph{正统派}——而我们说圣神既非受造也非受生。\par

\Emph{异端者}——绝非如此。\par

\Emph{正统派}——那么你敢称圣神为非受生吗？\par

\Emph{异端者}——不敢。\par

\Emph{正统派}——但为何拒绝称那本性非受造却非受生的为非受生呢？\par

\Emph{异端者}——因为我未从圣经中学会这样，我极害怕说出或使用圣经未用的言辞。\par

\Emph{正统派}——那么，我善士，我在救恩的苦难上也持同样的谨慎；你也要避免一切圣经在苦难中避免的天主名号，不要将苦难归于它们。\par

\Emph{异端者}——什么名号？\par

\Emph{正统派}——苦难从不与\Quote{天主}之名相连。\par

\Emph{异端者}——但就连我也不是说天主圣言脱离身体受苦，而是说祂在肉身内受苦。\par

\Emph{正统派}——那么你肯定了一种受苦的方式，而非无苦性。甚至对于人的身体，也无人会这样说。因为除非完全丧失理智，谁会保禄的灵魂死在肉身内？这话甚至对一个极大的恶人也不会说；因为恶人的灵魂也是不死的。我们说某个谋杀者被杀了，但从无人会说他的灵魂在肉身中被杀。但如果我们说谋杀者和盗墓者的灵魂免于死亡，那么我们更应当承认我们救主的灵魂是不死的，因为它从未尝过罪。如果那些犯了大错之人的灵魂因其本性而免于死亡，那灵魂本性不朽且从未沾染丝毫罪污的，怎能中了死亡的钩子呢？\par

\Emph{异端者}——你对我讲这些冗长的论证毫无用处。我们同意救主的灵魂是不死的。\par

\Emph{正统派}——但你如何不配受罚？你说那本性受造的灵魂是不死的，却要使圣言的天主性体成为可朽的；你否认救主的灵魂在肉身中尝过死味，却敢主张造物主天主圣言受了苦难？\par

\Emph{异端者}——我们说祂无苦地受了苦难。\par

\Emph{正统派}——哪个有理智的人能忍受这样可笑的谜语？谁曾听过无苦的苦难，或不死的死亡？无苦的从未受过苦，受过苦的不能是无苦的。但我们听见圣保禄的呼喊：\BibleQuote{独有不死不灭，住在不可接近的光中。}\par

\Emph{异端者}——那么我们为何还说无形的权势和人的灵魂，甚至魔鬼自身，是不死的？\par

\Emph{正统派}——我们确是这样说；但天主是绝对不死的。祂不死不是由于分有实体，而是实体本身；祂的不死并非从他人领受。正是祂将自己不死性赐予天使和你刚才提到的那些。此外，当圣保禄称祂为不死并说祂独有不死时，你怎能将死亡的苦难归于祂？\par

\Emph{异端者}——我们说祂在降生成人后尝了死味。\par

\Emph{正统派}——但我们一再承认祂是不变的。如果祂先前是不死的，后来却通过肉身经历了死亡，在祂经历死亡之前先有改变；如果祂的生命离开了祂三天三夜，这些说法岂非极度不虔？因为我认为甚至那些与不虔斗争的人，若不冒险，也不敢让这样的话从自己口中说出。\par

\Emph{异端者}——别再说我们不虔。连我们也说不是天主性受苦，而是人性；但我们确实说天主性与身体一同分受苦。\par

\Emph{正统派}——你所说的分受苦是什么意思？你是说当钉子刺入身体时，天主性感觉到了疼痛？\par

\Emph{异端者}——正是。\par

\Emph{正统派}——现在和在先前的探讨中，我们都已证明灵魂并非与身体的一切官能共享，但身体在接受生命力时，通过灵魂而有受苦的感觉。即使我们承认灵魂与身体一同分受痛苦，我们仍会发现天主性是无苦的，因为它不是代替灵魂与身体结合的。难道你不承认祂取了灵魂吗？\par

\Emph{异端者}——我常承认。\par

\Emph{正统派}——并且祂取了理性的灵魂？\par

\Emph{异端者}——是的。\par

\Emph{正统派}——那么，既然祂连同身体取了灵魂，我们承认灵魂与身体一同分受苦；那么与身体一同分受苦的是灵魂，而非天主性；灵魂分受苦，通过身体接受痛苦。但也许有人会同意灵魂与身体一同分受苦，但否认它分受死亡，因为它有不死的本性。因此主说：\BibleQuote{不要害怕那些杀死身体却不能杀死灵魂的。}那么，如果我们否认救主的灵魂与身体一同分受死亡，又有谁能接受你和你同伙肆意传扬的亵渎之语，即你们敢说天主性参与了死亡？这更加不可原谅，因为主一方面指出身体被交付，另一方面指出灵魂在忧闷。\par

\Emph{\Person{埃朗}}——主在哪里显示那身体曾被献上？或者你又要给我引那段老生常谈的经句\BibleQuote{你们拆毁这座圣殿，三天之内我要把它重建起来}？还是你傲慢自足地要给我引用福音作者的话？\BibleQuote{但耶稣所说的，是指他身体的圣殿。所以当他从死人中复活以后，他的门徒就想起他说过这话，便相信了圣经和耶稣所说的话。}\par

\Emph{\Person{正统}}——如果你如此厌恶那些宣讲降生成人奥迹的天主圣言，你为什么不像\Person{马尔基翁}、\Person{华伦提努}和\Person{摩尼}那样，不销毁这类经文呢？因为他们正是这样做的。但如果你认为这轻率而不敬，就不要把主的话当作笑柄，而要效法宗徒们，在复活后相信天主性将犹太人所毁坏的圣殿重建起来。\par

\Emph{\Person{埃朗}}——如果你有什么好的证据可以提出，就停止嘲笑，履行你的诺言吧。\par

\Emph{\Person{正统}}——请特别记住福音中那些话，主曾将吗哪与真粮作比较。\par

\Emph{\Person{埃朗}}——我记得。\par

\Emph{\Person{正统}}——在那段话里，主在详论生命之粮以后，接着说：\BibleQuote{我所要赐的粮，就是我的肉，为世人的生命而赐的。}这些话中既可以理解天主性的恩赐，也可以理解肉身的恩惠。\par

\Emph{\Person{埃朗}}——一段引文不足以解决这个问题。\par

\Emph{\Person{正统}}——\Person{埃塞俄比亚太监}并没有读过很多圣经，但当他从先知书中找到一个见证时，便受其引导而得救。然而，所有宗徒、先知以及此后一切宣讲真理的人，都不足以说服你。尽管如此，我还要为你带来一些关于主身体的进一步见证。你不可能不知道福音记载中的那段话：主与门徒吃过逾越节晚餐后，指出了那预表性羔羊的死，并教导那与这预表相对应的身体是什么。\par

\Emph{\Person{埃朗}}——是的，我知道。\par

\Emph{\Person{正统}}——那么请记住，我们的主拿起并擘开的是什么，以及他拿起后称它为什么。\par

\Emph{\Person{埃朗}}——为了未入门者的缘故，我要用奥秘的语言回答。他拿起、擘开并递给门徒后，说：\BibleQuote{这是我的身体，为你们而舍的——或照宗徒所说\InnerQuote{擘开的}——又说：\InnerQuote{这是我的血，新约的血，为多人流出来的。}}\par

\Emph{\Person{正统}}——那么，当展示受难的预表时，他没有提到天主性吗？\par

\Emph{\Person{埃朗}}——没有。\par

\Emph{\Person{正统}}——但他提到了身体和血。\par

\Emph{\Person{埃朗}}——是的。\par

\Emph{\Person{正统}}——那身体被钉在十字架上吗？\par

\Emph{\Person{埃朗}}——正是如此。\par

\Emph{\Person{正统}}——那么，请看这一点。复活后，当门都关着，主来到圣门徒们面前，见他们惊惧时，他用什么方法消除他们的恐惧，并在恐惧中注入信德呢？\par

\Emph{\Person{埃朗}}——他对他们说：\BibleQuote{你们看我的手和脚，就知道是我自己；你们摸摸看，因为幽灵没有肉和骨头，如同你们看见我有的。}\par

\Emph{\Person{正统}}——所以当他们不信时，他展示了身体？\par

\Emph{\Person{埃朗}}——是的。\par

\Emph{\Person{正统}}——因此身体复活了？\par

\Emph{\Person{埃朗}}——显然。\par

\Emph{\Person{正统}}——我想，复活的就是那已死去的？\par

\Emph{\Person{埃朗}}——正是如此。\par

\Emph{\Person{正统}}——那已死去的，就是被钉在十字架上的？\par

\Emph{\Person{埃朗}}——必然如此。\par

\Emph{\Person{正统}}——那么，根据你自己的论证，身体受了苦？\par

\Emph{\Person{埃朗}}——你的系列论证迫使我们得出这个结论。\par

\Emph{\Person{正统}}——再想想这一点。现在我来发问，你作为热爱真理的人来回答。\par

\Emph{\Person{埃朗}}——我会回答。\par

\Emph{\Person{正统}}——当圣神降临在宗徒们身上，那奇妙的景象和声音吸引了成千上万的人聚集到那房屋时，宗徒之长在他当时的讲论中，关于主的复活说了什么？\par

\Emph{\Person{埃朗}}——他引用了神性的\Person{达味}的话，说达味从天主那里得到应许，主基督将从他的腰腹所出的果实而生，并因信这些应许，他预言性地预见了主的复活，明确地说主的灵魂没有被遗弃在阴府，他的肉身也没有见腐朽。\par

\Emph{\Person{正统}}——因此，他的复活是属于这些的。\par

\Emph{\Person{埃朗}}——一个有理智的人怎么能说灵魂的复活呢？灵魂从未死去。\par

\Emph{\Person{正统}}——你这位将受难、死亡和复活归于不变且无限的天主性的人，怎么突然在我们面前显出理智，现在却反对把复活这个词与灵魂联系起来？\par

\Emph{\Person{埃朗}}——因为复活这个词适用于那已倒下的。\par

\Emph{\Person{正统}}——但身体若没有灵魂就不会获得复活，而是被天主的旨意更新，并与它的配偶结合后，才获得生命。主复活\Person{拉匝禄}不就是这样吗？\par

\Emph{\Person{埃朗}}——显然，复活的不仅仅是身体。\par

\Emph{\Person{正统}}——神性的\Person{厄则克耳}教导得更清楚，因为他指出主如何命令骨头聚拢，如何全部合适地连接起来，如何使筋、血管和动脉连同所有属于它们的肉和覆盖它们的皮生长起来，然后命令灵魂回到它们自己的躯体内。\par

\Emph{\Person{埃朗}}——这是真的。\par

\Emph{\Person{正统}}——但主的身体并没有经历这种腐朽，而是保持完好，并在第三天恢复了自己的灵魂。\par

\Emph{\Person{埃朗}}——同意。\par

\Emph{\Person{正统}}——那么死亡属于那受过苦的？\par

\Emph{\Person{埃朗}}——毫无疑问。\par

\Emph{\Person{正统}}——而当伟大的\Person{伯多禄}提到复活，以及神性的\Person{达味}也提到时，他们说他的灵魂没有被遗弃在阴府，但他的身体没有经历腐朽？\par

\Emph{\Person{埃朗}}——是的。\par

\Emph{\Person{正统}}——那么，不是天主性经历了死亡，而是身体因与灵魂分离而经历了死亡？\par

\Emph{\Person{埃朗}}——我无法容忍这些谬论。\par

\Emph{\Person{正统}}——但你是在与自己的论证作战；你称之为谬论的正是你自己的话。\par

\Emph{\Person{埃朗}}——你在诽谤我；这些话没有一句是我的。\par

\Emph{正教徒.}——假设有人问：哪种动物既是理性的又是会死的？另一个人回答：人。你会说这二人谁是那说法的解释者？提问者还是回答者？\par

\Emph{厄瓦尼斯特.}——回答者。\par

\Emph{正教徒.}——那么，我把那些论点说成是你的，这很正确吧？因为，我想，你在回答中，拒绝某些点，接受另一些点，从而确认了它们。\par

\Emph{厄瓦尼斯特.}——那么我不再回答了；你来回答。\par

\Emph{正教徒.}——我来回答。\par

\Emph{厄瓦尼斯特.}——你对宗徒的这些话有什么看法：\Quote{他们若知道，就不会把光荣的主钉在十字架上了}？在这段经文中，他没有提到身体或灵魂。\par

\Emph{正教徒.}——因此，你不应该把\Quote{在肉身内}这个词加进去——因为这是你为了贬低圣言的天主性而发明的巧妙手法——而必须把苦难归于圣言的天主性本身。\par

\Emph{厄瓦尼斯特.}——不，不。祂在肉身内受苦，但祂的非物质本性本身不能受苦。\par

\Emph{正教徒.}——啊！但不可在宗徒的话上加添什么。\par

\Emph{厄瓦尼斯特.}——当我们知道宗徒的意思时，加上省略的部分并没有什么荒谬的。\par

\Emph{正教徒.}——但在神圣的话语上加添什么，是野蛮和鲁莽的。解释所写的，揭示隐藏的意义，才是圣洁和虔诚的。\par

\Emph{厄瓦尼斯特.}——很对。\par

\Emph{正教徒.}——那么，我们两人在考查圣经的心意时，并不是在做不合理和不圣洁的事。\par

\Emph{厄瓦尼斯特.}——不是。\par

\Emph{正教徒.}——那么，让我们一起来看看似乎隐藏的内容。\par

\Emph{厄瓦尼斯特.}——当然。\par

\Emph{正教徒.}——伟大的保禄是否称神圣的雅各伯为主（的）兄弟？\par

\Emph{厄瓦尼斯特.}——他称了。\par

\Emph{正教徒.}——但我们应把祂视为兄弟，凭的是祂天主性的关系还是祂人性的关系？\par

\Emph{厄瓦尼斯特.}——我不会同意分开联合的本性。\par

\Emph{正教徒.}——但你在我们之前的探讨中常常分开它们，现在你也该这样做。告诉我：你是否说天主圣言是独生子？\par

\Emph{厄瓦尼斯特.}——我说是的。\par

\Emph{正教徒.}——而独生指的是唯一的儿子。\par

\Emph{厄瓦尼斯特.}——当然。\par

\Emph{正教徒.}——那么独生子不可能有兄弟？\par

\Emph{厄瓦尼斯特.}——当然不可能，因为如果祂有兄弟，就不会被称为独生子了。\par

\Emph{正教徒.}——那么，他们称雅各伯为主（的）兄弟，这是错的。因为主是独生子，而独生子不可能有兄弟。\par

\Emph{厄瓦尼斯特.}——不，但主并不是无形的，而真理的宣讲者只提及关乎天主性的事。\par

\Emph{正教徒.}——那么你如何证明宗徒的话是真的？\par

\Emph{厄瓦尼斯特.}——通过说雅各伯是按肉身与主有亲属关系。\par

\Emph{正教徒.}——看，你又引进了你反对的那种区分。\par

\Emph{厄瓦尼斯特.}——不可能以其他方式解释这种亲属关系。\par

\Emph{正教徒.}——那么，不要责备那些无法以其他方式解释类似困难的人。\par

\Emph{厄瓦尼斯特.}——现在你把论证引到岔路上了，因为你想回避问题。\par

\Emph{正教徒.}——一点也不，我的朋友。我们探讨过的要点也会解决那个问题。现在看：当你被提醒雅各伯是主的兄弟时，你说这关系不是指天主性，而是指肉身。\par

\Emph{厄瓦尼斯特.}——我确实说了。\par

\Emph{正教徒.}——那么，既然你被告知十字架的苦难，也请把这苦难归于肉身。\par

\Emph{厄瓦尼斯特.}——宗徒称被钉的那位为\Quote{光荣的主}，而同一位宗徒也称主为\Quote{雅各伯的兄弟}。\par

\Emph{正教徒.}——而这两处是同一个主。那么，如果你把亲属关系归于肉身的做法是正确的，你也必须把苦难归于肉身，因为不加区分地看待亲属关系，却把苦难不加区分地归于基督，这是完全可笑的。\par

\Emph{厄瓦尼斯特.}——我跟随宗徒，他称被钉的那位为\Quote{光荣的主}。\par

\Emph{正教徒.}——我也跟随，并且相信祂是\Quote{光荣的主}。因为被钉在木头上的身体不是任何普通人的，而是光荣的主的。但我们必须承认，结合使名号成为共通的。再说一次：你是否说主的肉身是从天降下的？\par

\Emph{厄瓦尼斯特.}——当然不是。\par

\Emph{正教徒.}——而是在童贞女的子宫中形成的？\par

\Emph{厄瓦尼斯特.}——是的。\par

\Emph{正教徒.}——那么，主如何说\Quote{如果你们看见人子升到祂以前所在之处}，以及\Quote{没有人升过天，除了那从天降下的人子，即是在天上的}？\par

\Emph{厄瓦尼斯特.}——祂不是在讲肉身，而是讲天主性。\par

\Emph{正教徒.}——是的，但天主性来自天主父。那么祂为什么称自己为人子？\par

\Emph{厄瓦尼斯特.}——本性的特有属性由位格共享，因为由于结合，同一位既是人子又是天主之子，既是永恒的又是有时间的，既是达味之子又是达味之主，其余以此类推。\par

\Emph{正教徒.}——很对。但认识到这个事实也很重要：两者共有一个名号并不会导致本性的混淆。因此，我们努力区分同一位怎样既是天主之子又是人子，以及祂怎样\Quote{昨天、今天、直到永远，都是一样的}，并且通过虔敬地区分术语，我们发现这些矛盾是一致的。\par

\Emph{厄瓦尼斯特.}——你说得对。\par

\Emph{正教徒.}——你说神圣本性从天降下，并因结合而被称为人子。因此，我们理应该说，肉身被钉在树上，但神圣本性即使在十字架和坟墓上，也与这肉身不可分离，虽然肉身感觉不到苦难，因为神圣本性自然地不能受苦和死亡，其本体是不朽和不受苦的。在这个意义上，被钉者被称为\Quote{光荣的主}，这是将不受苦本性的称号归于受苦的本性，因为我们知道，身体被描述为属于后者。\par

现在让我们就这样来审查此事。神圣宗徒的话是：\BibleQuote{他们若认识，就不会把光荣的主钉在十字架上了}。他们钉死他们所认识的本性，而非他们完全不知的本性；倘若他们知道那不知的，就不会钉死那所知的；他们钉死人性，因为不认识天主性。你们忘了他们自己的话吗？\BibleQuote{我们不是为了善事用石头砸你，而是为了亵渎，并且因为你是人，却自称为天主}。这些话清楚地证明，他们认出所见到的本性，而对那无形的却完全无知；倘若他们认识那本性，就不会把光荣的主钉死了。\par

\Emph{厄兰}——这非常可能，但教父们在尼西亚大公会议上订立的信仰说明称，独生子自己，真天主，与圣父同一实体，受难并被钉十字架。\par

\Emph{正统派}——你好似忘记了我们一再协议的事。\par

\Emph{厄兰}——你是什么意思？\par

\Emph{正统派}——我的意思是，在结合之后，圣经将高举与贬抑的术语都应用于同一位。但或许你也不知道，卓越的教父们先提到祂取了肉身并成为人，然后才加上祂受难并被钉十字架，并在设立了能受苦的本性之后才谈到受难。\par

\Emph{厄兰}——教父们说，天主之子，光之光，出自圣父的实体，受难并被钉十字架。\par

\Emph{正统派}——我不止一次观察到，神圣的和人性的都被归于同一位。正是基于此，三倍有福的教父们在教导我们应如何信仰圣父，然后转向圣子之位时，并没有立即加上\Quote{与天主之子}，虽然按常理，在界定关乎天主圣父的事之后，本应直接引入子的名号。但他们的目的是同时给予我们关于\OriginalTerm{神学}{theology}和\OriginalTerm{救恩计划}{œconomy}的教导，以免有人以为神性之位与人性之位有任何分别。为此，他们在有关圣父的声明中加上，我们也必须信我们的主耶稣基督，天主之子。现在，降生之后，天主圣言被称为基督，因为这名字同等地包含神性与人性的一切属性。然而我们承认，有些特性属于这一个本性，有些属于那一个，这可以立即从信经的实际措辞中看出。请告诉我：你把\Quote{出自圣父的实体}这句话应用于什么？是神性，还是那由达味的后裔所形成的本性？\par

\Emph{厄兰}——应用于神性，这是显然的。\par

\Emph{正统派}——而\Quote{出自真天主的真天主}这一句，你认为属于哪一个？是神性还是人性？\par

\Emph{厄兰}——属于神性。\par

\Emph{正统派}——因此，肉体和灵魂都不与圣父同一实体，因为它们是受造的，而那创造了万物的神性则不然。\par

\Emph{厄兰}——对。\par

\Emph{正统派}——很好，那么。当我们听到受难与十字架时，必须认出那顺从于受难的本性；我们必须避免将其归于不能受难者，而应归于那专为受苦而被取的本性。最卓越的教父们承认神性本性是不能受难的，而他们将受难归于肉体，这由信经的结尾所证明，其中写道：\Quote{但那些说'祂曾有一时不存在'，'受生以前不存在'，'祂从无中造成'，或声称天主之子是另一样本质或实体，可变的或可改变的，这些人，神圣公教会、宗徒教会予以绝罚。}那么，看哪，对那些将受难归于神性本性的人，宣布了何等刑罚！\par

\Emph{厄兰}——他们在这里谈论的是改变和变化。\par

\Emph{正统派}——但受难若不是改变和变化，又是什么呢？因为，如果祂在降生成人之前是不能受难的，在降生成人之后却受了难，那么祂必定是因经历了改变而受难；如果祂在成为人之前是不朽的，在成为人之后却尝了死味，如你所说，那么祂在成为人之后，经历了一场完全的变易，从有死变为不死。但这类说法及其作者，都被卓越的教父们逐出教会范围，并像腐烂的肢体一样从健康的身体上砍下。因此，我们劝你畏惧刑罚，憎恶亵渎。\par

现在我要向你指出，圣教父们在自己的著作中持守我们所表达的意见。在我要提出的见证人中，有些参与了那伟大的大公会议；有些在其后时代在教会中昌盛；有些则早在很久以前照亮了世界。但他们的和谐既不被时代的差异也不被语言的多样性所打破；如同竖琴，其琴弦多而分离，却奏出和谐的音乐。\par

\Emph{厄兰}——我渴望并乐于看到这样的引述。这类教导不可辩驳，且极为有益。\par

\Emph{正统派}——现在，敞开你的耳朵，接纳从属灵泉源涌出的溪流。\par

\Signature{圣依纳爵，安提约基亚主教，殉道者的见证}\par

\WorkTitle{出自他致斯弥尔纳人书}：\par

\Quote{他们不接纳感恩祭和奉献，因为他们不承认感恩祭是救主耶稣基督的肉体，这肉体为我们的罪而受难，并因祂的慈善而被圣父复活。}\par

\Signature{依勒内，里昂主教的见证}\par

\WorkTitle{出自他反驳异端第三卷（第二十章）}：\par

\Quote{显然，\Person{保禄}所知道的基督，唯有那受苦、被埋葬、复活并被生的那一位，他称之为人。因为说过\BibleQuote{如果基督被宣讲说祂从死者中复活了}之后，他接着说明祂降生成人的理由：\BibleQuote{因为死亡既由一人而来，死者的复活也由一人而来}，并且在所有涉及主的受难、人性和瓦解的场合，他都使用基督这个名字，如经上说：\BibleQuote{不可因你的食物使基督为他而死的那人败坏}，又说：\BibleQuote{但如今在基督内，你们从前远离的人，藉着基督的血，成了亲近的}，又说：\BibleQuote{基督由法律的咒骂中赎出了我们，为我们成了可咒骂的，因为经上记载：凡挂在木头上的，都是被咒骂的。}}\par

同一位作者，出自同一著作（第二十一章）：\par

\Quote{因为祂曾是人，为受试探；同样，祂曾是圣言，为受光荣。在祂受试探、被钉死及死亡中，圣言无所行动；但在祂的胜利、忍耐、良善、复活和被提升天中，圣言与人性合作。}\par

同一位作者，出自同一著作第五卷：\par

\Quote{主以自己的血赎回了我们，并以祂的灵魂代替我们的灵魂，以祂的肉代替我们的肉。}\par

\Emph{圣\Person{希坡律陀}主教兼殉道者的证言}\par

出自他写给某位皇后的信函：\par

\Quote{因此祂称祂为\BibleQuote{已眠者的初果}和\BibleQuote{死者中的首生者}。当祂复活后，愿意显示那复活的是与死亡相同的身体时，宗徒们怀疑，祂就召叫\Person{多默}来说：\BibleQuote{摸摸我，看看！因为鬼魂没有血肉，如同你们看到我所有的。}}\par

同一位作者，出自同一信函：\par

\Quote{祂称祂为初果，是为证明我们所说的：救主取了相同质料的肉身后，使之复活，使这肉成为义人肉身的初果，使我们所有相信的人，因着对那位复活者的信赖，期望我们的复活。}\par

同一位作者，出自他关于两个盗贼的讲道：\par

\Quote{主的身体把两样东西给了世界——圣血和圣水。}\par

同一位作者，出自同一讲道：\par

\Quote{按人的说法，这身体虽是一具尸体，其本身却蕴藏巨大的生命力，因为从它流出了从死尸中不会流出的东西——血和水——使我们知道那内住在身体中的能力蕴含着何等强大的生命力，以至于这一具尸体显然与其他尸体不同，且能为我们倾流出生命的泉源。}\par

同一位作者，出自同一讲道：\par

\Quote{圣羔羊的骨头一根也不可折断。预象表明，受难不能触及那能力，因为骨头是身体的能力所在。}\par

\Emph{圣\Person{欧斯大丢}，\Place{安提约基雅}主教兼精修者的证言}\par

出自他关于灵魂的著作：\par

\Quote{他们不虔敬的毁谤用几句话就能反驳；他们也许是对的，除非祂为了人类的救恩自愿将自己的身体交给死亡的毁坏。首先，他们归咎于祂一种异常的软弱，即不能击退敌人的攻击。}\par

同一位作者，出自同一著作：\par

\Quote{为何他们在炮制那些尘世骗局时，竭力证明基督取了一个没有灵魂的身体？为的是，如果他们能诱使某人接受这个说法，那么，通过将情感变化归给天主圣神，他们就能轻易说服那人，认为可变者不是出自不变者的本性。}\par

同一位作者，出自他关于\BibleQuote{主在祂道路的起点创造了我}的讲道：\par

\Quote{那死去的人第三日复活了，当\Person{玛利亚}急于触摸祂的圣肢时，祂制止并喊道：\BibleQuote{不要摸我，因为我还没有升到父那里去；但你去到我的弟兄那里，对他们说：我要升到我的父和你们的父那里，到我的天主和你们的天主那里。}现在，\BibleQuote{我还没有升到父那里去}这句话，不是由那从天上降下来、曾在父怀里的圣言和天主所说的，也不是由那包含一切受造物的智慧所说的，而是由那由各种肢体组成的人说的，这人从死者中复活，死后还未回到父那里，并为自己的进步保留了初果。}\par

同一位作者，出自同一著作：\par

\Quote{他写的时候，明确地将被钉十字架的人描述为光荣的主，宣认祂是主和基督，正如宗徒们同声对血肉之躯的以色列人说：\BibleQuote{因此，以色列全家应确知，天主已将你们钉死的这位耶稣立为主和基督了。}祂如此形成了受苦的耶稣基督。祂没有如此形成智慧，也没有如此形成那从起初就有统治权能的圣言，而是形成了那被高高举起、在十字架上伸出双手的那一位。}\par

同一位作者，出自同一著作：\par

\Quote{因为如果祂是无形的，不被手动接触，也不被肉眼看见，祂就不会受伤，不会被钉钉子，不会分享死亡，不会被埋在地下，不会被封在坟墓里，不会从坟墓中复活。}\par

同一位作者，出自同一著作：\par

\Quote{\BibleQuote{没有人夺去我的性命，是我甘心情愿舍掉它；我有权舍掉它，也有权再取回它。}如果作为天主，祂拥有双重权能，祂却向那些怀着恶意计划毁灭圣殿的人让步，但藉着复活，祂以更大的光荣重建了圣殿。无可辩驳的证据表明，祂自己复活并更新了自己的殿宇，圣子的大工应归于天主圣父；因为圣子不脱离父行事，正如圣经无可指责的声明所宣告的。因此，有时圣父被描述为使基督从死者中复活，有时圣子应许要重建自己的圣殿。那么，如果根据之前所确立的，基督的神性被证明是不可受苦的，那么那些被诅咒的人攻击宗徒的定义就是徒劳的。如果\Person{保禄}说光荣的主被钉十字架，显然是针对人性说的，我们就不应因此将受苦归于神性。那么他们为何将这两件事放在一起，说基督因软弱而被钉十字架呢？}\par

同一位作者，出自同一著作：—\par

\Quote{但若将任何软弱归给祂是合宜的，任何人都会说，将这些性质归于人性是自然的，尽管不归于天主性的圆满，也不归于至高智慧的尊荣，也不归于那按\Person{保禄}所描述的万有之上的天主。}\par

同一位作者，出自同一书：—\par

\Quote{这就是软弱的方式，按此方式\Person{保禄}描述祂走向死亡，因为当人明确地与天主的圣神结合时，人靠天主的德能生活，因为由前面的论述可知，那被认为在祂内的那位也被证明是至高者的德能。}\par

同一位作者，出自同一著作：—\par

\Quote{正如进入童贞玛利亚的胎并未减少祂的德能，同样，将祂的身体钉在十字架的木头上也未玷污祂的神。因为当身体被高举在十字架上时，神圣的智慧之神甚至住在身体内，踏足于天上的地方，充满全地，统治深渊，眷顾并审判每个人的灵魂，继续做天主不断做的一切，因为那高处的智慧并不被囚禁在物质身体内，如同湿的和干的材料被封闭在器皿内，被包含却不包含它们。但这智慧是神圣而不可言喻的德能，同样拥抱并坚固殿内和殿外的一切，由此超越，一举涵盖并掌管所有物质。}\par

同一位作者，出自同一著作：—\par

\Quote{但若太阳作为可见的物体，被感官所感知，到处忍受如此不利的影响，不变其秩序，也不觉任何打击，无论大小；我们岂能设想非物质的智慧被玷污并改变其本性，因为它的殿被钉在十字架上，或被毁坏，或被创伤，或被腐蚀？殿受苦，但本质毫无玷污而存留，并保持其完整的尊严。}\par

同一位作者，出自其\WorkTitle{登阶咏标题论}：—\par

\Quote{圣父是完美的、无限的、不可理解的，既不能增饰也不能减损，不接受获得的荣耀；同样，祂的圣言——由祂所生的天主——也不接受，藉着祂有天使、天和地的无边量体以及受造万物的所有形式与物质；但从死人中复活的基督耶稣被高举得荣耀，使祂的仇敌公然溃败。}\par

同一位作者，出自同一著作：—\par

\Quote{但那些对祂怀有仇恨的人，即使他们被祂仇敌的军力所包围，也被驱散到四方，而天主圣言荣耀地使祂自己的殿复活。}\par

同一位作者，出自其\WorkTitle{第九十二篇圣咏释义}：—\par

\Quote{此外，先知\Person{依撒意亚}跟随祂受苦的轨迹，在其他言论中大声呼喊：\InnerQuote{我们看见祂，祂没有形态，也没有美丽。祂的形态被羞辱、被众人弃绝。}以此清楚表明，耻辱和痛苦必须归于人性而非神性。紧接着他又加上：\InnerQuote{祂是受打击的人，能承受软弱。}正是这位，在受辱之后，被看见没有形态与美丽，然后又改变并披上美丽，因为住在祂内的天主并未像羔羊被牵到宰杀之地，也未像绵羊被屠宰，因为祂的本性是不可见的。}\par

\Emph{圣亚大纳削的见证，亚历山大主教，精修者。}\par

从他的\WorkTitle{致埃皮克提图书}：—\par

\Quote{谁竟堕落到如此不敬的地步，以至认为并说出那与圣父同体的天主性本身受了割损，并且从完美变成不完美；并否认被钉在十字架上的是身体，反而断言那是智慧的创造本质？}\par

同一位作者，出自同一论文：—\par

\Quote{圣言亲自关联并承受了圣言的人性所受的苦，以便我们能分享圣言的天主性。奇妙的是，受苦者与不受苦者是同一位；作为受苦者，因为祂自己的身体受苦，而祂就在那身体内受苦，但祂又不受苦，因为圣言按本性是天主，是不能受苦的。而祂自己这非物质的，就在可受苦的身体内，并且身体内包含着不能受苦的圣言，摧毁了祂身体的各种软弱。}\par

同一位作者，出自同一书信：—\par

\Quote{因为祂是天主和荣耀之主，却以不荣耀的方式在身体内被钉十字架；但身体在树上被击打时受苦，水与血从祂的肋旁流出；但因它是圣言的殿，便充满天主性。因此当太阳看见它的造物主在受辱的身体内受苦时，便收回了光芒，使大地黑暗。而那个有朽坏本性的身体，因着其中的圣言，超越了自身的本性，不再受其本性的朽坏所触及，反而穿上了超人的圣言，成为不朽的。}\par

同一位作者，出自其\WorkTitle{论信仰大讲章}：—\par

\Quote{从死者中复活的，是人还是天主？宗徒\Person{伯多禄}，比我们更了解，解释并说：\InnerQuote{他们既完成了论及祂的全部记载，就从木架上取下祂，放在墓穴里，但天主使祂从死者中复活了。}从木架上取下的耶稣的尸身，曾被放在墓穴里，由\Person{阿黎玛特雅人}若瑟埋葬，正是圣言所复活的身体，祂说：\InnerQuote{你们拆毁这座圣殿，三天之内，我要把它重建起来。}祂是使所有死者复活的，也使那人耶稣基督——由\Person{玛利亚}所生、祂所取的人性——复活了。因为若祂在十字架上使先前已分解的圣人的尸体复活，那么永生天主的圣言更能使祂所穿着的身体复活，如\Person{保禄}所说：\InnerQuote{因为天主的话是活泼的，是有力的。}}\par

同一位作者，出自同一著作：—\par

\Quote{生命不会死，却使死者复活；因为正如光在黑暗处不受损害，生命在探访有死本性时也不能受苦，因为圣言的天主性是不变不易的，正如主在关于祂自己的预言中所说：\InnerQuote{我是上主，我不改变。}}\par

同一位作者，出自同一著作：—\par

\Quote{生活的祂不能死，反而使死者复活。因此，祂因由圣父而来的\OriginalTerm{天主性}{Godhead}，是光的泉源；但那死了的，或更好说由死者中复活的，我们的中保，由童贞玛利亚所生，\OriginalTerm{圣言}{Word}的\OriginalTerm{天主性}{Godhead}为我们而取了祂，祂是人。}\par

同一位作者，同一著作：—\par

\Quote{事情是这样的：拉匝禄患病并死了；但那位神人没有患病，也没有违背自己的意愿而死，而是自愿来到\OriginalTerm{死亡的救恩计划}{dispensation of death}，由内住于祂内的天主圣言所加强，祂说：\BibleQuote{没有人夺去我的生命，而是我自愿舍弃。我有权舍弃，也有权再取回。}因此，那舍弃并取回祂所穿戴的人性生命的，是圣子的\OriginalTerm{天主性}{Godhead}，因为祂完整地取了\OriginalTerm{人性}{manhood}，为能完整地使之复活，并连同它使死者复活。}\par

同一位作者，摘自其驳斥亚略派的讲道：—\par

\Quote{因此，当蒙福的保罗说圣父\InnerQuote{复活了}圣子\InnerQuote{从死者中}时，若望告诉我们耶稣说：\BibleQuote{拆毁这座圣殿，三天内我要把它重建起来……但祂所说的是}祂自己的\InnerQuote{身体。}因此，对于留心的人而言，很清楚：在身体的复活上，保罗说圣子是从死者中复活了，因为他将属于身体的事归于圣子的人位；同样，当他说\InnerQuote{圣父赐生命给圣子}时，必须理解为生命是赐给了\OriginalTerm{肉身}{Flesh}。因为如果祂本身是生命，生命如何能领受生命呢？}\par

同一位作者，摘自其论降生成人之著作：—\par

\Quote{因为当\OriginalTerm{圣言}{Word}意识到，若不藉由彻底的死亡，人类的败坏就无法解除，而圣言既是永生且是圣父之子，不可能死亡，为达成此目的，祂取了可死的身体，使这身体与统御万有的圣言结合，能为众人而成为死亡的臣属，并因内住的圣言而保持\OriginalTerm{不朽的}{incorruptible}，如此，藉着复活的恩宠，腐败将来便失去对人类的权势。因此，祂将所取的身体作为无瑕的\OriginalTerm{祭献}{sacrifice}和\OriginalTerm{牺牲}{victim}献给死亡，藉此相应的祭献，祂立即摧毁了死亡对祂所有同类的权势；因为祂是超越并凌驾于所有人之上的天主圣言，祂以正义献上并偿付了自己的圣殿和身体工具，作为所有该死亡灵魂的\OriginalTerm{赎价}{ransom}。如此，藉着相似（的身体）与所有人结合，不朽的天主圣子藉着复活的应许以正义给所有人穿上\OriginalTerm{不朽}{incorruption}，因为死亡内固有的腐败不再存在于人中，这是由于那藉着同一身体而住于他们内的圣言之故。}\par

同一位作者，同一著作：—\par

\Quote{因此，在祂大能的圣迹显示之后，现在祂也为所有人献上了祭献，将自己的圣殿代替众人交给死亡，为使众人对远古的过犯不负责任并得自由，并藉着显示自己的身体为人类复活的\OriginalTerm{不朽的}{incorruptible}\OriginalTerm{初果}{firstfruits}，表明自己比死亡更强。因为那身体，由于具有共同的实体——因为它是人的身体，尽管藉着新的奇迹，它完全由童贞女构成——因此是必死的，按同类之例而死亡；但由于\OriginalTerm{圣言}{Word}降入其中，它不再按其本性遭受腐败，反而因内住于它的天主圣言，得以脱离腐败。}\par

同一位作者，同一著作：—\par

\Quote{因此，如我所说，既然\OriginalTerm{圣言}{Word}是不朽的，不可能死亡，祂便取了可死的身体，为能将同一身体为众人献上，而祂自己在为众人的苦难中，藉着降入这身体，\BibleQuote{毁灭那握有死亡权势的。}}\par

同一位作者，同一著作：—\par

\Quote{因为那身体在受难时，如同身体的本性，死了，但它藉着内住在它内的\OriginalTerm{圣言}{Word}，获得了不朽的应许。因为当身体死亡时，圣言并未受损；但祂本身是无所受难、不朽、且永生的，作为天主的圣言，且因与身体结合，祂保住了它免于身体自然的\OriginalTerm{腐败}{corruption}，正如圣神向祂所说：\BibleQuote{祢不会让祢的圣者见到腐朽。}}\par

\Emph{圣\Person{达玛稣}，\Place{罗马}主教的见证}：—\par

\Quote{若有人说，在十字架的受难中，天主圣子经历了痛苦，而非那取了奴仆形像的、祂所穿上的肉身与灵魂，正如圣经所说，那么，让他受诅咒。}\par

\Emph{圣\Person{盎博罗削}，\Place{米兰}主教的见证。}\par

摘自其论公教信仰之书：—\par

\Quote{有一些人达到了如此不敬的程度，竟认为主的\OriginalTerm{天主性}{Godhead}受了割损，从完美变成了不完美；并且认为那创造万有的\OriginalTerm{神圣本体}{divine substance}，而非肉身，被悬在木头上。}\par

同一位作者，同一著作：—\par

\Quote{肉身受了苦；但\OriginalTerm{天主性}{Godhead}是免于死亡的。祂按人性的律令将自己的身体交出，使之受苦。因为当灵魂不能死时，天主如何能死？祂说：\BibleQuote{不要害怕那些杀害身体而不能杀害灵魂的。}如果灵魂不能被杀害，那么天主性如何能成为死亡的臣属呢？}\par

\Emph{圣\Person{巴西略}，\Place{凯撒勒雅}主教的见证：—}\par

\Quote{但凡对宗徒话语的含义稍有了解的人都知道，他并非在教导我们一种神学模式，而是在解释\OriginalTerm{救恩计划}{œconomy}的理由，因为他说：\BibleQuote{天主已经立了你们钉死的这位耶稣为主，为基督。}因此，他明确地将论据指向祂的人性及可见的本性。}\par

\Emph{圣\Person{额我略}，\Place{纳齐盎}主教的见证。}\par

摘自其致蒙福的\Person{乃克达留}，\Place{君士坦丁堡}主教的信函：—\par

\Quote{在教会所遭遇的事情中，最可悲的是\Person{阿波利纳里}及其同伙的言论之放肆。我无法理解您的圣座如何允许他们篡夺与我们同等集会的权力。}\par

稍后又说：—\par

\Quote{我不再称此为严重；这实在是最可悲的，即独生子天主本身，一切存在者的审判者，生命之君，死亡之毁灭者，竟被他说成会死，并声称在其自己的天主性中承受苦难。他使天主性与身体在那三天死亡的身体分解中有分，并如此在死亡之后又由圣父复活。}\par

同一位作者，出自其先前致\Person{克莱多尼奥}的论述：—\par

\Quote{亚略派争论说，那人性是没有灵魂的，为要把苦难归于天主性，并声称同一能力既移动身体又受苦。}\par

同一位作者，出自其关于圣子的讲道：—\par

\Quote{我们本应处理祂所被命令之事，以及祂遵守诫命、做一切中悦祂的事；此外还有祂的成全、高举、因所受的苦难学习服从、祂的司祭职、祂的奉献、祂的被出卖、祂向那能救祂免于死亡者的恳求、祂的忧苦、祂的血汗、祂的祈祷及类似的表现，若非人人清楚，所有这些与祂受难相关的表达绝不意指那不变且超越苦难的本性。}\par

同一位作者，出自其复活节讲道（第二篇）：—\par

\Quote{\InnerQuote{那从\Place{厄东}而来的是谁？}从地上而来，那无血无体者的衣裳怎能红如榨酒池中践踏者？请以受苦的身体之美作答，它在苦难中变得美丽，并被天主性所荣耀，无物比之更可爱、更美。}\par

\Emph{\Place{尼撒}的\Person{额我略}主教的证言。}\par

出自其教理演讲：—\par

\Quote{这就是天主关于人性及从死者中复活的救恩计划的奥迹：不是阻止灵魂按照人性必然的法则因死亡与身体分离，而是通过复活使之重新合一。}\par

同一位作者，出自同一著作：—\par

\Quote{那接受了天主性、并藉复活与天主性一同被高举的血肉，并非由其它材料形成，而是由我们的材料形成。因此，如同我们自己的身体一样，一个感官的活动将整个与该部分相连的人引至全身感觉，同样，仿佛整个自然界是一个单一活物，部分的复活渗透全体，通过本性的连续与合一从部分传递到全体。那么，在这奥迹中，有何异常之处：直立者屈身于堕落者，为扶起那低垂者？}\par

同一位作者，出自同一著作：—\par

\Quote{在此部分，自然也不可只注意其一而忽略另一；而是在不朽者中看到人性，并对那人性中的神圣品质加以精细考察。}\par

同一位作者，出自其驳斥欧诺弥的著作：—\par

\Quote{并非人性使\Person{拉匝禄}复活；并非不能受苦的能力为死者流泪。眼泪属于人；生命来自生命本身。数千人并非由人的贫困所饱饫；全能并未匆匆走向无花果树。谁在途中疲倦，而谁用祂的话支撑全世界而不疲倦？祂荣耀的光辉是什么，什么被钉子刺穿？什么形体在受难中被击打，什么被永远荣耀？答案显而易见，无需解释。}\par

同一位作者，出自同一论著：—\par

\Quote{他责备那些将苦难归于人性者。他自己愿意完全使天主性本身服从苦难，因为命题是双重的且可疑的：神性或人性是否涉及苦难，对其中一方的否定就成为对另一方的肯定指控。因此，当他们责备那些在人性中看见苦难者时，却会对那些主张天主之子可受苦者给予无保留的赞扬。但由此确立的观点却成为他们教义荒唐性的确证；因为，如他们所说，圣子的天主性受苦，而圣父的天主性按其本质保持完全不能受苦，那么，不能受苦的本性与承受苦难的本性就是相矛盾的。}\par

\Emph{圣\Place{依科尼雍}的主教\Person{安斐洛西}的证言。}\par

出自其关于经文\Quote{我实实在在告诉你们：那听我的话，并相信派遣我者，就有永生}的讲道：—\par

\Quote{那么，苦难是谁的？是血肉的。因此，你若将苦难归于血肉，也要将卑微的话归于它；而将崇高的话归于你所赋予奇迹的那一位。因为天主在行奇事时，自然会以与其事迹相称的高超伟大之言说话；而人在受苦时，则适切地发出与苦难相称的卑微之言。}\par

同一位作者，出自其关于\Quote{我的父比我大}的讲道：—\par

\Quote{但当你将苦难归于血肉，将奇迹归于天主时，你必然，虽不情愿，也要将卑微的话归于那由\Person{玛利亚}所生的人，而将崇高伟大、合于天主的话，归于那起初就有的圣言。我有时说崇高的话，有时说卑微的话，原因是：借崇高的话，我显明内住的圣言的尊贵；借卑微的话，我表明卑微血肉的软弱。因此，我一时称自己与父相等，一时又称父比我大；我在此并不自相矛盾，而是表明我是天主又是人：藉崇高的话是天主，藉卑微的话是人。你若想知道我的父在哪方面比我大，我是在血肉中说话，而不是以天主性的位格说话。}\par

同一位作者，出自其关于\Quote{若是可能，就让这杯离开我}的讲道：—\par

\Quote{不要把肉身的苦难归给不受苦的天主，因为我是天主也是人，异端者；借着奇迹显为天主，借着苦难显为人。既然我是天主也是人，告诉我，谁受了苦？如果天主受苦，你说了亵渎的话；但如果肉身受苦，你为什么把苦难归于你所畏惧的那一位？因为一位受苦时，另一位并不感到畏惧；人受十字架时，天主并不困扰。}\par

同一位，在其驳斥亚略人的讲道中：——\par

\Quote{为免我所说的冗长，我将简短地问你，异端者，在万世之前由天主所生的那一位受苦了，还是末日在达味所生的耶稣受苦？如果天主性受苦，你说了亵渎的话；如果正如真理所是，人性受苦，你为何犹豫把苦难归于人？}\par

同一位，在其论及圣子的讲道中：——\par

\Quote{\Person{伯多禄}说：\BibleQuote{天主已经立这位耶稣为主，为基督了？}又说：\BibleQuote{这位你们钉死的耶稣，天主已经复活了。}如今成为尸体的不是天主性，而是人性；使之复活的是圣言，天主的德能，祂在福音中说：\BibleQuote{你们拆毁这座圣殿，三天之内我要把它重建起来。}因此，当圣经说天主使成为尸体又从死者复活的那一位为主为基督时，所指的乃是肉身，而非圣子的天主性。}\par

同一位，在其论及\BibleQuote{子不能凭自己做什么}的讲道中：——\par

\Quote{因为祂的本性并不是那样，使祂的生命能被腐朽所拘禁，因为祂的天主性并非被强制地陷入苦难。如何能呢？但人性在朽坏中得以更新。因此祂说：\BibleQuote{因为这必朽坏的必须穿上不朽，这必朽的必须穿上不朽。}你注意这精确性；祂清楚地指明\Quote{这必朽的}，好使你不致想到任何其他肉身的复活。}\par

\Emph{圣\Person{弗拉维亚诺}的见证，\Place{安提约基雅}主教。}\par

在复活节：——\par

\Quote{因此，我们也勇敢地宣讲十字架，并在我们中间承认主的死亡，虽然天主性丝毫没有受苦，因为天主是不受苦的，但救恩计划借着身体得以完成。}\par

同一位，论及叛徒犹达斯：——\par

\Quote{因此，当你听到主被出卖时，不要将天主的尊严贬低为卑微，也不要把身体的苦难归于天主的权能。因为天主是不受苦的、不变的。如果因祂对人的爱，祂取了奴仆的形象，祂的本性并未改变。祂仍是祂向来所是的，将神圣的身体交付给死亡的经验。}\par

\Emph{\Person{德奥斐禄}的见证，\Place{亚历山大里亚}主教。}\par

摘自其节期卷：——\par

\Quote{无理之生灵的灵魂不被取走也不被放回：它们与身体同腐朽，化为尘土。但救主在十字架之时，从祂自己的身体中取了灵魂，当祂从死者中复活时，又将灵魂归还给身体。为向我们保证这一点，祂说出了圣咏作者的话，那预言的呼喊：\BibleQuote{祢不会将我的灵魂遗留在阴间，也不让祢的圣者见腐朽。}}\par

\Emph{真福\Person{杰拉修斯}的见证，\Place{巴勒斯坦}的\Place{凯撒勒雅}主教：——}\par

\Quote{祂被捆绑，祂被击伤，祂被钉十字架，祂被触摸，祂被留下伤疤，祂被长矛刺伤，所有这些侮辱都由\Place{玛利亚}所生的身体承受，而从圣父在万世之前所生的那一位，无人能够伤害，因为圣言没有这样的本性。因为谁能约束天主性？如何能击伤它？如何使无形体的本性染红血液？如何用坟墓的布条围绕它？现在承认你无法反驳的，并被无法抗拒的理性所约束，尊崇天主性。}\par

\Emph{圣\Person{若望}的见证，\Place{君士坦丁堡}主教。}\par

摘自其论及\BibleQuote{我父作工直到现在，我也作工}的讲道：——\par

\Quote{\BibleQuote{祢显什么神迹给我们看，因为祢做这些事？}那么祂自己回答什么？\BibleQuote{你们拆毁这座圣殿，三天之内我要把它重建起来。}祂说的是祂自己的身体，但他们不明白。}\par

稍后又说：——\par

\Quote{为什么福音作者不略过这句话？为什么他加上纠正：\InnerQuote{但祂说的是祂身体的圣殿}？因为祂没有说\Quote{拆毁这个身体}，而是\Quote{圣殿}，为要表明内住的天主。拆毁这座圣殿，它远比犹太人的圣殿卓越。犹太人的圣殿包含法律；这座圣殿包含立法者；前者是杀死的文字；后者是使人生存的圣神。}\par

同一位，摘自讲道\Quote{那在谦卑中所说所行的，并非因为权能软弱，而是不同的救恩计划}：——\par

\Quote{那么祂怎么说\BibleQuote{如果可能}？祂是在向我们指明人性的软弱，这人性不愿脱离今世的生命，反而因天主起初植入的对此生的爱而退后和退缩。如果当主自己如此多次这样说时，还有人胆敢说祂没有取了肉身，那么如果祂没有说过这些话，他们会说什么呢？}\par

同一位，摘自同一著作：——\par

\Quote{注意他们如何谈论祂早年的年龄。问异端者：天主会害怕吗？祂会退后吗？祂会畏缩吗？祂会忧伤吗？如果他说是，从此远离他，把他与魔鬼同列，甚至比魔鬼更低，因为魔鬼也不敢说这话。但如果他说每一件事都不配于天主，回答——天主也不祈祷；因为除此之外，如果这些话是天主的话，将是另一个荒谬，因为这些话不仅表明一种苦闷，而且表明两个意志；一个是圣子的意志，另一个是圣父的意志，彼此对立。因为\BibleQuote{不要照我的意愿，而要照祢的意愿}这话，正是表明了这一点。}\par

同一位，摘自同一著作：——\par

\Quote{因为如果这是就\OriginalTerm{天主性}{Godhead}而言，便会出现某种矛盾，由此产生许多荒谬之处。相反，如果是就肉体而言，这些表达是合理的，无可指摘。因为肉体不愿死亡并不招致谴责；肉体本性如此，祂彰显肉体的一切属性，唯除罪过，而且确实极其充分，为要堵住异端者的口。因此，当祂说\InnerQuote{若是可能，就让这杯离开我}和\InnerQuote{不要照我的意愿，而照你的意愿}，祂只是表明祂确实穿上了那惧怕死亡的肉体，因为惧怕死亡、退缩和受苦是肉体的本性。祂让肉体被遗弃、剥去其自身活动，为要藉展示其软弱也让我们确信其本性。然而，有时祂又隐藏它，因为祂不仅是人。}\par

\Emph{\Person{塞维里亚努斯}之证言，\Place{加巴拉}主教。}\par

摘自其论封印的讲辞：——\par

\Quote{犹太人抗拒可见的，不知不可见的；他们钉死肉体；他们并不毁灭天主性。因为倘若我的言语连同作为言语外衣的文字不被毁灭，那么天主圣言、生命之源又怎能与肉体一同死亡？受苦属于身体，而无受苦性却属于尊位。}\par

那么请看，那些在东方、西方、南方和北方耕种的人，都已被我们表明谴责你们虚妄的异端，并公开宣告神圣本性的无受苦性。请看，两种语言——我指希腊语和拉丁语——关于天主之事作出了何等和谐的宣认。\par

\Emph{赝辩者：}——我本人对他们的和谐感到惊讶，但我观察到他们所用术语存在显著差异。\par

\Emph{正统者：}——请勿动怒。正是他们与对手抗争的激烈本身，导致了他们看似过激。这在种植者身上也可观察到；当他们看到一棵树向某侧弯曲时，他们不会满足于将其拉直，反而会将其向相反方向进一步弯曲，为的是通过这样过度弯曲而使其恢复笔直姿态。为使你明白，正是这多种异端的推动者和支持者，试图以其亵渎之巨超越往昔异端，请再听\Person{阿波利纳里}的著作，它们宣告神圣本性的无受苦性，并承认受苦属于身体。\par

\Emph{\Person{阿波利纳里}之证言。}\par

摘自其纲要：——\par

\Quote{若望论及那被毁的圣殿，即那复活圣殿者的身体，而身体与祂完全结合，祂在他们中并非另一位。如果主的身体与主是一体，身体的属性便因身体而成为祂的属性。}\par

又：——\par

\Quote{事实上，祂与身体的结合并非通过圣言受限制而发生，以致祂除道成肉身之外一无所有。因此，即使在死亡中，不朽与祂同在；因为如果祂超越这一结合，那么祂也超越解散。而死亡是解散。但祂并不包含在结合中；若如此，宇宙便会虚空；在解散中，祂也不像灵魂那样遭受解散后的丧失。}\par

又：——\par

\Quote{正如救主说死者的身体从坟墓出来，而他们的灵魂并非从那里出来，同样地，祂说自己将要从死者中复活，尽管复活的只是祂的身体。}\par

在另一部类似著作中，他写道：——\par

\Quote{从死者中复活属于人；使之复活属于天主。而基督既复活又使之复活，因为祂是天主和人。倘若基督只是人，祂就不会使死者复生；倘若祂只是天主，祂就不会独自、离开圣父使任何死者复生。但基督两者都做了；同一位存有既是天主又是人。倘若基督只是人，祂就不会拯救世界；倘若祂只是天主，祂就不会藉受苦拯救世界；但基督两者都做了，因此祂是天主和人。倘若基督只是人，或只是天主，祂就不能作人与天主之间的中保。}\par

稍后又：——\par

\Quote{而今，肉体是生命的工具，按天主旨意适合承受苦难。言语不属于肉体，行为也不属于。由于肉体按其本性被置于承受苦难的能力之下，它因是天主的肉体而胜过苦难。}\par

又稍后：——\par

\Quote{圣子从童贞女取了肉体，来到世界。祂以圣神充满这肉体，以圣化我们众人。祂就这样将死亡交给死亡，并藉复活摧毁死亡，以复活我们众人。}\par

摘自其论信德的短论：——\par

\Quote{既然苦难涉及肉体，祂的大能自有其无受苦性，因此将苦难归于大能是亵渎的谬误。}\par

在其论道成肉身的短论中，他进一步写道：——\par

\Quote{在此祂表明，从死者中复活的那同一位人，即是统御万有的天主。}\par

你现在看到，虚妄异端的教授之一公开宣讲天主性的无受苦性，称身体为圣殿，并坚持认为这身体是由天主圣言所复活。\par

\Emph{赝辩者：}——我听了并感到惊讶；我实在羞愧，我们的教义竟似乎不如\Person{阿波利纳里}的创新牢靠。\par

\Emph{正统者：}——但我还要从另一异端群体为你带来一位证人，他明确宣讲独生者天主性的无受苦性。\par

\Emph{赝辩者：}——你指的是谁？\par

\Emph{正统者：}——你可能听说过腓尼基人\Person{欧瑟伯}，他曾是\Place{黎巴嫩}旁的\Place{艾梅沙}的主教。\par

\Emph{赝辩者：}——我读过他的一些著作，发现他是亚略学说的支持者。\par

\Emph{正统者：}——是的；他确实属于那一派，但在试图证明圣父大于独生者时，他宣称被贬低的圣子的天主性是无受苦的，并极其持久地为此意见争辩。\par

\Emph{赝辩者：}——如果您也能引用他的话，我将不胜感激。\par

\Emph{Orth.}——为满足你的愿望，我将引述稍长一些的证据。现在请听他说什么，设想那人亲自在向我们讲话。\par

\Emph{\Person{厄梅萨的欧瑟伯}之见证：}——\par

为何祂怕死？怕死会带给祂什么？死亡对祂是什么？岂非能力与肉身的分离？能力受了钉子吗，以致它该怕？若我们的灵魂与肉身相连时并不承受肉身的软弱，眼目虽盲，理智仍强；足虽被砍，推理能力并不跛行——这本性为此作证，主也以\Quote{不要怕那杀害身体却不能杀害灵魂的}这番话加以确认（他们既不能杀害灵魂，并非因为他们不愿，而是因为他们不能，尽管他们愿意使灵魂与同轭的肉身分担痛苦）——那创造灵魂且形成肉身的主宰，虽承担了肉身的痛苦，岂会如肉身般受苦？但基督为我们受苦，我们不说谎。\BibleQuote{我将赐给的食粮就是我的肉。}这是祂为我们所赐的。\par

\Quote{那能被制伏的已被制伏；那能被钉死的已被钉死，但那能同样居住其中又离开它的能力却说：\BibleQuote{父啊，我把我的灵魂交托在你手中}，不是交在那些试图加速祂死亡的人手中。我不喜欢争辩；我宁愿避免；我愿意以全然的温柔，如兄弟般探究我们之间的分歧。我说能力不能受肉身之苦，这难道不真确吗？我什么都不说；谁愿意，就让他说能力受了什么苦。它衰竭了吗？请看危险。它消灭了吗？请看亵渎。它不再存在了吗？这是能力的死亡。告诉我，什么能如此制伏它，以致它受苦，我就退下。但如果你不能告诉我，你为何反对我不告诉你？你不能告诉我的，就是它没有接受。把钉子钉进灵魂里，我就承认它能被钉进能力里。但它是同情的。告诉我，你所谓的\InnerQuote{同情}是什么意思。当钉子进入肉身，痛苦也进入能力。让我们这样理解\InnerQuote{同情}。那么，未被击打的能力感到了痛苦。因为痛苦总是随着受苦而来。但如果身体常常因思想的力量而藐视痛苦，当心智健全时，那么在这种情况下，让某位公正地解释什么是受苦，什么是同情或一起受苦。那又怎样？基督没有为我们死吗？祂怎么死的？\BibleQuote{父啊，我把我的灵魂交托在你手中。}灵魂离开了；身体留下了；身体没有气息。那么祂没有死吗？祂为我们死了。牧人献出了羊，司铎献出了祭品，祂为我们舍了自己。\BibleQuote{祂没有怜惜自己的儿子，反而为我们众人把祂交出来。}我不拒绝这些话，但我想要这些话的含义。主说天主的食粮从天降下，虽然我因奥秘的缘故不能更清楚地表达，但祂解释道：\BibleQuote{这是我的肉。}圣子的肉是从天降下的吗？不。那么祂为何说，而且是解释，天主的食粮是活着的，并且从天降下？祂把能力的属性归于肉身，因为那取了肉身的能力从天降下。那么换词吧：祂把肉身所受的归于能力。基督如何为我们受苦？祂被吐唾沫，祂被掌掴，他们给祂戴上了荆棘冠冕，祂的手脚被刺穿。所有这些痛苦都是身体的，但它们被归给居住在其中者。向皇帝的雕像扔石头。喊叫是什么？\InnerQuote{你侮辱了皇帝。}撕破皇帝的袍子。喊叫是什么？\InnerQuote{你背叛了皇帝。}钉死基督的身体。喊叫是什么？\InnerQuote{基督为我们死了。}但我和你还需要什么？让我们到福音作者那里去。你们从主那里如何领受主是如何死的？他们读到\BibleQuote{父啊，我把我的灵魂交托在你手中。}灵魂在高处，身体在十字架上为我们。就身体归属于祂自己而言，祂献出了羊。}\par

同一位作者，同一著作：\par

\Quote{祂来拯救我们的本性，而不是毁灭祂自己的本性。如果我说骆驼会飞，你直接会认为奇怪，因为这与它的本性不符；你这样做很正确。如果我说人生活在海里，你也不会接受；你做得对。这违反本性。那么，如果我对这些本性说出奇怪的事，你便认为奇怪；如果我说那在万世之前、本性\OriginalTerm{无身体的}{incorporeal}、尊严上\OriginalTerm{不受苦的}{impassible}、与圣父同在且坐在圣父右边、在荣耀中的大能，如果我说这无身体的本性受苦，你岂能不掩耳？如果你听到这些而不掩耳，我将掩心。我们对天使能做什么？用剑击打他？或把他切成碎片？为什么我说天使？我们对灵魂能吗？灵魂会接受钉子吗？灵魂既不割也不烧。你问为什么？因为它被如此创造。祂的作为是不受苦的，而祂自己却是\OriginalTerm{能受苦的}{passible}？我不拒绝那\OriginalTerm{经纶}{œconomy}；相反，我欢迎那虐待。基督为我们而死并被钉十字架。经上如此记载；本性也承认。我不涂去这些话，也不亵渎那本性。但这不是真的。很好，那么就让更真实的话被说出。老师是恩人，从不严厉，从不敌对，除非学生倔强。你有什么好话要说？我的耳朵欣然敞开。有人要争吵吗？让他从容争吵。犹太人能把天主之子钉死并使大能本身成为死尸吗？活着的能死吗？这大能的死亡就是它的失败。即使我们死，我们的身体还留下。但如果我们使那大能成为死尸，我们就把它归于无有。我怕你听不进去。如果身体死了，灵魂就与它分离而存留；但如果灵魂死了，既然它没有身体，它就完全停止存在。灵魂因死亡而完全停止存在。因为\OriginalTerm{不死的}{immortals}的死亡是对它们存在的否定。考虑另一种情况；因为我甚至不敢提它。我们按我们所理解的谈论这些，但如果有人好争论，我们不定什么法律。但我知道一件事：每个人都必须收获他意见的果实。每个人来到天主面前，将他对天主的所说所想带到祂面前。不要以为天主读书本，或为必须回忆你说过什么或谁听过而烦恼：一切都显明了。审判者坐在宝座上。\Person{保禄}被带到他面前。\InnerQuote{你说我是人；你不能与我同享生命。你不认识我；我不认识你。}另一个人上前。\InnerQuote{你说我是受造之物之一。你不认识我的尊严；我不认识你。}另一个人上前。\InnerQuote{你说我没有取了身体。你轻视了我的恩宠。你不得分享我的不死。}另一个人上前。\InnerQuote{你说我不是由童贞女所生以拯救童贞女的身体；你不得得救。}每人都收获他对信德意见的果实。}\par

你看，你所师从的另一派，你曾以为在其中学到了独生子的\OriginalTerm{神性}{Godhead}受苦，而这一派憎恶这亵渎，宣讲神性的\OriginalTerm{不受苦性}{impassibility}，并离开那敢于把\OriginalTerm{苦难}{passion}归于神性之人的行列。\par

\Emph{厄兰}：是的；我对这冲突感到惊奇，并且我钦佩这人的见识和意见。\par

\Emph{正统}：那么，我亲爱的先生，请仿效蜜蜂。当你在心智的飞翔中飞过圣经的草地，在这些卓越教父的美丽花朵中，请在你心中建造信德的蜂巢。如果你偶然发现某处草苦涩不可食，如同这些\Person{阿波林}和\Person{欧瑟比}之流，但仍然并非完全没有可用于酿蜜的东西，你应当啜取甘甜而留下有毒之物，正如蜜蜂常落在有害的灌木上，却留下所有致命的毒物而收集一切好的。我们以兄弟之爱给你这忠告，亲爱的朋友。接受它，你便会做得好。如果你不听，我们将用宗徒的话对你说：\Quote{我们是洁净的。}正如先知所说，我们说了所吩咐的。\par

\SourceRecord{Translated by Blomfield Jackson.；Nicene and Post-Nicene Fathers, Second Series；Vol. 3.；Edited by Philip Schaff and Henry Wace.；Buffalo, NY: Christian Literature Publishing Co.,；1892.；<http://www.newadvent.org/fathers/27033.htm>.}
